\chapter{Classical and Electromagnetic Readings of X}
\label{ch:ch03_readings}

The question of this chapter is what the two readings of one vector $X\in\R^{3,3}$---a displacement and a field strength---share,
and what keeps them apart. Both carry the metric $\eta$ and the Clifford action; they differ in units, in their natural Lorentz
algebras, in their canonical conjugates and in their observables, and no equivariant linear map joins them (\cref{r1:thm:type}). The
chapter derives the Lagrangians of both readings, the four-momentum read by a clock (under \cref{assum:clockred}), the field-octant map
and Maxwell's equations as one first-order equation on the octants, the bridge $T=\hat nJ(\hat n)$ between the Dirac structure and the
octants, spin one in both readings, the orthogonality of $\bE$ and $\bB$, and the undivided energy $E_6$ with its rest part $E_0$. The
rest energy $E_0$ is a pointwise invariant, not a generated mass; the even-Clifford obstruction is specific to the octant representation;
the coupling classification holds in its stated class.


% --- source: two_readings/r0_intro#0 ((head))
% Introduction of the two-readings note.

% --- source: two_readings/r0_intro#1 (Introduction)
\section{Introduction}
\label{r0:sec:main}

The letter $X$ is used in this arena for two different things. In the momentum--force construction it is a displacement,
a vector of $V=\R^{3,3}$ with the dimension of a length; in the field-strength reading of \cref{sec:emfield} the same vector space
carries the electric and magnetic fields, $X=\sum_i(B_ie_i+E_ie_{i+3})$. One may keep the two
readings strictly apart, $X_{\rm cl}$ for the displacement and $X_{\rm EM}=(\bB,\bE)$ for the six components of the field two-form, and to
build the four-momentum of each separately. This chapter answers four questions put on that basis:
\begin{enumerate}[label=(\arabic*),leftmargin=2.2em]
\item construct the Lagrangians for $X$, for the Dirac field and for the energies;
\item differentiate between the two $X$'s, the first a length and the second the pair $(\bB,\bE)$;
\item construct the four-momentum, the length-valued $X$ being taken along the three directions of the spatial triple $V^-$, of signature
$(-,-,-)$;
\item place $X_{\rm EM}=(\bB,\bE)$ on the octants, each of the six components $E_x,E_y,E_z$, $B_x,B_y,B_z$ on an octant.
\end{enumerate}
Question (4) is read as the pair: $X_{\rm EM}$ is not a sum of two vectors but the six independent components of one field strength.
Questions (1) and (2) are answered in \cref{r1:sec:main}, question (3) in \cref{r2:sec:main} and question (4) in \cref{r3:sec:main}.
\Cref{r4:sec:main} takes up the element that the answers share: the temporal volume $T=e_1e_2e_3$, which is the Dirac mass form of the
classical reading for every clock $\hat n$, through $T=\hat nJ(\hat n)$ with $J(\hat n)$ the Dirac complex structure of that clock, and the duality operator
of the field reading on the octants. \Cref{r5:sec:main} collects the answers.

\paragraph{What is established, in brief.} The two readings share a vector space and a Clifford action, but their units, their natural
Lorentz algebras, their canonical conjugates and their observables differ, and no equivariant linear map relates them
(\cref{r1:thm:type}). The Lagrangians are the six-dimensional worldline and the interior shifted Lagrangian for $X_{\rm cl}$, the real
Grassmann Lagrangian for the Dirac field, and the Maxwell Lagrangian $-\tfrac12\Re q$ for the field, whose Hamiltonian is the undivided
energy $E_6$; the rest energy $E_0=\tfrac12|q|$ is a Lorentz-invariant and conformal Lagrangian, but not a duality-invariant one and not an
admissible Hamiltonian (\cref{r1:thm:L0,r1:prop:ham}). With a clock $\hat n$ on the temporal triple, the four-momentum $(E/c,\bp)$ has $\bp$
conjugate to the three lengths of $V^-$ and $E/c$ conjugate to the time read by $\hat n$, and the four-dimensional mass read by $\hat n$ never exceeds
the six-dimensional one (\cref{r2:prop:cons,r2:prop:rotate}); the field's four-momentum comes from its stress tensor and is a four-vector
only for an isolated whole (\cref{r2:sec:field}). On the octants, the joint eigenlines of the pole operators, $\bE$ sits on the three
octants of Fock weight one and $\bB$ on the three of weight two; Maxwell's four equations become one first-order equation with the two
Gauss laws on the two poles; and the boosts and the duality rotation of the field act on the octants only through odd operators, so that
the octant representation of the field is not the spin representation in which a Dirac field transforms (\cref{r3:thm:noeven,r3:prop:odd}).

\paragraph{Conventions and classes.} The conventions are those of \cref{sec:conventions}: $\eta=\operatorname{diag}(1,1,1,\allowbreak-1,-1,-1)$, the temporal
triple $V^+=\langle e_1,e_2,e_3\rangle$ carries the converted times of the classical reading and $\bB$ of the field reading, the spatial
triple $V^-=\langle e_4,e_5,e_6\rangle$ the lengths and $\bE$; the paper's $\bB$ is minus the textbook magnetic field. Each section states
its units. Every claim carries a class stated in words: \emph{fixed} (an identity or theorem that could have failed), \emph{constructed}
(admitted and exhibited, not selected), \emph{obstructed} (impossible under stated hypotheses), \emph{open}, \emph{inherited} (true of
ordinary physics in any notation, cited), \emph{interpretive} (a reading), or \emph{restated} (established in another chapter and re-verified here).
Every finite-dimensional and symbolic claim is checked by a script; the checks are described in \cref{r9:sec:main}.

% --- source: tex/s12_emfield#0 ((head))


% --- source: tex/s12_emfield#1 (The split space as a space of field strengths)
\section{The split space as a space of field strengths}
\label{sec:emfield}

The arena was introduced with an electromagnetic picture in mind: six directions, three of one sign and three of the other, paired
one to one like the components of a magnetic and an electric field. This section turns the picture into a theorem and measures its
price. The binding root $G$ of \cref{prop:twoK} is the key, and its role in the interior dynamics, where it is the operator that
makes bound motion possible, is established in \cref{sec:neutral}. Here it is used only as a linear map with $G^2=-\id$ and
$G=G^\dagger$.

% --- source: tex/s12_emfield#2 (Arithmetic first)
\subsection{Arithmetic first}

The space $\Lambda^2\R^{1,3}$ with the induced metric $\langle\theta^a\wedge\theta^b,\theta^a\wedge\theta^b\rangle=\eta^{aa}\eta^{bb}$ is
negative on the three planes $\theta^0\wedge\theta^i$ and positive on the three planes $\theta^j\wedge\theta^k$. It therefore has
signature $(3,3)$, the Lorentz group acts on it by the adjoint representation, and the Hodge star satisfies $\star^2=-\id$. This is
counting, and it must be said before anything is built on it: every six-dimensional space of split signature is isometric to a
Lorentzian field-strength space, so exhibiting one isometry predicts nothing. What is not arithmetic is which subalgebra
$\so{1,3}\subset\so{3,3}$ realises the action, and whether the operator playing the Hodge star is one the arena already had.

% --- source: tex/s12_emfield#3 (The Lorentz algebra of the binding root)
\subsection{The Lorentz algebra of the binding root}

\begin{theorem}[The centraliser of $G$]
\label{thm:GLorentz}
The centraliser $\lG$ of $G$ in $\so{3,3}$ is isomorphic to $\so{3,\C}\cong\slr{2,\C}\cong\so{1,3}$. Its maximal compact
subalgebra is the $\so3$ of \emph{diagonal} rotations, each acting identically on a temporal plane $(e_i,e_j)$ and on its spatial
partner $(e_{i+3},e_{j+3})$.
\end{theorem}

\begin{proof}
Since $G^2=-\id$, the identification $z_j=x_j+ix_{j+3}$ makes $\R^6$ a complex space $\C^3$ on which $G$ acts as $i$, and
commuting with $G$ is $\C$-linearity. Under this identification
\[
\eta(X,X)=\sum_jx_j^2-\sum_jx_{j+3}^2=\Re q(z),\qquad q(z):=\sum_jz_j^2,
\]
with $q$ holomorphic. A $\C$-linear map $A$ preserving $\Re q$ preserves it on every phase rotation $e^{i\vartheta}z$, and
$q(e^{i\vartheta}z)=e^{2i\vartheta}q(z)$, so $\Re\bigl(e^{2i\vartheta}[q(Az)-q(z)]\bigr)=0$ for all $\vartheta$ forces $q(Az)=q(z)$. The
group is $O(3,\C)$ and its algebra $\so{3,\C}$, whose underlying real Lie algebra is $\slr{2,\C}\cong\so{1,3}$ with maximal compact
$\so3$. A rotation of $(e_i,e_j)$ alone does not commute with $G$, because $G$ carries $e_i$ to $e_{i+3}$; only the rotation
performed identically on the partner plane does.
\end{proof}

\begin{proposition}[Two Lorentz algebras, not conjugate]
\label{prop:notconj}
The interior algebra $\lint$ and $\lG$ are isomorphic as Lie algebras but not as subalgebras of $\so{3,3}$. The interior algebra
fixes the two-dimensional surplus plane pointwise, while $\lG$ contains simultaneous rotations of both triples and has no nonzero
common fixed vector. Conjugation preserves the dimension of the common fixed subspace, so the two are not conjugate in $GL(6,\R)$,
let alone in $O(3,3)$.
\end{proposition}

Write $L_i$ for the diagonal rotations and put, for cyclic $(i,j,k)$, $K_i=\tfrac12\bigl(e_ke_{j+3}-e_je_{k+3}\bigr)$. The six
operators $L_i,K_i$ span $\lG$.

\begin{theorem}[The field-strength reading]
\label{thm:emfield}
Write a vector of $\R^{3,3}$ as $X=\sum_i\bigl(B_ie_i+E_ie_{i+3}\bigr)$. Then
\[
[L_i,L_j]=\epsilon_{ijk}L_k,\qquad [L_i,K_j]=\epsilon_{ijk}K_k,\qquad [K_i,K_j]=-\epsilon_{ijk}L_k,
\]
both triples $\bB$ and $\bE$ transform as vectors under the $L_i$, and the $K_i$ act by
\begin{equation}
\label{eq:emboost}
\delta\bE=\bv\times\bB,\qquad \delta\bB=-\bv\times\bE .
\end{equation}
Identify $X$ with the two-form $F_X$ of components $(F_X)_{0i}=E_i$ and $(F_X)_{jk}=\epsilon_{ijk}B_i$. Then $G$ is the Hodge star,
$F_{GX}=\star F_X$ with $(\star F)_{ab}=\tfrac12\varepsilon_{abcd}F^{cd}$ and $\varepsilon_{0123}=+1$, and
\[
\eta(X,X)=|\bB|^2-|\bE|^2,\qquad \eta(X,GX)=-2\,\bE\cdot\bB
\]
are the only invariant quadratic forms. With $\lG$ as its Lorentz algebra, $\R^{3,3}$ \emph{is} $\Lambda^2\R^{1,3}$.
\end{theorem}

The orientation matters once, and we fix it here. With $F_{jk}=\epsilon_{ijk}B_i$ the vector $\bB$ is minus the textbook magnetic
field of the convention $F_{jk}=-\epsilon_{ijk}B_i$ \cite{Jackson}, so we write $\Btext=-\bB$, as announced in \cref{sec:conventions};
the electric field is the textbook one. Passing to a frame moving with velocity $\mathbf u$ changes the frame variables by
$\delta\bB=\mathbf u\times\bE$ and $\delta\bE=-\mathbf u\times\bB$, which is \eqref{eq:emboost} with $\bv=-\mathbf u$; in textbook
variables this is $\bE'=\bE+\mathbf u\times\Btext$ and $\Btext'=\Btext-\mathbf u\times\bE$, the familiar first-order transformation.

\begin{example}[A boost along the third axis]
\label{ex:emboost}
Take the pure magnetic field $X=e_1$, that is $\bB=(1,0,0)$ and $\bE=0$, and the boost generator $K_3=\tfrac12(e_2e_4-e_1e_5)$,
acting on vectors by the commutator. Since $e_1$ anticommutes with $e_2$ and with $e_4$ it commutes with $e_2e_4$, while
$[e_1e_5,e_1]=-2e_1^2e_5=-2e_5$; hence $[K_3,e_1]=e_5$. Under $\delta X=\delta\,[K_3,X]$ the field acquires $\bE=(0,\delta,0)$, in accordance
with $\delta\bE=\bv\times\bB$ for $\bv=(0,0,\delta)$. In the same way $[K_3,e_5]=e_1$, so a pure electric field $\bE=(0,1,0)$ acquires
$\bB=(\delta,0,0)=-\bv\times\bE$. The invariant $\eta(X,X)=|\bB|^2-|\bE|^2$ changes only at
second order, as it must.
\end{example}

\begin{corollary}[The invariant complex structure]
\label{cor:Gselected}
$\R^6$ is irreducible under $\lG$, and the commutant of $\lG$ in $\End(\R^6)$ is $\operatorname{span}\{\id,G\}\cong\C$. Hence, once
$\lG$ is taken as the Lorentz algebra, its invariant complex structure is $\pm G$ and nothing else. The interior algebra, by contrast,
has a commutant of dimension five, $\R$ on the four-plane and all of $M_2(\R)$ on the surplus plane.
\end{corollary}

The logic of this corollary runs from $G$ to $\lG$ and back to $\pm G$, so the choice of $G$ is recovered, not removed. It is a
consistency statement, a uniqueness conditional on taking $\lG$ as the Lorentz algebra. Removing the choice would require a
characterisation of $\lG$ that does not mention $G$, and we know none.

% --- source: tex/s12_emfield#4 (The two readings exclude one another)
\subsection{The two readings exclude one another}

\label{sec:readingsexclude}
The price of the field-strength reading is not small.

\begin{proposition}[Exclusion]
\label{prop:emexclusive}
Under $\lint$ the six directions decompose as $V\oplus2\cdot\mathbf1$, a four-vector and two scalars; under $\lG$ they are
irreducible, with no fixed vector. The two representations are not similar, so a single theory cannot use both algebras as the
Lorentz action on the same six directions. If the vector space $\R^{3,3}$ is read as the electromagnetic field strength, those six
directions are not a four-plane plus a surplus pair, and the Dirac structure, the projector census, $P_*$ and the diagonal sector
all lose the arena on which they are built.
\end{proposition}

What is excluded is reading one six-dimensional space in two ways, not a theory that contains both kinds of module. A theory may
contain spacetime vectors and field strengths as different representations of one Lorentz algebra, which is the ordinary
arrangement, and $\Cl(3,3)$ realises it with the vectors at grade one and the field strength at grade two.

\begin{proposition}[The field strength at grade two]
\label{prop:bivectorF}
Let $\Lambda^2E$ be the span of the six bivectors of the interior four-plane. It is closed under the adjoint action of $\lint$, on
which it is the adjoint representation; the trace form $\tfrac18\Tr(XY)$ on it has signature $(3,3)$; and the interior volume
$\Omega=\gamma^0\gamma^1\gamma^2\gamma^3$, with $\Omega^2=-\id$, commutes with $\lint$, preserves $\Lambda^2E$ by left multiplication,
exchanges $\gamma^0\gamma^i$ with the spatial bivectors and acts on $\Lambda^2E$ as \emph{minus} the Hodge star of \cref{thm:emfield}
($\varepsilon_{0123}=+1$, the bivector $\tfrac12F_{ab}\gamma^a\gamma^b$ standing for the two-form $F$). Equivalently, the space-time-algebra
pseudoscalar $I_4:=\gamma_0\gamma_1\gamma_2\gamma_3=-\Omega$, written with the lower-index generators $\gamma_\mu=\eta_{\mu\nu}\gamma^\nu$, acts as
the Hodge star.
\end{proposition}

Grade three behaves differently under the two algebras, and the difference sharpens the exclusion. Under $\lG$ one has
$(\Lambda^3\R^6)_\C=2\,(0,0)\oplus2\,(1,1)$, of dimension $2+18$, so no field-strength module $(1,0)\oplus(0,1)$ occurs at grade three
at all: trivectors cannot carry the field for this action. Under $\lint$, $\R^6=E\oplus\R^2$ gives $\Lambda^3\R^6\cong2E\oplus2\Lambda^2E$,
which contains two field-strength modules. The quadratic Casimir distinguishes the cases cleanly, with spectrum $\{0^{(2)},8^{(18)}\}$
on $\Lambda^3\R^6$ under $\lG$ and $\{3^{(8)},4^{(12)}\}$ under $\lint$, the value $4$ being that of $(1,0)\oplus(0,1)$ and $3$ that of a
four-vector. An explicit attempt confirms it: the six trivectors $C_i=\sum_{j,k}\epsilon_{ijk}e_jN_k$, $N_j=e_{k+3}e_{l+3}$ cyclically,
together with $D_i=\mathcal VC_i$, transform as two vectors under the $L_i$, but all eighteen brackets $[K_i,C_j]$, $[K_i,D_j]$ leave
their span, for instance $\tfrac12[K_1,C_1]=e_2e_3e_4+e_4e_5e_6$, and the Lorentz closure of the span has dimension $18$.

% --- source: tex/s12_emfield#5 (How much of Maxwell's theory is forced)
\subsection{How much of Maxwell's theory is forced}
\label{sec:maxwell}

\Cref{thm:emfield} supplies a field with the right transformation law and nothing that the field satisfies. How much of the
dynamics does the transformation law fix by itself? The answer is: the form of the equations completely, and nothing else.

\begin{theorem}[Exactly two first-order covariant operators]
\label{thm:twoops}
Consider linear first-order differential operators with constant coefficients on Minkowski space, acting on a field $F$ valued in
$\Lambda^2\R^{1,3}$ and taking values in the vector representation. The derivative $\partial F$ lies in $\Lambda^2\otimes V$, of
dimension $24$, and under $\so{1,3}$
\[
\Lambda^2\otimes V=\underbrace{(\tfrac32,\tfrac12)\oplus(\tfrac12,\tfrac32)}_{16}\ \oplus\ \underbrace{2\,(\tfrac12,\tfrac12)}_{8},
\]
so the vector module occurs with multiplicity exactly two, and no other module of dimension four occurs. The two copies are realised
by
\[
(dF)_{\mu\nu\rho}=3\,\partial_{[\mu}F_{\nu\rho]}\quad(\text{through }\Lambda^3\cong V),\qquad (\delta F)_\nu=\partial^\mu F_{\mu\nu},
\]
and every covariant operator of the class is $\mathcal DF=a\,\delta F+b\,\delta(\star F)$, since $\delta\star=\pm\star d$.
\end{theorem}

The Casimir operator confirms the count: on $\Lambda^2\otimes V$ its spectrum is $\{3^{(8)},9^{(16)}\}$, and the value $3$ is the
vector one. Nothing in the theorem constrains operators of higher order, operators with variable coefficients, or operators valued
in the remaining $(\tfrac32,\tfrac12)\oplus(\tfrac12,\tfrac32)$.

\begin{corollary}[The form of Maxwell's equations]
\label{cor:maxwellform}
Within this class the two available equations are $dF=\mathcal J_m$ and $d\star F=\mathcal J_e$. Covariance does not rule out a
magnetic current, which is a covariant field exactly as an electric current is; it identifies the two operators and leaves open
which equations are imposed and which sources are admitted. The source-free system $dF=\delta F=0$ is the only choice that needs no
external data, and within this linear class it is the only one invariant under duality, because the Hodge star $G$ exchanges the two
operators. Finally, on a contractible region $dF=0$ gives $F=dA$ by the Poincar\'e lemma, so the gauge potential is a consequence of
the first equation rather than an assumption.
\end{corollary}

In the variables of the frame the system takes a compact form, which we use repeatedly below.

\begin{proposition}[Maxwell's equations on the frame]
\label{prop:RS}
Read $\R^{3,3}$ as $\C^3$ through $G$, $z_j=B_j+iE_j$. Then $dF=0$ and $d\star F=0$ for the two-form $F_X$ of
\cref{thm:emfield} are
\begin{equation}
\label{eq:RS}
\partial_tX=-G\,(\curl X),\qquad \nabla\cdot X=0,
\end{equation}
that is $\partial_tz=-i\,\curl z$ and $\nabla\cdot z=0$, with the curl built from the rotation generators $L_i$. Splitting into the
two triples gives $\partial_t\bB=\curl\bE$, $\partial_t\bE=-\curl\bB$ and $\nabla\cdot\bE=\nabla\cdot\bB=0$; in textbook variables these
are Faraday's law $\partial_t\Btext=-\curl\bE$ and the Amp\`ere--Maxwell law $\partial_t\bE=\curl\Btext$ in vacuum. With a source,
$\partial_\mu F^{\mu\nu}=(\rho,\bj)$ reads
\begin{equation}
\label{eq:RSsource}
\partial_tz=-i\,\curl z-i\,\bj,\qquad \nabla\cdot z=i\rho ,
\end{equation}
and taking the divergence of the first equation gives $\partial_t\rho+\nabla\cdot\bj=0$: continuity is a consequence of the system and
the compatibility condition that any proposed current must satisfy.
\end{proposition}

Up to a factor $i$, $z$ is the Riemann--Silberstein vector $\bE+i\Btext$, which obeys
$i\partial_t(\bE+i\Btext)=\curl(\bE+i\Btext)$ \cite{Silberstein1907,BialynickiBirula13}, and the holomorphic form of the proof of
\cref{thm:GLorentz} has
\[
\Re q(z)=|\bB|^2-|\bE|^2,\qquad \Im q(z)=2\,\bE\cdot\bB ,
\]
so the two Maxwell invariants are the real and imaginary parts of one holomorphic quadratic form. Three classical facts fall into
place. The duality rotation is the phase rotation $z\mapsto e^{i\vartheta}z$ that the proof of \cref{thm:GLorentz} already used, and
it turns the pair of invariants by $2\vartheta$. The two quadratic forms that will compete in the interior dynamics of
\cref{sec:shifted,sec:neutral} reappear with physical names: the $\eta$-form $\Re q$ is the Lorentz invariant, while the Hermitian
form $|z|^2=|\bB|^2+|\bE|^2$ is twice the energy density and depends on the frame. And the remaining sesquilinear covariant is the
Poynting vector, $\tfrac i2\,\bar z\times z=\bE\times\bB$, so that $\bE\times\Btext=-\tfrac i2\,\bar z\times z$.

\begin{example}[A circularly polarised wave on the frame]
\label{ex:circular}
The field $z=(1,i,0)\,e^{ik(x_3-t)}$ solves \eqref{eq:RS}: $\nabla\cdot z=0$ because $z$ has no third component, and
$\curl z=ik\,(-i,1,0)e^{ik(x_3-t)}=k\,(1,i,0)e^{ik(x_3-t)}$, so $-i\curl z=-ik\,z=\partial_tz$. Its invariant is
$q(z)=(1+i^2)e^{2ik(x_3-t)}=0$, so $|\bB|=|\bE|$ and $\bE\cdot\bB=0$ at every point: the wave is null. Its Poynting vector
$\tfrac i2\bar z\times z=\tfrac i2(1,-i,0)\times(1,i,0)=\tfrac i2(0,0,2i)=(0,0,-1)$ is $\bE\times\bB$, so $\bE\times\Btext=(0,0,1)$
points along the direction of propagation, as it should.
\end{example}

% --- source: tex/s12_emfield#6 (Frame vectors as spacetime planes)
\subsection{Frame vectors as spacetime planes}

Each pairing plane $(e_i,e_{i+3})$ carries one complex component $z_i=B_i+iE_i$, and a frame vector is at the same time a bivector of
spacetime. The two descriptions meet in the following way.

\begin{proposition}[Frame vectors as planes]
\label{prop:planes}
Let $X=\sum_i(B_ie_i+E_ie_{i+3})$ and $F_X$ as above. Then $\Pf F_X=\bE\cdot\bB$ and $\det F_X=(\bE\cdot\bB)^2$. A nonzero $F_X$ is a
single plane $u\wedge v$ exactly when $\bE\cdot\bB=0$, that is when $\Im q=0$. Otherwise there is a Lorentz frame in which
$F_X=a\,\theta^0\wedge\theta^1+b\,\theta^2\wedge\theta^3$ is a pair of orthogonal planes, and the duality rotation $e^{-\vartheta G}$ with
$\tan\vartheta=a/b$ carries it to the single plane $\sqrt{a^2+b^2}\,\theta^2\wedge\theta^3$, the complexion of Rainich and of Misner
and Wheeler \cite{Rainich1925,MisnerWheeler57}.
\end{proposition}

\Cref{fig:qplane} shows the content of this proposition in the plane of the complex invariant. The real axis is the locus of single
planes, the positive half of magnetic type and the negative half of electric type, the origin is the null cone, and the duality
rotation moves a field around a circle at twice its own angle. In \cref{sec:neutral} the bound orbits of the interior dynamics will
turn out to be exactly such rotations.

\begin{figure}[t]
\centering
\begin{tikzpicture}[scale=1.15,>=Stealth]
\draw[->,gray] (-3.0,0) -- (3.2,0) node[right] {\small $\Re q=|\bB|^2-|\bE|^2$};
\draw[->,gray] (0,-2.6) -- (0,2.8) node[above] {\small $\Im q=2\,\bE\cdot\bB$};
\draw[scol,line width=2.2pt] (-2.9,0) -- (-0.08,0);
\draw[tcol,line width=2.2pt] (0.08,0) -- (2.9,0);
\node[tcol,anchor=north west] at (0.15,-0.08) {\small single plane, magnetic};
\node[scol,anchor=north east] at (-0.15,-0.08) {\small single plane, electric};
\fill[black] (0,0) circle (2pt);
\node[anchor=north west] at (0.05,0.45) {\small null};
\draw[gcol,thick,dashed] (0,0) circle (1.8);
\coordinate (q0) at ({1.8*cos(35)},{1.8*sin(35)});
\fill[gcol] (q0) circle (2.3pt);
\node[gcol,right=3pt] at (q0) {\small $q_0$};
\draw[->,gcol,thick] ({2.2*cos(35)},{2.2*sin(35)}) arc (35:95:2.2);
\node[gcol,align=left,anchor=west] at (1.75,2.05) {\small duality $e^{\vartheta G}$:\\ \small $q\mapsto e^{2i\vartheta}q$};
\foreach \a in {0,180}{\fill[ocol] ({1.8*cos(\a)},{1.8*sin(\a)}) circle (2.0pt);}
\node[ocol,anchor=north east,align=right] at (-1.0,-1.55) {\small each crossing of the real axis\\ \small is met twice per turn of $\vartheta$};
\end{tikzpicture}
\caption{The plane of the complex invariant $q=z\cdot z$ of a frame vector (\cref{thm:emfield,prop:planes}). The real axis is the
locus $\bE\cdot\bB=0$ on which the spacetime two-form is a single plane, magnetic on the right and electric on the left; the origin is
the null cone. A duality rotation $e^{\vartheta G}$ preserves $|q|$ and turns $q$ at twice the rate, so a non-null field becomes a single
plane four times as $\vartheta$ runs once around.}
\label{fig:qplane}
\end{figure}

Everything in this section is a statement about a representation and its invariants. \Cref{thm:twoops} is a multiplicity count that
fixes the form of the equations and nothing about their content: it does not produce a source, does not say that a source is the
current of anything, supplies no coupling constant and no Lagrangian. Nor is any of it special to this arena, since the count uses only
$\so{1,3}$ and $\Lambda^2$. What the arena adds is that the objects the count needs, the module, the duality and the invariants, were
already present, so that the system can be written in the arena's own variables \eqref{eq:RS}. The deflationary test of
\cref{sec:deflation} applies to it in full.

% --- source: two_readings/r1_lagrangians#0 ((head))
% Section r1: the two readings of X and their Lagrangians.  Verification: r1_lagrangians.py (ids r1-n).

% --- source: two_readings/r1_lagrangians#1 (Two readings of XX and their Lagrangians)
\section{Two readings of \texorpdfstring{$X$}{X} and their Lagrangians}
\label{r1:sec:main}

The question asked in the introduction contains two requests that this section answers together: that the two uses of the letter $X$
be told apart, the first being a length and the second the pair $(\bB,\bE)$ of three-vectors, and that Lagrangians be written for $X$, for the Dirac field
and for the energies. The section types the two readings (\cref{r1:sec:types}), writes the Lagrangians of the classical reading
(\cref{r1:sec:cl}), of the Dirac field (\cref{r1:sec:dirac}) and of the energies (\cref{r1:sec:energies}), assembles them
(\cref{r1:sec:total}) and states what is fixed and what is not (\cref{r1:sec:answer}). Identifiers r1-$\hat n$ refer to the lines
\texttt{PASS r1-$\hat n$} printed by \texttt{r1\_lagrangians.py}, which checks every finite-dimensional and symbolic claim in exact
arithmetic; Euler--Lagrange equations, Noether currents and stress tensors are derived there on general symbolic fields, and the
Grassmann-valued Dirac field is handled by an exact Grassmann jet algebra.

Conventions are those of \cref{sec:conventions}: $V=\R^{3,3}$ with $\eta=\operatorname{diag}(1,1,1,\allowbreak-1,-1,-1)$, the temporal triple
$V^+=\langle e_1,e_2,e_3\rangle$ and the spatial triple $V^-=\langle e_4,e_5,e_6\rangle$, the realisation of \cref{app:realisation},
the binding root $Ge_i=e_{i+3}$, $Ge_{i+3}=-e_i$, which satisfies $G^2=-\id$, $\eta G=(\eta G)^{\!\top}$ and $G^{\!\top}\eta G=-\eta$
(r1-1), the interior Lorentz algebra $\lint$ on $E=\langle e_1,e_4,e_5,e_6\rangle$ with $\gamma^\mu=(e_1,e_4,e_5,e_6)$, $J=e_2e_3$ and
$T=e_1e_2e_3=\gamma^0J$. Spacetime signature is $(+,-,-,-)$. The paper's $\bB$ is minus the textbook magnetic field,
$\Btext=-\bB$ (\cref{sec:conventions}); every Maxwell formula below names which of the two it uses. Each claim carries a class in
words: \emph{fixed}, \emph{constructed}, \emph{obstructed}, \emph{open}, \emph{inherited}, \emph{interpretive} or \emph{restated}.

% --- source: two_readings/r1_lagrangians#2 (The two readings, typed)
\subsection{The two readings, typed}
\label{r1:sec:types}

\begin{definition}[The classical reading]
\label{r1:def:cl}
A classical displacement in the arena is
\begin{equation}
\label{r1:eq:Y}
X_{\rm cl}=Y=\sum_{i=1}^3c\,t^i\,e_i+\sum_{a=1}^3x^a\,e_{a+3},\qquad \eta(Y,Y)=\sum_i(ct^i)^2-\sum_a(x^a)^2,
\end{equation}
with $[t^i]=\mathrm s$, $[x^a]=\mathrm m$, $[c\,t^i]=\mathrm m$ and $[\eta(Y,Y)]=\mathrm m^2$. Its length-valued part $\sum_ax^ae_{a+3}$ lies on the spatial triple: this is the precise content of the phrase that $X$ is taken in the
three directions of sign $(-,-,-)$. The temporal components are times, converted to lengths by the constant $c$.
\end{definition}

The conversion is not a choice. A map $C:V^+\to V^-$ that commutes with the common rotations of the two triples is $C=c\,\id$ with
$[c]=\mathrm{m\,s^{-1}}$ (\cref{thm:onec}; \cref{en:onec}), and the binding root acts on the unconverted
coordinates $(t;x)$ as
\begin{equation}
\label{r1:eq:GC}
G_C=\begin{pmatrix}0&-C^{-1}\\ C&0\end{pmatrix}=\Phi^{-1}G\,\Phi,\qquad \Phi=\operatorname{diag}(c,c,c,1,1,1),\qquad \Phi(t;x)=Y .
\end{equation}
Hence $G$ is dimensionless on $Y$, while $G_C$ carries entries of dimension $\mathrm{m\,s^{-1}}$ and $\mathrm{s\,m^{-1}}$: $G$ maps
lengths to lengths only after the conversion, and on $Y$ it sends the length $x^a$ to the converted time slot $e_a$ (r1-2; restated from \cref{sec:clock}, the commutation criterion being fixed there).

\begin{definition}[The field reading]
\label{r1:def:em}
The field-strength reading of \cref{thm:emfield} is, with SI units and $c$ explicit,
\begin{equation}
\label{r1:eq:XEM}
X_{\rm EM}=\sum_{i=1}^3\bigl(cB_i\,e_i+E_i\,e_{i+3}\bigr),\qquad z_c=c\bB+i\bE,\qquad
\eta(X_{\rm EM},X_{\rm EM})=c^2|\bB|^2-|\bE|^2=\Re\,(z_c\cdot z_c),
\end{equation}
with $[E_i]=[cB_i]=\mathrm{V\,m^{-1}}$, so $[\eta(X_{\rm EM},X_{\rm EM})]=\mathrm{V^2\,m^{-2}}$. $X_{\rm EM}$ is the pair $(c\bB;\bE)$ of
three-vectors, the coordinates of one vector of $\R^{3,3}$ in the basis $e_1,\dots,e_6$, and not a sum of $\bE$ and $\bB$: its six
entries label the six components of the field two-form, $F_{0i}=E_i$ and $F_{jk}=\epsilon_{ijk}B_i$ in Heaviside--Lorentz units with
$c=1$. In these units $X=\sum_i(B_ie_i+E_ie_{i+3})$,
$z=\bB+i\bE$, $q=z\cdot z$, $E_6=\tfrac12|z|^2$, $E_0=\tfrac12|q|$, and $E_6^2=E_0^2+|\mathbf S|^2$ with $\mathbf S=\bE\times\bB$
(\cref{en:split}). The same matrix $\Phi$ of \eqref{r1:eq:GC} converts $(\bB;\bE)$ to $(c\bB;\bE)$.
\end{definition}

\begin{theorem}[The type rule]
\label{r1:thm:type}
The two readings share the vector space $\R^{3,3}$, its form $\eta$ and the Clifford action on $S=\R^8$, and differ in the following
four respects.
\begin{enumerate}[label=(\roman*),leftmargin=2.2em]
\item \emph{Units.} A linear map $X_{\rm EM}\mapsto X_{\rm cl}$ has entries of dimension $\mathrm{m^2\,V^{-1}}$. No monomial
$c^a\ell^b\hbar^d$, with real exponents, has this dimension, where $\ell$ is the length of the shift operator (equation~\cref{eq:kappa}); with the vacuum permittivity admitted, the only such monomial is $\ell^3(\varepsilon_0/\hbar c)^{1/2}$.
\item \emph{Lorentz action.} Once the clock $e_1$ is chosen, $\lint$ acts on $Y$ as $E\oplus2\cdot\mathbf1$: the four-vector
$(ct^1;x^1,x^2,x^3)$, which complexifies to $(\tfrac12,\tfrac12)$, and the two scalars $ct^2,ct^3$. The algebra that acts naturally on
$X_{\rm EM}$ is $\lG$, the centraliser of $G$ in $\so{3,3}$, under which $V$ is irreducible and complexifies to $(1,0)\oplus(0,1)$
(\cref{sec:emfield}). The two subalgebras of $\so{3,3}$ meet in $0$.
\item \emph{No equivariant map.} For every Lie algebra isomorphism $\phi:\lint\to\lG$ let $V_\phi$ be $V$ with $\lint$ acting through
$\phi$. Then $\Hom_{\lint}(E,V_\phi)=\Hom_{\lint}(V_\phi,E)=0$ and $\Hom_{\lint}(E\oplus2\cdot\mathbf1,V_\phi)=\Hom_{\lint}(V_\phi,E\oplus2\cdot\mathbf1)=0$.
By contrast $\Hom_{\lint}(\Lambda^2E,V_\phi)=\{(a+bG)M_0\}$ is two-dimensional: the field strength of the interior four-plane, at
grade two (\cref{prop:bivectorF}), is carried to $X_{\rm EM}$ uniquely up to a scale and a duality angle.
\item \emph{Conjugates and observables.} The six components of $Y$ are configuration coordinates that Poisson-commute, with conjugate
covector $K_A$ of dimension $\mathrm{kg\,m\,s^{-1}}$; its observables are a worldline, a four-momentum and a mass. In Maxwell's theory the
configuration variable is the potential $\mathbf a$ with $\bB=\curl\mathbf a$, and $\bE$ is (up to $\varepsilon_0$) the momentum conjugate
to $\mathbf a$, so $X_{\rm EM}$ is a point of phase space, not of configuration space: $\{\int\mathbf f\cdot\bE,\int\mathbf g\cdot\bB\}=-\int\mathbf
f\cdot\curl\mathbf g$ (Heaviside--Lorentz) for test fields $\mathbf f,\mathbf g$, while the $\bE$--$\bE$ and $\bB$--$\bB$ brackets vanish. Its
observables are $E_6$, $\mathbf S$, the invariants $q$ and the stress tensor.
\end{enumerate}
\end{theorem}

\begin{proof}
(i) Over the base $(\mathrm m,\mathrm{kg},\mathrm s,\mathrm A)$ the volt is $\mathrm{kg\,m^2\,s^{-3}\,A^{-1}}$; none of $c$, $\ell$,
$\hbar$ contains the ampere, so the ampere exponent of $\mathrm{m^2V^{-1}}$ cannot be matched, and with $\varepsilon_0$ the linear system
for the exponents has the unique solution stated (r1-3). (ii) $\lint$ is spanned by $R_i=\tfrac12e_{j+3}e_{k+3}$ and
$B_i=\tfrac12e_1e_{i+3}$ ($(i,j,k)$ cyclic), acting on vectors by the commutator; they annihilate $e_2,e_3$ and preserve $E$. The
centraliser of $G$ in $\so{3,3}$ has dimension $6$ and is spanned by $L_i=\tfrac12(e_{j+3}e_{k+3}-e_je_k)$ and
$K_i=\tfrac12(e_ke_{j+3}-e_je_{k+3})$; its commutant on $V$ is $\operatorname{span}\{\id,G\}\cong\C$, a division algebra, and $V$ is
irreducible (\cref{cor:Gselected}). The spans of $\{R_i,B_i\}$ and $\{L_i,K_i\}$ inside $\so{3,3}$ have trivial intersection (r1-4).
(iii) Both triples satisfy $[R_i,R_j]=\epsilon_{ijk}R_k$, $[R_i,B_j]=\epsilon_{ijk}B_k$, $[B_i,B_j]=-\epsilon_{ijk}R_k$ and the same with
$(L,K)$, so $R_i\mapsto L_i$, $B_i\mapsto K_i$ is an isomorphism (r1-4). For any isomorphism $\phi$, $V_\phi$ is irreducible of dimension
$6$ and $E$ irreducible of dimension $4$. The image of an equivariant $M:E\to V_\phi$ is a submodule of dimension at most $4$, hence $0$;
the kernel of $M:V_\phi\to E$ is a submodule of dimension at least $2$, hence all of $V_\phi$. On the trivial summand an equivariant map
lands in the fixed vectors of $V_\phi$, of which there are none, and a map out of $V_\phi$ into it is zero by irreducibility. An intertwiner
$\Lambda^2E\to V_\phi$ exists for every $\phi$: the six bivectors of $E$ are $2R_i$ and $2B_i$, so $\Lambda^2E$ is the adjoint module of $\lint$
(\cref{prop:bivectorF}), $\phi$ intertwines it with the adjoint module of $\lG$ pulled back by $\phi$, and that module is
isomorphic to $V$ (r1-5). Schur's lemma \cite{GoodmanWallach} on the irreducible $\Lambda^2E\cong V_\phi$ then gives an intertwiner space
equal to the commutant $\operatorname{span}\{\id,G\}$.
The exact solution of the intertwining equations confirms all these dimensions for $\phi$ and for its parity twist $B_i\mapsto-K_i$ (r1-5),
and the Casimirs separate the modules: $\sum R_i^2-\sum B_i^2=-3$, $\sum R_iB_i=0$ on $E$, against $-4$ and $-2G$ (eigenvalues $\pm2i$) on
$V$ (r1-6). (iv) The bracket follows from $\{a_i(\mathbf x),E_j(\mathbf y)\}=\delta_{ij}\delta^3(\mathbf x-\mathbf y)$, read off the
first-order action of \cref{r1:prop:E6}(ii), and from $\mathbf g\cdot\curl\mathbf a-\mathbf a\cdot\curl\mathbf g=\nabla\cdot(\mathbf
a\times\mathbf g)$ (r1-31).
\end{proof}

Parts (i) and (iii) are fixed: (iii) could have failed for a different pair of subalgebras, and (i) obstructs only conversions built
from $c$, $\ell$ and $\hbar$ in SI units. Since $(\varepsilon_0/\hbar c)^{1/2}\bE$ has dimension $\mathrm m^{-2}$, which is the field of
Heaviside--Lorentz units with $\hbar=c=1$, in those units $\ell^3$ alone supplies the dimension (r1-3); what obstructs a dictionary between
the two readings is (iii), not (i). Part (ii) is restated from \cref{prop:notconj} and~\cref{prop:emexclusive} except its last sentence, which is
fixed (r1-4), and (iv) is inherited electrodynamics \cite{HenneauxTeitelboim92}. The identity map of $\R^{3,3}$ is linear but not equivariant between the two
actions, so sharing the vector space identifies nothing.

\begin{remark}[Coordinates and modules of the two readings]
\label{r1:rem:fmb1}
The surplus coordinates $u^a$ are
the converted times $ct^2,ct^3$ of the temporal triple, length-valued only through $C=c\,\id$. The distinction between the cotangent momentum $K$ and the momentum $\pi$ of the interior Lagrangian is a distinction between two Lagrangians for the same coordinates, as
\cref{r1:sec:cl} shows: applied to $Y$, the interior Lagrangian has canonical momentum $p=\eta\pi$. As a statement about $\lint$-modules,  $V_{\rm EM}\cong\Lambda^2E_{\hat n}^*$ holds through an isomorphism $\phi$, unique up to the factor $a+bG$, but
$\Lambda^2E$ is not a subspace of $V$, and $\lG$, under which $V$ is irreducible (\cref{cor:Gselected}), preserves no proper
subspace of $V$, while $G$ itself preserves no Lorentzian four-plane (\cref{r2:prop:nogo}). Parts (i) and (iii) give the rule of interpretation: no linear equivariant map exists in either
direction, for any isomorphism of the two Lorentz algebras, and in SI units no monomial in $c$, $\ell$, $\hbar$ supplies its dimension; in particular there is no linear equivariant map $\Lambda^2E_{\hat n}^*\to E_{\hat n}$, nor one in the opposite direction.
\end{remark}

% --- source: two_readings/r1_lagrangians#3 (Lagrangians for the classical reading)
\subsection{Lagrangians for the classical reading}
\label{r1:sec:cl}

\begin{proposition}[The free worldline in the arena]
\label{r1:prop:free}
Let $Y(\lambda)$ be a curve with $\eta(Y',Y')>0$, $\lambda$ an arbitrary parameter, $[M]=\mathrm{kg}$, and
\begin{equation}
\label{r1:eq:Lfree}
L_{\rm free}=-Mc\sqrt{\eta(Y',Y')},\qquad S=\int L_{\rm free}\,d\lambda,\quad [S]=\mathrm{J\,s}.
\end{equation}
\begin{enumerate}[label=(\roman*),leftmargin=2.2em]
\item $K_A=\partial L/\partial Y'^A=-Mc\,\eta_{AB}Y'^B/\sqrt{\eta(Y',Y')}$, with $\eta^{AB}K_AK_B=M^2c^2$ and $K_AY'^A-L=0$: the
canonical Hamiltonian vanishes, $L$ is homogeneous of degree one in $Y'$, and the dynamics is the constraint
$H=\tfrac N2(\eta^{AB}K_AK_B-M^2c^2)$. The kinetic six-momentum is $P^A=-\eta^{AB}K_B=McY'^A/\sqrt{\eta(Y',Y')}$.
\item The Euler--Lagrange equations are $dK_A/d\lambda=0$; their solutions are the straight lines, in every parametrisation.
\item The timelike set $\{\eta>0\}$ and the shell $\{\eta(K,K)=M^2c^2\}$ are connected: the curve $(\cos s,\sin s,0;0,0,0)$
joins $e_1$ to $-e_1$ with $\eta=1$, so there is no past and no future. With the clock $e_1$ and $P=(\mathsf E/c,q_1,q_2;\bp)$, the shell reads
$\mathsf E^2=c^2|\bp|^2+M^2c^4-c^2|\mathbf q|^2$; for $M=c=1$ and $\mathbf q=(2,0)$ the four-dimensional mass squared is $-3$.
\end{enumerate}
\end{proposition}

\begin{proof}
(i)--(ii) by differentiation (r1-7). A straight line $Y_0+Vf(\lambda)$ with $f'>0$ solves the equations, while the timelike curve
$(2\cos\lambda,2\sin\lambda,0;\lambda,0,0)$, with $\eta(Y',Y')=3$ and $Y''$ not parallel to $Y'$, does not (r1-8). (iii) Write $Y=(u;v)$
with $u\in V^+$, $v\in V^-$. Along $(u;(1-s)v)$, $0\le s\le1$, $\eta$ stays positive, so $\{\eta>0\}$ retracts onto $\{(u;0):u\neq0\}$
and hence onto the unit sphere of $V^+$, which is connected; the shell is the image of $S^2\times\R^3$ under
$(\omega,v)\mapsto(\sqrt{M^2c^2+|v|^2}\,\omega;v)$, with inverse $(u;v)\mapsto(u/|u|,v)$. Finally $\eta(P,P)=\mathsf
E^2/c^2+|\mathbf q|^2-|\bp|^2$ (r1-9).
\end{proof}

The form \eqref{r1:eq:Lfree} is inherited \cite{LandauLifshitz2}; (iii) is restated from \crefrange{prop:toy}{prop:orient} and \cref{ex:tachyons}, where the surplus momenta $\mathbf q$ produce the unstable band and the first-class constraints $q_a\approx0$ (\cref{prop:constraints},
\cref{thm:wlgauge}) remove it. This constraint is Option~II of \cref{r2:prop:shellII}.

The interior Lagrangian of \cref{ch:ch04_rigidbody} is the split-form Lagrangian \eqref{eq:splitL}. Applied to the length-valued $Y$ it reads
\begin{equation}
\label{r1:eq:Lkappa}
L_\kappa=\frac m2\,\eta\bigl(Y'+\kappa GY,\,Y'+\kappa GY\bigr),\qquad \kappa=\frac c\ell,\quad [\tau]=\mathrm s,\quad [m]=\mathrm{kg},\quad [L_\kappa]=\mathrm J,
\end{equation}
with $Y'=dY/d\tau$; the sum $Y'+\kappa GY$ is homogeneous in dimension only because $G$ acts on the converted coordinates (\eqref{r1:eq:GC},
r1-2, r1-36).

\begin{proposition}[The interior Lagrangian on lengths]
\label{r1:prop:shift}
\begin{enumerate}[label=(\roman*),leftmargin=2.2em]
\item $p=\partial L_\kappa/\partial Y'=m\eta(Y'+\kappa GY)$, $\pi:=\eta^{-1}p=m(Y'+\kappa GY)$, and the Euler--Lagrange equation is
$Y''=-\kappa^2Y$ on all six components; on shell $p'=\kappa G^{\!\top}p$ and $\pi'=\kappa G\pi$.
\item $L_\kappa=\tfrac m2\bigl(\eta(Y',Y')-\kappa^2\eta(Y,Y)\bigr)+\tfrac{d}{d\tau}\Lambda$ with
$\Lambda=\tfrac{m\kappa}2\eta(Y,GY)=-m\kappa\sum_a(ct^a)\,x^a$, and $p-m\eta Y'=\partial\Lambda/\partial Y$.
\item The energy $\mathcal E=p\cdot Y'-L_\kappa=\tfrac m2\bigl(\eta(Y',Y')+\kappa^2\eta(Y,Y)\bigr)$ is conserved and equals
$H(Y,p)=\tfrac1{2m}\eta^{-1}(p,p)-\kappa\,p^{\!\top}GY$, which generates the flow. Explicitly
$\mathcal E=\tfrac m2\sum_i\bigl((c\dot t^i)^2+\kappa^2(ct^i)^2\bigr)-\tfrac m2\sum_a\bigl((\dot x^a)^2+\kappa^2(x^a)^2\bigr)$: the
three converted-time oscillators carry positive and the three length oscillators negative energy, and both $\mathcal E$ on $(Y,Y')$ and
$H$ on $(Y,p)$ have signature $(6,6)$.
\end{enumerate}
\end{proposition}

\begin{proof}
(i) Differentiation and $G^\dagger=G$ (r1-10). (ii) $\eta(GY,GY)=-\eta(Y,Y)$, $\tfrac{d}{d\tau}\eta(Y,GY)=2\eta(Y',GY)$ because $\eta G$ is
symmetric, and $\eta(Y,GY)=-2\sum_iy_iy_{i+3}$ (r1-11). (iii) Substitution and Hamilton's equations; the inertia of the two $12\times12$
Hessians is computed exactly (r1-12).
\end{proof}

Bounded solutions do not make the energy bounded below. Every solution oscillates with the single frequency $\kappa$ because the system is linear and its
two signs of energy are decoupled; $\mathcal E$ and $H$ have signature $(6,6)$, so once the length oscillators are coupled to a system of positive
energy, or quantised with a positive norm, they are ghosts whose excitation lowers the energy. The surplus constraints of \cref{r2:sec:main} concern the
free worldline and do not remove these interior modes. Which constraints of the coupled interior system, if any, remove the negative-energy
oscillators, and which modes are physical, is not derived here (X-04).

Item (i) is restated (equations~\eqref{eq:splitL} and~\eqref{eq:splitEg}, and \cref{thm:neutral}); in particular the
raised momentum obeys $\pi'=\kappa G\pi$ for constant $G$. Item (ii) is \cref{thm:EBsplit} read on lengths: the
boundary term that is $-m\kappa\,\bE\cdot\bB$ in the field reading is $-m\kappa\sum_a(ct^a)x^a$, the correlation of converted time and
length along each paired axis. Item (iii) is fixed; the indefinite sign is that of \cref{bf:krein}.

\begin{proposition}[The spatial triple is not invariant]
\label{r1:prop:mix}
Every solution is $Y(\tau)=e^{-\kappa G\tau}Y_0+\frac{\sin\kappa\tau}{\kappa m}\pi_0$, $\pi(\tau)=e^{\kappa G\tau}\pi_0$, with
$e^{sG}=\cos s\,\id+\sin s\,G$.
\begin{enumerate}[label=(\roman*),leftmargin=2.2em]
\item Phase-space data of lengths and spatial momenta, $Y_0=\sum_ax_0^ae_{a+3}$ and $\pi_0\in V^-$, acquire the temporal part
$\sin(\kappa\tau)\sum_ax_0^ae_a$: lengths are rotated into converted times at once.
\item Configuration data $(Y_0,Y'_0)\in V^-\times V^-$ stay in $V^-$, since $Y''=-\kappa^2Y$ is diagonal, but their canonical momentum
has the temporal part $\pi(0)^+=-m\kappa\sum_ax^a_0e_a$, nonzero for $Y_0\neq0$, and $\pi'=\kappa G\pi$ maps $V^-$ into $V^+$; for example
$\pi(0)=e_5$ gives
$\pi(\tau)=\cos(\kappa\tau)e_5-\sin(\kappa\tau)e_2$.
\item The zero-momentum orbits $e^{-\kappa G\tau}Y_0$ are duality rotations; $e^{-\kappa G\tau}e_4=\cos(\kappa\tau)e_4+\sin(\kappa\tau)e_1$.
\end{enumerate}
\end{proposition}

\begin{proof}
The formula solves $Y''=-\kappa^2Y$ and $m(Y'+\kappa GY)=\pi$ with $Y(0)=Y_0$; the components are read off (r1-13).
\end{proof}

This is fixed. Restricting $X_{\rm cl}$ to the spatial triple is therefore not preserved by the interior dynamics as a statement about phase
space, while it is preserved as a statement about configurations. \Cref{r2:prop:leak} extends the rank of the leakage of the example in (ii) to every clock. Its reading needs one correction. By (ii) of \cref{r1:prop:shift} the rotation of $\pi$ out of a subspace
is carried by the gradient $m\kappa\eta GY$ of a boundary term, which a canonical transformation removes (\cref{thm:EBsplit}); the
kinetic momentum $mY'$ obeys $mY''=-m\kappa^2Y$ and preserves $V^-$ and every four-plane $E_{\hat n}$. The obstruction of \cref{r2:prop:nogo} and of the leakage is therefore an obstruction to taking the canonical $\pi$ as the observable momentum, not a property of the
motion. Neither $\pi$ nor $mY'$ is conserved, so neither is a four-momentum.

% --- source: two_readings/r1_lagrangians#5 (Lagrangians for the energies)
\subsection{Lagrangians for the energies}
\label{r1:sec:energies}

In this subsection units are Heaviside--Lorentz with $c=1$, $A_\mu$ is covariant, $\mathbf a=(A_1,A_2,A_3)$, and
$F_{\mu\nu}=\partial_\mu A_\nu-\partial_\nu A_\mu$ with $F_{0i}=E_i$ and $F_{jk}=\epsilon_{ijk}B_i$ (\cref{thm:emfield}). In SI
every density below is multiplied by $\varepsilon_0$ with $\bB\to c\bB$, which makes it a density in $\mathrm{J\,m^{-3}}$ (r1-36).

\begin{proposition}[Maxwell's Lagrangian in the field reading]
\label{r1:prop:maxwell}
\begin{enumerate}[label=(\roman*),leftmargin=2.2em]
\item $\bE=\partial_t\mathbf a-\nabla A_0$, $\bB=\curl\mathbf a$, and
$L_M=-\tfrac14F_{\mu\nu}F^{\mu\nu}=\tfrac12(|\bE|^2-|\bB|^2)=-\tfrac12\eta(X_{\rm EM},X_{\rm EM})=-\tfrac12\Re q$; in SI,
$L_M=\tfrac{\varepsilon_0}2(|\bE|^2-c^2|\bB|^2)$.
\item The Euler--Lagrange equations of $L_M-A_\mu j^\mu$ are $\partial_\mu F^{\mu\nu}=j^\nu$: $\nabla\cdot\bE=\rho$ and
$\partial_t\bE=-\curl\bB-\bj$ (in textbook variables $\partial_t\bE=\curl\Btext-\bj$), together with the identities $\partial_t\bB=\curl\bE$,
$\nabla\cdot\bB=0$.
\item The symmetric tensor $\Theta^{\mu\nu}=F^{\mu\lambda}F_\lambda{}^\nu+\tfrac14\eta^{\mu\nu}F_{\alpha\beta}F^{\alpha\beta}$ is traceless, with
$\Theta^{00}=E_6$ and $\Theta^{0i}=-(\bE\times\bB)_i=(\bE\times\Btext)_i$: the momentum density is $-\mathbf S$ in the paper's variables. It obeys
$\partial_\mu \Theta^{\mu\nu}=-F^\nu{}_\lambda\,\partial_\mu F^{\mu\lambda}$ identically, and $\Theta^\mu{}_\lambda T^\lambda{}_\nu=E_0^2\delta^\mu_\nu$,
so that $E_0=\tfrac12\sqrt{T_{\mu\nu}\Theta^{\mu\nu}}$.
\item The second invariant $\bE\cdot\bB=\tfrac12\Im q=-\bE\cdot\Btext$ is the divergence $\bE\cdot\bB=-\partial_\mu K^\mu$ of the
Chern--Simons current $K^\mu=\tfrac12\varepsilon^{\mu\nu\rho\sigma}A_\nu\partial_\rho A_\sigma$ ($\varepsilon_{0123}=+1$), with
$K^0=-\tfrac12\mathbf a\cdot\curl\mathbf a$; its Euler--Lagrange expression vanishes identically, so $\theta\,\bE\cdot\bB$ changes no field
equation.
\end{enumerate}
\end{proposition}

\begin{proof}
Symbolic differentiation on general $A_\mu(t,\mathbf x)$ and $j^\mu(t,\mathbf x)$ (r1-20, r1-21, r1-23); (iii) on general $(\bE,\bB)$ (r1-22).
\end{proof}

All of this is inherited \cite{Jackson,LandauLifshitz2,Rainich1925}; the Rainich identity and $E_6^2=E_0^2+|\mathbf S|^2$ are restated from
\crefrange{en:split}{en:rainich}. Two signs are fixed by the paper's orientation and are easy to lose: the second invariant is
$+\tfrac12\Im q$ in the paper's variables and $-\tfrac12\Im q$ in textbook ones, and the momentum density $\Theta^{0i}$ is $-\mathbf S$, not
$\mathbf S$. Maxwell's Lagrangian is the $\eta$-form of $X_{\rm EM}$, a member of the Lorentz class of \cref{thm:inertiareq}.

\paragraph{The rest energy as a Lagrangian.}
Write $S=\tfrac12(|\bE|^2-|\Btext|^2)$ and $P=\bE\cdot\Btext$, so that $E_0=\tfrac12|q|=\sqrt{S^2+P^2}$, and set $L_0:=E_0$.
For a Lagrangian $L(S,P)$ put $\mathbf D=\partial L/\partial\bE$ and $\mathbf H=-\partial L/\partial\Btext$.

\begin{theorem}[The rest-energy Lagrangian]
\label{r1:thm:L0}
\begin{enumerate}[label=(\roman*),leftmargin=2.2em]
\item $L_0$ is Lorentz invariant (annihilated by $\lG$), homogeneous of degree two, and invariant pointwise under the duality rotation
$z\mapsto e^{i\theta}z$, which does not leave $L_M=-\tfrac12\Re q$ invariant. It is not differentiable at null fields: along
$\Btext=(1,0,0)$, $\bE=(0,e,0)$ it equals $\tfrac12|1-e^2|$, with one-sided derivatives $-1$ and $+1$ at $e=1$.
\item $\mathbf D=(S\bE+P\Btext)/E_0$, $\mathbf H=(S\Btext-P\bE)/E_0$, and the Hamiltonian density is
$\mathbf D\cdot\bE-L_0=\tfrac S{E_0}E_6=-\cos(\arg q)\,E_6$, which ranges over $[-E_6,E_6]$ and equals $-\tfrac12$ on the magnetic field
$\Btext=(1,0,0)$, $\bE=0$, where $L_0$ is not convex in $\bE$.
\item For every $L(S,P)$ the gauge-invariant tensor $\Theta^\mu{}_\nu=\frac{\partial L}{\partial(\partial_\mu A_\lambda)}F_{\nu\lambda}-\delta^\mu_\nu L$
equals $L_ST_M{}^\mu{}_\nu+\delta^\mu_\nu(SL_S+PL_P-L)$ and obeys $\partial_\mu \Theta^\mu{}_\nu=-\mathcal E^\lambda F_{\nu\lambda}$, with
$\mathcal E^\lambda$ the Euler--Lagrange expression of $A_\lambda$. Since $L_0$ is homogeneous of degree one in $(S,P)$, its tensor
$T_0=(S/E_0)T_M$ is traceless: $L_0$ is conformal.
\item For any $L(S,P)$, $\mathbf D\cdot\mathbf H-\bE\cdot\Btext=P(L_S^2-L_P^2-1)-2SL_SL_P$. The defect is $-2P$ for $L_0$ and
$-(1+\lambda^2)P$ for $\lambda L_0$: no multiple of $E_0$ is a duality-invariant Lagrangian in the sense of Gaillard and Zumino.
\item The ModMax Lagrangians $L_\gamma=\cosh\gamma\,S+\sinh\gamma\,E_0=\cosh\gamma\,L_M+\sinh\gamma\,L_0$ have vanishing defect for every
$\gamma$, are conformal, and for $\gamma\ge0$ have $L_S\ge e^{-\gamma}>0$ and $\partial_\gamma L_\gamma=\tfrac12\sqrt{T^{(\gamma)}_{\mu\nu}T^{(\gamma)\mu\nu}}$;
in particular $L_0=\partial_\gamma L_\gamma|_{\gamma=0}=\tfrac12\sqrt{T^M_{\mu\nu}T^{M\mu\nu}}$. As $\gamma\to\infty$,
$L_\gamma/\sinh\gamma\to S+E_0=L_M+L_0$, whose defect is $-P$.
\end{enumerate}
\end{theorem}

\begin{proof}
(i) Under $L_i,K_i$ the variations of $S$ and $P$ vanish; $|q|$ is invariant under $q\mapsto e^{2i\theta}q$ (r1-24). (ii) Differentiation
and $\mathbf D\cdot\bE-L_0=(SE^2+P^2-S^2-P^2)/E_0$ with $E^2-S=E_6$; the Hessian in $\bE$ at the stated point has a negative eigenvalue
(r1-25). (iii) With $G^{\mu\lambda}=\partial L/\partial(\partial_\mu A_\lambda)=L_S\partial S/\partial(\partial_\mu A_\lambda)+L_P\partial
P/\partial(\partial_\mu A_\lambda)$ the identity is polynomial in $(\partial A,L_S,L_P,L)$; conservation uses
$G^{\mu\lambda}\partial_\mu F_{\nu\lambda}=\tfrac12G^{\mu\lambda}\partial_\nu F_{\mu\lambda}$, which holds for antisymmetric $G$ by the
Bianchi identity (r1-27). (iv) $\mathbf D=L_S\bE+L_P\Btext$, $\mathbf H=L_S\Btext-L_P\bE$ (r1-25, r1-26). (v) $L_S=\cosh\gamma+\sinh\gamma\,(S/E_0)$
with $S/E_0\in[-1,1]$; by (iii) and \cref{r1:prop:maxwell}(iii), $\tfrac12\sqrt{T^{(\gamma)}\!\cdot T^{(\gamma)}}=L_SE_0=\partial_\gamma L_\gamma$
(r1-26, r1-28).
\end{proof}

The general defect formula in (iv) is restated from \cref{sec:orth}; every other computation in the theorem is fixed. The
duality invariance and conformality of ModMax are inherited
\cite{BandosLechnerSorokinTownsend20,Kosyakov20}, as is the Gaillard--Zumino criterion \cite{GaillardZumino81,GibbonsRasheed95} (\cref{sec:orth}) and the generation of ModMax by the root-$T\bar T$ flow \cite{BabaeiAghbolagh22}; the scripts re-derive them. Calling $L_0$
``the rest-energy Lagrangian'' is interpretive, and the theorem says what the name costs. $E_0$ is a pointwise invariant, the energy density in the
local rest frame of a non-null field; the invariant mass of a localised configuration is the Lorentz length of the integrated $T^{0\mu}$, which is at least
$\int E_0$, with equality only when all points share one rest frame. The mass $m$ of the Dirac Lagrangian below is an input, not generated (X-05). $L_0$ is an admissible Lagrangian, Lorentz invariant
and conformal, but it is not duality invariant although $E_0$ is pointwise duality invariant; its energy density is indefinite; and it is
singular exactly on the null fields that carry radiation. A positive-energy field theory with Lagrangian $\lambda E_0$ is obstructed. What the
rest energy is, in the Lagrangian formalism, is the operator that deforms Maxwell's theory into ModMax: $L_0=\partial_\gamma L_\gamma|_0$,
equal to $\tfrac12\sqrt{T_{\mu\nu}\Theta^{\mu\nu}}$ by Rainich's identity. The limit $\gamma\to\infty$ does not isolate it but gives $L_M+L_0$.

\paragraph{Hamiltonians.} In Hamiltonian variables $(\mathbf D,\Btext)$ put $w=\Btext+i\mathbf D$, $s=E_6(w)=\tfrac12(|\mathbf D|^2+|\Btext|^2)$,
$p=|\mathbf D\times\Btext|=|\mathbf S(w)|$ and $E_0(w)=\tfrac12|w\cdot w|=\sqrt{s^2-p^2}$, the last by the identity of \cref{en:split}. For $\mathcal H(\mathbf D,\Btext)$ the field equations are $\dot{\mathbf D}=\curl\mathbf H$, $\dot\Btext=-\curl\bE$ with
$\bE=\partial\mathcal H/\partial\mathbf D$, $\mathbf H=\partial\mathcal H/\partial\Btext$; the energy flux is $\bE\times\mathbf H$, the
momentum density $\mathbf D\times\Btext$, and Lorentz invariance requires $\bE\times\mathbf H=\mathbf D\times\Btext$; duality invariance is
$\bE\cdot\Btext=\mathbf D\cdot\mathbf H$ and conformal invariance $\mathbf D\cdot\bE+\Btext\cdot\mathbf H=2\mathcal H$
\cite{BialynickiBirula83,BandosLechnerSorokinTownsend20}.

\begin{proposition}[Which energies are Hamiltonians]
\label{r1:prop:ham}
For $\mathcal H=a\,s+b\,E_0(w)$ one has $\bE\times\mathbf H=(a^2-b^2)\,\mathbf D\times\Btext$, and $\mathcal H$ is duality invariant and
conformal for all $(a,b)$. Hence $\mathcal H$ is Lorentz invariant exactly when $a^2-b^2=1$; with $a>0$ these are
$\cosh\gamma\,E_6-\sinh\gamma\,E_0$, the Maxwell Hamiltonian $s=E_6$ at $\gamma=0$ and the ModMax Hamiltonians for $\gamma>0$. The rest energy alone, $\mathcal H=E_0$, gives $\bE\times\mathbf H=-\mathbf D\times\Btext$
and is not Lorentz invariant. The Bialynicki-Birula Hamiltonian $\mathcal H=p=|\mathbf S(w)|=\sqrt{E_6^2-E_0^2}$ satisfies all three
conditions.
\end{proposition}

\begin{proof}
With $\xi=\tfrac12(|\mathbf D|^2-|\Btext|^2)$, $\eta_w=\mathbf D\cdot\Btext$ and $E_0=\sqrt{\xi^2+\eta_w^2}$,
$\bE=a\mathbf D+b(\xi\mathbf D+\eta_w\Btext)/E_0$ and $\mathbf H=a\Btext+b(\eta_w\mathbf D-\xi\Btext)/E_0$; the cross terms cancel and the
$b^2$ term is $-(\xi^2+\eta_w^2)\mathbf D\times\Btext/E_0^2$ (r1-29).
\end{proof}

The classification of conformal duality-invariant Hamiltonians, ModMax and Bialynicki-Birula's, is inherited
\cite{BandosLechnerSorokinTownsend20,BialynickiBirula83}. What is fixed here is its form in the variables of \cref{sec:energy}: the differentiation
$E_6^2=E_0^2+|\mathbf S|^2$ of the undivided energy supplies all three duality invariants, the admissible Hamiltonians are $\cosh\gamma\,E_6-\sinh\gamma\,E_0$
and $|\mathbf S|$, and the rest energy is admissible only together with $E_6$, never alone.

\begin{proposition}[The undivided energy is a Hamiltonian, not a Lagrangian]
\label{r1:prop:E6}
\begin{enumerate}[label=(\roman*),leftmargin=2.2em]
\item Used as a Lagrangian of the potential, $E_6=\tfrac12(|\bE|^2+|\bB|^2)$ with $\bE=\partial_t\mathbf a-\nabla A_0$, $\bB=\curl\mathbf a$
gives $\nabla\cdot\bE=0$ and $\partial_t\bE=+\curl\bB$; a transverse plane wave then needs $\omega^2+k^2=0$. The equations are elliptic.
\item The first-order action $L_H=\bE\cdot\partial_t\mathbf a-E_6+A_0\nabla\cdot\bE$, with $\bE$ independent, has Euler--Lagrange equations
$\bE=\partial_t\mathbf a-\nabla A_0$, $\partial_t\bE=-\curl\bB$, $\nabla\cdot\bE=0$, and on the first of them
$L_H=L_M+\nabla\cdot(A_0\bE)$: $E_6$ is the Legendre transform of $L_M$, with $\bE$ conjugate to $\mathbf a$.
\item With two potentials, $\bB=\curl\mathbf A$ and $\bE=\curl\mathbf C$, $X_A=(\mathbf A;\mathbf C)$, the Schwarz--Sen Lagrangian
$L_{SS}=\tfrac12(\curl X_A)\cdot G\dot X_A-\tfrac12|\curl X_A|^2=\tfrac12(\bE\cdot\dot{\mathbf A}-\bB\cdot\dot{\mathbf C})-E_6$ is first order in
time with Hamiltonian $E_6$, is exactly invariant under $X_A\mapsto e^{\theta G}X_A$, and the duality charge density is
$-\tfrac12X_A\cdot\curl X_A$, minus the optical helicity density. $L_{SS}-(\bE\cdot\dot{\mathbf A}-E_6)$ has identically vanishing
Euler--Lagrange expression.
\end{enumerate}
\end{proposition}

\begin{proof}
Symbolic Euler--Lagrange equations on general fields, with the plane wave $\mathbf a=(0,e^{i(kx-\omega t)},0)$ and $A_0=0$ in (i)
(r1-30, r1-31, r1-32).
\end{proof}

Item (i) is fixed, (ii) inherited \cite{HenneauxTeitelboim92}, and (iii) restated from equation~\cref{eq:SSL} and \cref{thm:SS}, the
duality-symmetric form of Deser and Teitelboim and of Schwarz and Sen \cite{DeserTeitelboim76,SchwarzSen94}. In (iii) duality is a symmetry
off shell because both terms are $G$-invariant: $E_6=\tfrac12|z|^2$ has duality weight zero, and $G$ is orthogonal for the Euclidean product
of the symplectic term. This is the field form of the dichotomy of \cref{thm:inertiareq} and \cref{en:unique}: the Lagrangian
$-\tfrac12\eta(X,X)$ lies in the Lorentz class, the Hamiltonian $E_6$ in the duality class, and the two classes meet in $0$, so no quadratic
form is both. $E_0$ enters neither $L_M$ nor $L_H$ nor $L_{SS}$; it enters as the deformation of \cref{r1:thm:L0}(v) and of
\cref{r1:prop:ham}.

% --- source: two_readings/r1_lagrangians#6 (The total Lagrangian)
\subsection{The total Lagrangian}
\label{r1:sec:total}

With $\hbar=c=1$ the total Lagrangian of the field theory and of a classical charged test particle is
\begin{equation}
\label{r1:eq:Ltot}
\begin{gathered}
\mathcal L_{\rm tot}=-\tfrac14F_{\mu\nu}F^{\mu\nu}+\tfrac12\psi^{\!\top}\bigl(K^\mu D_\mu-mT\bigr)\psi,\\
L_{\rm pp}=-M\sqrt{\eta(Y',Y')}-e\,A_\mu(x)\,x'^\mu,\qquad x=(Y^1;Y^4,Y^5,Y^6).
\end{gathered}
\end{equation}

\begin{proposition}[Coupled equations and conservation]
\label{r1:prop:total}
\begin{enumerate}[label=(\roman*),leftmargin=2.2em]
\item $\partial\mathcal L_{\rm tot}/\partial A_\nu=q\,j^\nu$, so that Maxwell's equations read $\partial_\mu F^{\mu\nu}=J_e^\nu:=-q\,j^\nu$;
the Dirac equation is $K^\mu D_\mu\psi=mT\psi$.
\item The gauge-covariant tensor $\Theta_D{}^\mu{}_\nu=\tfrac12\psi^{\!\top}K^\mu D_\nu\psi-\delta^\mu_\nu L_D$ satisfies
$\partial_\mu\Theta_D{}^\mu{}_\nu=qF_{\lambda\nu}j^\lambda$ on the Dirac shell, and $\partial_\mu T_M{}^\mu{}_\nu=-F_{\nu\lambda}\partial_\mu F^{\mu\lambda}$,
which equals $qF_{\nu\lambda}j^\lambda$ on the Maxwell shell. Hence $\partial_\mu(T_M+\Theta_D)^\mu{}_\nu=0$ on the coupled shell.
\item For $L_{\rm pp}$ the surplus momenta $K_2,K_3$ are conserved, and the interior equations are the Lorentz force
$\frac{d}{d\lambda}\bigl(Mx'_\mu/\sqrt{\eta(Y',Y')}\bigr)=eF_{\mu\nu}x'^\nu$.
\end{enumerate}
\end{proposition}

\begin{proof}
(i)--(ii) in the Grassmann jet algebra with coefficients depending on general $A_\mu(t,\mathbf x)$, the time derivatives of $\psi$ being
eliminated with the Dirac equation (r1-33, r1-34). (iii) Symbolic Euler--Lagrange equations for general $A_\mu(x)$ (r1-35).
\end{proof}

The coupled statement is inherited \cite{WeinbergQFT}; in the arena it is the one-constituent version of \cref{thm:coupled}, with
the Grassmann kernels of \eqref{r1:eq:LD} in place of the commuting ones there, and the checks make it fixed for this realisation.  The four-momentum of the field is $P^\mu=\int T^{0\mu}_{\rm tot}\,d^3x$, developed in \cref{r2:sec:field}.

Each reading enters $\mathcal L_{\rm tot}$ in its own place. The classical reading is the argument: the coordinates $x^\mu=(ct^1;x^a)$ on which
$A_\mu$ and $\psi$ depend, and the worldline $Y(\lambda)$, whose surplus components $ct^2,ct^3$ couple to nothing but the six-dimensional
mass shell. The field reading is a value: $X_{\rm EM}$ is the field strength $F=dA$ of the interior potential, entering through
$L_M=-\tfrac12\eta(X_{\rm EM},X_{\rm EM})$, through $\Theta^{\mu\nu}$ and through the Lorentz force, and it becomes a vector of $\R^{3,3}$ with
the action of $\lG$ only through the two-parameter dictionary of \cref{r1:thm:type}(iii). The two meet only in the derivative
$\partial_\mu$, of dimension $\mathrm m^{-1}$ in the classical reading, which carries the potential to the field strength.

\begin{table}[t]
\centering
\footnotesize
\setlength{\tabcolsep}{3pt}
\renewcommand{\arraystretch}{1.15}
\begin{tabular}{@{}>{\raggedright\arraybackslash}p{1.45cm}>{\raggedright\arraybackslash}p{1.95cm}>{\raggedright\arraybackslash}p{1.25cm}%
>{\raggedright\arraybackslash}p{2.2cm}>{\raggedright\arraybackslash}p{2.9cm}>{\raggedright\arraybackslash}p{2.35cm}>{\raggedright\arraybackslash}p{2.4cm}@{}}
\toprule
symbol & reading & SI unit & natural Lorentz action & Lagrangian & canonical conjugate & Noether charges\\
\midrule
$Y=X_{\rm cl}$ & displacement; lengths on $V^-$ & m & $\lint$ on $(ct^1;x)$ after the clock; $ct^2,ct^3$ scalars & $-Mc\sqrt{\eta(Y',Y')}$ &
$K_A$, $\mathrm{kg\,m\,s^{-1}}$ & $K_A$ (translations), $\so{3,3}$ charges $Y_AK_B-Y_BK_A$\\
$Y$, interior & configuration & m & none preserving a Lorentzian four-plane & $\tfrac m2\eta(Y'+\kappa GY,Y'+\kappa GY)$ & $p=m\eta(Y'+\kappa GY)$ & energy, signature $(6,6)$\\
$A_\mu=(A_0;\mathbf a)$ & potential & $\mathrm{V\,s\,m^{-1}}$ ($A_0=\phi/c$) & $\lint$, $(\tfrac12,\tfrac12)$ & $L_M$ through $F=dA$ & $\varepsilon_0\bE$ conjugate to $\mathbf a$ & $P^\mu=\int T^{0\mu}d^3x$\\
$X_{\rm EM}=(c\bB;\bE)$ & field strength & $\mathrm{V\,m^{-1}}$ & $\lG$, $(1,0)\oplus(0,1)$ & $L_M=-\tfrac{\varepsilon_0}2\eta(X_{\rm EM},X_{\rm EM})$ &
none: a phase-space point & observables $E_6$, $\mathbf S$, $q$, $\Theta^{\mu\nu}$\\
$X_A=(\mathbf A;\mathbf C)$ & dual potentials & $\mathbf A$: $\mathrm{V\,s\,m^{-1}}$, $\mathbf C$: V & boosts only on shell; duality $e^{\theta G}$ & $L_{SS}$ & $\tfrac12G^{\!\top}\curl X_A$ (second class) & energy $\int E_6$, helicity\\
$\psi$ & Dirac field & $\mathrm m^{-3/2}$ & $\lint$, $(\tfrac12,0)\oplus(0,\tfrac12)$ & $\tfrac12\psi^{\!\top}(K^\mu D_\mu-mT)\psi$ & $\tfrac12\psi^{\!\top}$ (second class) & $\int j^0$, $P^\mu$\\
$E_6$ & undivided energy & $\mathrm{J\,m^{-3}}$ & $\Theta^{00}$, a tensor component & Hamiltonian of $L_M$, $L_{SS}$; not a Lagrangian & -- & generates time translation\\
$E_0$ & rest energy & $\mathrm{J\,m^{-3}}$ & scalar, duality invariant & $L_0$ (interpretive); ModMax deformation & -- & --\\
\bottomrule
\end{tabular}
\caption{The objects of \cref{r1:sec:main}: reading, unit, Lorentz action, Lagrangian, conjugate and charges.}
\label{r1:tab:objects}
\end{table}

% --- source: two_readings/r1_lagrangians#7 (Answer)
\subsection{Answer}
\label{r1:sec:answer}

The Lagrangians asked for in the introduction, and the status of each statement, are the following.
\begin{enumerate}[label=(\arabic*),leftmargin=2.2em]
\item \emph{The two $X$'s.} $X_{\rm cl}=Y$ is a displacement in metres whose length-valued part lies on the spatial triple, the temporal
part being times converted by the unique $c$; $X_{\rm EM}=(c\bB;\bE)$ is a field strength in volts per metre. They share $\R^{3,3}$, $\eta$
and the Clifford action, and nothing else: no nonzero linear equivariant map joins them for any identification of their Lorentz algebras,
and in SI units no monomial in $c$, $\ell$, $\hbar$ supplies the unit $\mathrm{m^2V^{-1}}$ such a map would need (\cref{r1:thm:type};
fixed, the first statement obstructing a dictionary; the second holds only in SI units, since in Heaviside--Lorentz units with
$\hbar=c=1$ the monomial $\ell^3$ supplies the conversion). They are kept apart by treating $X_{\rm cl}$ as the argument and $X_{\rm EM}$ as a value of the fields
(\cref{r1:sec:total}; interpretive).
\item \emph{Lagrangians for $X_{\rm cl}$.} The free worldline \eqref{r1:eq:Lfree} (inherited form; its split-signature consequences restated)
and the interior Lagrangian \eqref{r1:eq:Lkappa} on lengths, with $\kappa=c/\ell$, Euler--Lagrange equation $Y''=-\kappa^2Y$, $\pi'=\kappa G\pi$
(restated) and an energy of signature $(6,6)$ (fixed). The interior flow does not preserve the spatial triple in phase
space (fixed); the obstruction concerns the canonical momentum, which differs from the kinetic one by a boundary-term gradient (fixed).
\item \emph{The Dirac Lagrangian.} \eqref{r1:eq:LD}, with the Grassmann variation rule (fixed), the field equation
$J\gamma^\mu\partial_\mu\psi=m\psi$ (restated), the current $\tfrac12\psi^{\!\top}K^\mu J\psi$ and the canonical and Belinfante tensors
(inherited, checked).
\item \emph{Lagrangians for the energies.} Maxwell's $L_M=-\tfrac12\eta(X_{\rm EM},X_{\rm EM})$ (inherited); $E_6$ is its Hamiltonian, in the
first-order and the duality-symmetric forms (inherited, restated), and as a Lagrangian it gives elliptic equations (fixed); $E_0$ is a
Lorentz-invariant conformal Lagrangian with indefinite energy, pointwise duality invariant but not duality invariant in the sense of
Gaillard and Zumino (fixed), and it is the generator of ModMax and the
partner of $E_6$ in the ModMax Hamiltonian (fixed; the ModMax facts inherited). Reading $L_0$ as ``the rest-energy Lagrangian'' is
interpretive, and as the Lagrangian of a positive-energy theory it is obstructed.
\item \emph{The total Lagrangian.} \eqref{r1:eq:Ltot}, with total conservation on the coupled shell (inherited, checked) and
\cref{r1:tab:objects}.
\item \emph{Open.} A dynamical relation between the two readings, which would have to be a new, dimensionful, covariant map, and the
quantisation of either reading in the arena's own terms, remain open.
\end{enumerate}

% --- source: two_readings/r2_fourmomentum#0 ((head))
% Track r2 -- The four-momentum with X on the spatial triple.
% New references: sections/bib_r2.tex (entries copied verbatim from tex/refs.tex).

% --- source: two_readings/r2_fourmomentum#1 (The four-momentum with X on the spatial triple)
\section{\texorpdfstring{The four-momentum with $X$ on the spatial triple}{The four-momentum with X on the spatial triple}}
\label{r2:sec:main}

The third question asked in the introduction is how the four-momentum is constructed when $X$ is taken in the three directions
$(-,-,-)$. We read the sign triple as the values $\eta(e_{a+3},e_{a+3})=-1$ of the metric on the spatial triple $V^-=\langle e_4,e_5,e_6\rangle$:
the classical configuration $X_{\rm cl}$ has its three lengths there, and its three remaining components are times converted by $c$ on the
temporal triple $V^+$. In brief:
(1) a clock $\hat n\in V^+$ selects the four-plane $E_{\hat n}=\R \hat n\oplus V^-$, and these are exactly the Lorentzian four-planes that contain the three
lengths; \cref{thm:i} holds at every clock, with complex structure $J(\hat n)=e_1e_2e_3\,\hat n$ (\cref{r2:sec:clock});
(2) the four-momentum of a particle is the projection $P_{\hat n}=(E/c,\bp)$ of the six-momentum, conjugate to $(ct_{\hat n},x^a)$ and conserved by every dynamics
that is translation invariant along $E_{\hat n}$; of the three mass shells the arena's metric fixes the sign of the ambient one, and the gauge
principle of \cref{sec:constraints} selects the reduced one (\cref{r2:sec:particle});
(3) rotating the clock reads a fixed six-momentum as a four-dimensional mass with $m_4(\hat n)^2\le\mu^2$, with equality exactly at the clocks
parallel to the temporal momentum (\cref{r2:sec:rotate});
(4) the field $X_{\rm EM}$ has no linear, and no quadratic, covariant map to a four-vector: its four-momentum is the integrated first row of the
stress tensor, a four-vector only for an isolated whole (\cref{r2:sec:field});
(5) the no-go for the binding root concerns $G$ on the vector module $E_{\hat n}$, where it leaks two of the three spatial directions into the surplus
plane at every clock; it does not concern $G=\star$ on the field module (\cref{r2:sec:nogo}).
\Cref{r2:sec:answer} collects the verdicts.

\paragraph{Conventions and units.} As in \cref{sec:arena}, \cref{sec:clock} and \cref{sec:emfield}: $V=\R^{3,3}$, $\eta=\operatorname{diag}(1,1,1,\allowbreak-1,-1,-1)$, the
realisation of equation~\cref{eq:realisation}, $s_a:=e_{a+3}$, the binding root $Ge_i=e_{i+3}$, $Ge_{i+3}=-e_i$ ($G^2=-\id$, $\eta$-self-adjoint, an
anti-isometry) and the para-complex $\Kc e_i=e_{i+3}$, $\Kc e_{i+3}=e_i$ ($\Kc^2=\id$, $\eta$-skew; this is $-\cK$ for the grading root $\cK$ of
\cref{prop:twoK}). A bivector $a\wedge b$ acts on $V$ by
$v\mapsto a\,\eta(b,v)-b\,\eta(a,v)=\tfrac12[ab,v]$. In the particle reading units are SI with $c$ explicit: times $t^i$ in $\mathrm s$, lengths $x^a$
and the coordinates $Y^A$ in $\mathrm m$, $[c]=\mathrm{m\,s^{-1}}$, momenta $p_a,q_a$ and $K$ in $\mathrm{kg\,m\,s^{-1}}$, energies in $\mathrm J$,
masses $M_6,m_4,\mu$ in $\mathrm{kg}$; clocks $\hat n$ and frames are dimensionless. In the field reading the computations use Heaviside--Lorentz units with
$c=1$, with $\bB$ measured as $c\bB$; the SI form is restored through $z_c=c\bB+i\bE$, $[\bE]=[c\bB]=\mathrm{V\,m^{-1}}$ (\cref{en:SI}).
The paper's $\bB$ is minus the textbook magnetic field, $\Btext=-\bB$ (\cref{sec:conventions}); the frame Poynting vector is $\mathbf S=\bE\times\bB
=-\mathbf S_{\rm t}$, with $\mathbf S_{\rm t}=\bE\times\Btext$ the textbook one (equation~\cref{en:defs}). Identifiers r2-$\hat n$ refer to
\texttt{r2\_fourmomentum.py}; r2-0a--c check these conventions, the Fock modes $a^\pm_i=\tfrac12(e_i\pm e_{i+3})$ (CAR, $N$ counting the minus
signs of the pole operators $\omega_i=e_ie_{i+3}$, which are diagonal) and the generators $L_i,K_i$ of $\lG$.

% --- source: two_readings/r2_fourmomentum#3 (The four-momentum of a particle)
\subsection{The four-momentum of a particle}
\label{r2:sec:particle}

The four-momentum read by the clock $\hat n$ is the projection
\begin{equation}
\label{r2:eq:P}
P_{\hat n}:=\pr_{\hat n}K^\sharp=\frac Ecn+\sum_ap_as_a,\qquad P_{\hat n}^\mu=(E/c,\bp),\qquad P_{n,\mu}=(E/c,-\bp),
\end{equation}
whose three spatial components are the conjugates of the three lengths and whose time component is the conjugate of the clock's time.

\begin{proposition}[Four-vector and conservation]
\label{r2:prop:cons}
\begin{enumerate}[label=(\roman*),leftmargin=2.2em]
\item Under the boost $\exp(\chi\,s_1\wedge \hat n)$, extended by the identity on $U_{\hat n}$, the components $(E/c,\bp)$ transform by the $4\times4$ boost of
rapidity $\chi$ along the first spatial axis, an element of $SO^+(1,3)$.
\item The Hamiltonian flows of $E/c$ and of $-p_a$ translate $ct_{\hat n}$ and $x^a$ by one unit and leave every other coordinate fixed.
\item For every Hamiltonian $H(K_A,u^1,u^2)$ that does not depend on $(ct_{\hat n},x^a)$, $\{E/c,H\}=\{p_a,H\}=0$. Adding a term $V(ct_{\hat n})$ destroys only the
conservation of $E$; $q_a$ is conserved iff $H$ does not depend on $u^a$.
\end{enumerate}
\end{proposition}

\begin{proof}
(i) is \cref{r2:prop:cov}(ii) in the basis $(\hat n,s_a)$. (ii) is the bracket table above. (iii) $\{P_\mu,H\}=-\partial H/\partial x^\mu$ in the adapted
canonical coordinates.
\end{proof}

Fixed (r2-5c at $\hat n=(1,2,2)/3$; r2-5b for a symbolic frame and an arbitrary function $H$). The axiom alone does not make $E/c$ a Noether energy, since the axiom fixes no dynamics; (iii) supplies what is
missing: $E/c$ and $\bp$ are the generators of the translations of $E_{\hat n}$, hence the Noether charges of every dynamics invariant under them. 

\paragraph{Three mass shells.} The six-momentum is free until a dynamics is chosen. \Cref{r2:prop:shellI,r2:prop:shellII,r2:prop:shellIII} treat three options.

\begin{proposition}[Option I: the ambient shell]
\label{r2:prop:shellI}
For every clock, $\eta(K^\sharp,K^\sharp)=E^2/c^2-|\bp|^2+|\mathbf q|^2$. On the shell $\eta(K^\sharp,K^\sharp)=M_6^2c^2$,
\[
E^2=c^2|\bp|^2+M_6^2c^4-c^2|\mathbf q|^2,\qquad m_{\rm eff}^2=M_6^2-|\mathbf q|^2/c^2 ,
\]
so the four-dimensional mass vanishes at $|\mathbf q|=M_6c$ and is imaginary (tachyonic) beyond.
\end{proposition}

\begin{proof}
$(\hat n,s_a,u_a)$ is $\eta$-orthonormal with norms $(+1,-1,+1)$.
\end{proof}

Fixed (r2-6, symbolic frame). The sign of $|\mathbf q|^2$ is not a choice: the normal plane of a
four-plane of signature $(1,3)$ in $\R^{3,3}$ has the sign of time (\cref{thm:selector}(ii)). The option is \cref{en:temporal} in another notation.

\begin{proposition}[Option I is the momentum reading of the undivided energy]
\label{r2:prop:RE}
Let $K^\sharp=\sum_i(\mathsf E_i/c)e_i+\sum_ap_as_a$ with $\eta(K^\sharp,K^\sharp)=\mu^2c^2$, as in \cref{en:temporal}, and write
$\boldsymbol{\mathsf E}=(\mathsf E_1,\mathsf E_2,\mathsf E_3)$, in $\mathrm J$. For every clock $\hat n$,
\[
E=\hat n\cdot\boldsymbol{\mathsf E},\qquad c\,q_a=u_a\cdot\boldsymbol{\mathsf E},\qquad c^2|\mathbf q|^2=|\boldsymbol{\mathsf E}|^2-(\hat n\cdot\boldsymbol{\mathsf E})^2
=|\hat n\times\boldsymbol{\mathsf E}|^2,
\]
so Option I holds with $M_6=\mu$ and $m_{\rm eff}$ equal to the mass $m_4(\hat n)$ read by the clock. At $\hat n=\widehat W$, $c\,|\mathbf q|=|\boldsymbol{\mathsf E}_\perp|$
and $m_4^2c^4=\mu^2c^4-|\boldsymbol{\mathsf E}_\perp|^2$, which is the statement of \cref{en:temporal} that splitting the temporal energy lowers the mass. For
$\boldsymbol{\mathsf E}=(3,4,0)$, $\bp=(5,0,0)$, $c=1$: $\mu=0$, $m_4^2=-16$ at $e_1$ and $m_4^2=-\tfrac{26}3$ at $\widehat W$.
\end{proposition}

\begin{proof}
$E=c\,\eta(\hat n,K^\sharp)$, $q_a=\eta(u_a,K^\sharp)$, and $(\hat n,u_1,u_2)$ is orthonormal; the last identity is Lagrange's.
\end{proof}

Restated from \cref{en:temporal}, re-verified for a symbolic clock (r2-7). The surplus momenta $q_a$ are the relative-time energies of \cref{sec:energy} divided by $c$, and the tachyonic shell is realised by the example of \cref{sec:energy-temporal}. The obstruction stated there, that in the momentum reading the
differentiation of the temporal energy subtracts from the mass, is the same obstruction as the sign in \cref{r2:prop:shellI}.

\begin{proposition}[Option II: the first-class reduction]
\label{r2:prop:shellII}
In adapted canonical coordinates take the action
\[
S_{\rm red}=\int d\tau\Bigl(P_\mu\dot x^\mu+q_a\dot u^a-\tfrac N2\bigl(P_\mu P^\mu-M_6^2c^2+|\mathbf q|^2\bigr)-\lambda^aq_a\Bigr),
\]
with $[N]=\mathrm{kg^{-1}}$ and $[\lambda^a]=\mathrm{m\,s^{-1}}$ for $\tau$ in seconds. Then $\{q_a,q_b\}=0$ and $\{q_a,H\}=0$, so the constraints $q_a\approx0$ are
first class; $\varepsilon^aq_a$ generates $\delta u^a=\varepsilon^a$ and nothing else; $S_{\rm red}$ is exactly invariant under $\delta u^a=\varepsilon^a(\tau)$,
$\delta\lambda^a=\dot\varepsilon^a$; the reduced phase space is $T^*E_{\hat n}$, of dimension $12-2\times2=8$, with the shell $P_\mu P^\mu=M_6^2c^2$. With a normal
connection $A_\mu(x)$, $k_\mu=P_\mu-A_\mu\ell$ and $\ell=vp_u-up_v$ at $\hat n=e_1$, where $\mathbf q=(p_u,p_v)$, the brackets become $\{p_u,H\}=-Cp_v$, $\{p_v,H\}=Cp_u$ with $C=A_\mu
h^{\mu\nu}k_\nu$, and the constraints stay first class.
\end{proposition}

\begin{proof}
$H$ does not depend on $u$; the variation of the surplus part of the Lagrangian is $q_a\dot\varepsilon^a-\dot\varepsilon^aq_a=0$. For the last
statement, $H$ depends on $(u,v)$ only through $\ell$, with $\partial_u\ell=-p_v$ and $\partial_v\ell=p_u$, so $\{p_u,H\}=-\partial_uH=-Cp_v$ and
$\{p_v,H\}=-\partial_vH=Cp_u$, while $\{p_u,p_v\}=0$. The same computation for a curved four-metric is \cref{prop:constraints} in
\cref{ch:ch11_curvature}.
\end{proof}

At $A=0$ the reduction is that of \cref{thm:wlgauge} (\cref{ch:ch12_causal}) and its brackets are those of \cref{prop:constraints} (\cref{ch:ch11_curvature}); it is the implementation of \cref{assum:clockred} with the clock $e_1$ replaced by $\hat n$ (r2-8a; r2-8b re-verifies the brackets with
$A\neq0$). Option~II is the first-class surplus-momentum construction of \cref{sec:constraints}. Two consequences are restated from \cref{sec:constraints}: gauging only the rotation of $U_{\hat n}$ leaves the unstable band
(\cref{prop:isometryonly}), and in Dirac quantisation the physical states are functions on $E_{\hat n}$ obeying the Klein--Gordon equation (\cref{thm:wlgauge} and \cref{cor:weak}). By
\cref{r2:prop:RE}, $q\approx0$ means $\boldsymbol{\mathsf E}\parallel \hat n$: the reduction does not select a clock in $V^+$, it ties the temporal momentum
to whichever clock is chosen.

\begin{proposition}[Option III: a separate internal mass]
\label{r2:prop:shellIII}
For the action $S_{\hat n}=\int d\tau\,[P_\mu\dot x^\mu+q_a\dot u^a-\tfrac N2(P_\mu P^\mu-c^2\mathcal M^2(u,q))]$ with an arbitrary
function $\mathcal M^2$ (dimension $\mathrm{kg^2}$):
\begin{enumerate}[label=(\roman*),leftmargin=2.2em]
\item $\dot P_\mu=0$, $\tfrac{d}{d\tau}\mathcal M^2=0$, and the constraint is $E^2=c^2|\bp|^2+c^4\mathcal M^2$;
\item for the toy $\mathcal M^2=m_0^2+|\mathbf q|^2/c^2$ the surplus velocity is $\dot u^a=-Nq_a$, against $+Nq_a$ on the ambient shell, and the constraint
form is $P_\mu P^\mu-|\mathbf q|^2=\eta(R_{\hat n}K^\sharp,K^\sharp)$ with $R_{\hat n}=\pr_{\hat n}-\pr^{\perp}_{\hat n}=\exp(\pi\,u_1\wedge u_2)\in SO(3,3)$, a quadratic form of
signature $(1,5)$.
\end{enumerate}
\end{proposition}

\begin{proof}
(i) $H$ has no $x$-dependence, and $\{\mathcal M^2,H\}=-\tfrac{Nc^2}2\{\mathcal M^2,\mathcal M^2\}=0$. (ii) $\partial H/\partial q_a$; $R_{\hat n}$ is the
identity on $E_{\hat n}$ and $-\id$ on $U_{\hat n}$, the half-turn of $U_{\hat n}$.
\end{proof}

Option~III is constructed: $\mathcal M^2$ is supplied from outside the arena (r2-9a for an arbitrary
function; r2-9b at $\hat n=e_1,(2,3,6)/7$). The positive sign of the toy is opposite to the ambient one, and (ii) makes this precise: the toy is the free particle of the twisted form $\eta(R_{\hat n}\cdot,\cdot)$, which has one time and five space directions, so its sign is
that of surplus \emph{lengths}, the Kaluza--Klein sign, and not that of the arena's surplus times.

\begin{remark}[Which option the arena prefers]
\label{r2:rem:prefer}
The arena's metric fixes the sign of Option~I. For fields the same sign is the ill-posed, unbounded-below baseline of \cref{prop:toy}, whose
unstable band Option~II removes together with every other nonzero surplus momentum (\cref{fig:toy} and \cref{thm:wlgauge}; r2-8a). Option~III needs a
functional whose surplus sign is opposite to that of the metric; the clock supplies one, the twisted form $\eta(R_{\hat n}\cdot,\cdot)$ of
\cref{r2:prop:shellIII}(ii), but nothing in the arena selects it in place of $\eta$. So the arena prefers Option~II, as the only one of the three that has the arena's sign and no
instability; the preference is conditional, because the gauge principle of the normal translations is chosen and not derived (\cref{sec:constraints}). This is interpretive; whether a principle makes the normal translations gauge is open.
\end{remark}

% --- source: two_readings/r2_fourmomentum#4 (Rotating the clock)
\subsection{Rotating the clock}
\label{r2:sec:rotate}

Fix a six-momentum $K^\sharp=\sum_i(\mathsf E_i/c)e_i+\sum_ap_as_a$ with temporal energy vector $\boldsymbol{\mathsf E}$ and six-dimensional mass
$\mu^2c^4=|\boldsymbol{\mathsf E}|^2-|\bp|^2c^2$, and let the clock vary. The energy and the four-dimensional mass read by $\hat n$ are
\[
E(\hat n):=c\,\eta(K^\sharp,\hat n),\qquad m_4(\hat n)^2c^4:=E(\hat n)^2-|\bp|^2c^2 .
\]

\begin{proposition}[The mass read by a rotating clock]
\label{r2:prop:rotate}
\begin{enumerate}[label=(\roman*),leftmargin=2.2em]
\item On the clock sphere $E(\hat n)=\hat n\cdot\boldsymbol{\mathsf E}$ and
\[
m_4(\hat n)^2c^4=\mu^2c^4-|\hat n\times\boldsymbol{\mathsf E}|^2=\mu^2c^4-|\boldsymbol{\mathsf E}|^2\sin^2\angle(\hat n,\boldsymbol{\mathsf E}) .
\]
Hence $m_4(\hat n)^2\le\mu^2$, with equality iff $\hat n\parallel\boldsymbol{\mathsf E}$ (in either orientation), and $E(\hat n)\le|\boldsymbol{\mathsf E}|$, with equality iff
$\hat n=\boldsymbol{\mathsf E}/|\boldsymbol{\mathsf E}|$: the clock that reads the largest energy is the direction of the temporal part of $K$, and there $m_4=\mu$.
\item On the Fisher--Rao octant, $E(\hat n)=\sum_i\sqrt{w_i}\,\mathsf E_i$ and $m_4^2c^4=\bigl(\sum_i\sqrt{w_i}\,\mathsf E_i\bigr)^2-|\bp|^2c^2$. Let
$\boldsymbol{\mathsf E}^+$ be the componentwise positive part. If $\boldsymbol{\mathsf E}^+\neq0$, $\max E(\hat n)=|\boldsymbol{\mathsf E}^+|$, attained at
$w_i=(\mathsf E_i^+)^2/|\boldsymbol{\mathsf E}^+|^2$; if $\boldsymbol{\mathsf E}^+=0$, the maximum is $\max_i\mathsf E_i\le0$, at a vertex, and no clock of the
octant reads a positive energy. If $\boldsymbol{\mathsf E}$ has no negative component the maximising clock of (i) lies in the octant; otherwise, if $\boldsymbol{\mathsf E}^+\neq0$, the largest mass read by a future-reading
clock of the octant is $m_4^2c^4=\mu^2c^4-|\boldsymbol{\mathsf E}-\boldsymbol{\mathsf E}^+|^2<\mu^2c^4$.
\item At $\widehat W$, $E(\widehat W)=(\mathsf E_1+\mathsf E_2+\mathsf E_3)/\sqrt3$ and $m_4^2c^4=\mu^2c^4-|\boldsymbol{\mathsf E}_\perp|^2$, the $\mathsf E_W$ of \cref{sec:energy-temporal}.
\end{enumerate}
\end{proposition}

\begin{proof}
(i) $|\hat n|=1$ and Lagrange's identity $|\boldsymbol{\mathsf E}|^2=(\hat n\cdot\boldsymbol{\mathsf E})^2+|\hat n\times\boldsymbol{\mathsf E}|^2$. (ii) For $\hat n_i\ge0$,
$\hat n\cdot\boldsymbol{\mathsf E}\le \hat n\cdot\boldsymbol{\mathsf E}^+\le|\boldsymbol{\mathsf E}^+|$, with equality iff $\hat n=\boldsymbol{\mathsf E}^+/|\boldsymbol{\mathsf E}^+|$
when $\boldsymbol{\mathsf E}^+\neq0$; if $\boldsymbol{\mathsf E}^+=0$, then $\sum_in_i\ge|\hat n|=1$ gives $\hat n\cdot\boldsymbol{\mathsf E}\le\max_i\mathsf E_i\sum_in_i\le\max_i\mathsf E_i$,
with equality at the vertex of the largest component; $m_4^2$ increases with $E(\hat n)\ge0$, and $|\boldsymbol{\mathsf E}^+|^2=|\boldsymbol{\mathsf E}|^2-|\boldsymbol{\mathsf E}-\boldsymbol{\mathsf E}^+|^2$. (iii) is (i) at $\hat n=\widehat W$.
\end{proof}

Fixed (r2-10a symbolic; r2-10c a random test of (ii), 200 vectors $\boldsymbol{\mathsf E}$ against 4000 octant clocks each, in floating point with tolerance
$10^{-12}$; r2-10d the case $\boldsymbol{\mathsf E}^+=0$, exactly at $\boldsymbol{\mathsf E}=(-1,-2,-3)$ and by a random test). An exact example (r2-10b): $\boldsymbol{\mathsf E}=(2,3,6)$, $\bp=(0,0,3)$, $c=1$ give $\mu^2=40$; the aligned clock $(2,3,6)/7$ lies in the
Fisher--Rao octant with weights $(4,9,36)/49$ and reads $m_4^2=40$, as does the opposite clock, which reads the energy $-7$; the clock $e_1$ reads $m_4^2=-5$, and $\widehat W$ reads $\tfrac{94}3=40-|\boldsymbol{\mathsf E}_\perp|^2$.

\begin{remark}[The mass as a projection]
\label{r2:rem:projection}
Interpretive. In the six-momentum the temporal energy is one vector $\boldsymbol{\mathsf E}$, undivided; a clock reads its component $E(\hat n)$, and what the
reading leaves out, $c\,\pr^{\perp}_{\hat n}K^\sharp$, is the part orthogonal to the clock. The energy read by the clock is the projection of the undivided
temporal energy onto it; the four-dimensional mass computed from it is maximal, and equal to the six-dimensional $\mu$, only for the clocks parallel
to $\boldsymbol{\mathsf E}$, and is lowered by every misalignment, which is \cref{r2:prop:RE} read on the clock sphere. Nothing in this reading prefers a direction of $V^+$: the aligned clock is fixed relative to the momentum,
and for several particles read by one clock the reduction of \cref{r2:prop:shellII} requires all their temporal momenta to be parallel to it.
\end{remark}

% --- source: two_readings/r2_fourmomentum#5 (The four-momentum of the field)
\subsection{The four-momentum of the field}
\label{r2:sec:field}

In the field reading the field is a function on spacetime $E_{\hat n}$, with coordinates $x^\mu=(ct_{\hat n},x^a)$, valued in a \emph{second} copy $V_{\rm EM}$ of
$\R^{3,3}$ on which $X_{\rm EM}=\sum_i(B_ie_i+E_ie_{i+3})$ is read as a two-form through $\Phi:X\mapsto F_X$, $(F_X)_{0i}=E_i$, $(F_X)_{jk}=\epsilon_{ijk}B_i$, with
$\lG$ as Lorentz algebra (\cref{thm:emfield}). Both copies are needed: by \cref{prop:emexclusive} one six-dimensional space cannot carry
$\lint(\hat n)$ and $\lG$ as its Lorentz action at once.

\begin{proposition}[The dictionary is equivariant]
\label{r2:prop:dict}
$\Phi$ intertwines $\lG$ on $V_{\rm EM}$ with $\lint(\hat n)$ on $\Lambda^2E_{\hat n}^*$ through a unique isomorphism $\phi:\lG\to\lint(\hat n)$, namely
$L_i\mapsto s_j\wedge s_k$ and $K_i\mapsto s_i\wedge \hat n$ for cyclic $(i,j,k)$. Moreover $F_{GX}=\star F_X$ with $\varepsilon_{0123}=+1$; on $\Lambda^2E_{\hat n}^*$ the
Hodge star satisfies $\star^2=-\id$, commutes with $\lint(\hat n)$, and is self-adjoint and anti-isometric for $\langle F,F\rangle=\tfrac12F_{ab}F^{ab}=|\bB|^2-|\bE|^2
=\eta(X,X)$, a form of signature $(3,3)$.
\end{proposition}

Restated from \cref{thm:emfield} and \cref{prop:bivectorF} (r2-12a, which also finds $\phi$; r2-12b). With $\phi$, $E_{\hat n}$ becomes a four-vector module of
$\lG$, and the question of a covariant map from the field to a four-vector becomes a question about $\lint(\hat n)$-modules.

\begin{proposition}[No covariant map from the field to a four-vector]
\label{r2:prop:nolinear}
\begin{enumerate}[label=(\roman*),leftmargin=2.2em]
\item For $\hat n=e_1$ and $\hat n=(2,3,6)/7$ the space $\Hom_{\lint(\hat n)}(\Lambda^2E_{\hat n}^*,E_{\hat n})$ is zero, and so is $\Hom_{\lG}(V_{\rm EM},E_{\hat n})$ for the module
structure given by $\phi$.
\item For quadratic maps, $\Hom_{\lG}(\mathrm{Sym}^2V_{\rm EM},E_{\hat n})=0$, $\dim\Hom_{\lG}(\mathrm{Sym}^2V_{\rm EM},\R)=2$ (the invariants $|\bB|^2-|\bE|^2$ and
$\bE\cdot\bB$), $\dim\Hom_{\lG}(\mathrm{Sym}^2V_{\rm EM},\mathrm{Sym}^2E_{\hat n})=3$, and the traceless part of the last space is one-dimensional.
\end{enumerate}
\end{proposition}

\begin{proof}
Complexified, $V_{\rm EM}\cong(1,0)\oplus(0,1)$ and $E_{\hat n}\cong(\tfrac12,\tfrac12)$, and $\mathrm{Sym}^2V_{\rm EM}\cong(2,0)\oplus(0,2)\oplus(1,1)\oplus2(0,0)$
contains no $(\tfrac12,\tfrac12)$; Schur's lemma \cite{GoodmanWallach} gives the dimensions, and the exact computation of the intertwiner spaces
confirms them.
\end{proof}

Part (i) is inherited representation theory, computed (r2-11a, with two controls: the endomorphisms of $\Lambda^2E^*$ commuting with $\lint$ are spanned by the identity
and the Hodge star, and $\dim\Hom_{\lint}(V,E)=1$). Part (ii) is inherited as well,
computed (r2-11b), and shows that a field momentum must be quadratic in the field and integrated over space: both halves are needed. Quadratic alone is not enough: the only traceless symmetric quadratic covariant is the stress tensor, up to scale, and a four-vector arises only after it is
contracted with the unit normal of a spacelike hyperplane, which is a clock, and integrated.

\begin{proposition}[Pointwise energy--momentum is a tensor row]
\label{r2:prop:stress}
For the stress tensor $\Theta^{\mu\nu}=\eta^{\mu\alpha}F_{\alpha\lambda}F^{\lambda\nu}+\tfrac14\eta^{\mu\nu}F_{\alpha\beta}F^{\alpha\beta}$ of $F_X$
\cite{Jackson}:
\begin{enumerate}[label=(\roman*),leftmargin=2.2em]
\item $\Theta^{00}=E_6$ and $\Theta^{0a}=(\mathbf S_{\rm t})_a=-S_a$; $\Theta$ is symmetric and traceless, $\Theta^\mu{}_\lambda\Theta^\lambda{}_\nu=E_0^2\delta^\mu{}_\nu$ and
$E_6^2=E_0^2+|\mathbf S|^2$.
\item The rest field $X=e_1$ ($E_6=E_0=\tfrac12$, $\mathbf S=0$) boosted by $\exp(\chi\ad_{K_3})$ has $\bB=(\cosh\chi,0,0)$, $\bE=(0,\sinh\chi,0)$,
$E_6=\tfrac12\cosh2\chi$ and $\mathbf S=-\tfrac12\sinh2\chi\,e_3$, while a four-vector with rest components $(\tfrac12,0)$ would have time component
$\tfrac12\cosh\chi$.
\item For every $X$, $\Theta[\exp(\chi\ad_{K_3})X]=L\,\Theta[X]\,L^{\!\top}$ with $L$ the boost of rapidity $\chi$ in the $(0,3)$ plane, the passage to the
frame moving with velocity $-\tanh\chi\,e_3$.
\end{enumerate}
\end{proposition}

Restated from \cref{en:split} and~\cref{en:rainich} and \cref{en:min}, re-verified symbolically (r2-13, r2-14a,b); the identities are inherited
\cite{Rainich1925,LandauLifshitz2}. In the paper's variables the momentum density is $-\mathbf S=-\bE\times\bB$, the textbook $\bE\times\Btext$, that is minus the frame
Poynting vector. So $(E_6,-\mathbf S)$ at a point is the first row of a tensor, quadratic in the boost, and not a four-vector, as \cref{sec:energy} states; the density seen by the clock $\hat n$ is $\Theta(\cdot,\hat n)$, different for each clock and of Lorentz length $E_0$
(\cref{en:rainich}).

\begin{proposition}[The four-momentum of an isolated field]
\label{r2:prop:whole}
\begin{enumerate}[label=(\roman*),leftmargin=2.2em]
\item If $\partial_\mu\Theta^{\mu\nu}=0$ and the Cauchy data have compact support, $P^\mu=(U/c,\mathbf P)$ with $U=\int u\,d^3x$ and $\mathbf P=\int\mathbf g\,d^3x$,
$\mathbf g=\varepsilon_0\,\bE\times\Btext$ in SI, is independent of the spacelike hyperplane and is a future causal four-vector, with
$M^2c^4=U^2-c^2|\mathbf P|^2\ge0$.
\item For two free solutions the interference part $\Theta[X_1+X_2]-\Theta[X_1]-\Theta[X_2]$ is divergence free. Hence $P[X_1+X_2]=P[X_1]+P[X_2]$
whenever the two fields have disjoint supports at some time, and two null wholes with momenta at angle $\vartheta$ have
$M^2c^4=2U_1U_2(1-\cos\vartheta)$, timelike unless $\vartheta=0$, that is unless the two four-momenta are proportional.
\item In Heaviside--Lorentz units with $c=1$ and $f(s)=(1-s^2)^2$ on $|s|\le1$, zero elsewhere, the field $E_x=af(z-t)+bf(z+t)$,
$B_{{\rm t}y}=af(z-t)-bf(z+t)$ has, per unit transverse area, $U=\tfrac{256}{315}(a^2+b^2)$ and $P_z=\tfrac{256}{315}(a^2-b^2)$ at every $t$, so
$M=\tfrac{512}{315}|ab|$. Boosting the fields along $z$ with rapidity $\chi$ gives exactly $(U\cosh\chi-P_z\sinh\chi,\,P_z\cosh\chi-U\sinh\chi)$. For $|t|\ge1$ the
pulses are disjoint, $q\equiv0$ and $\int E_0=0$, while $M>0$.
\end{enumerate}
\end{proposition}

\begin{proof}
(i) Integrate the conservation law over the region between two spacelike hyperplanes; the flux through the timelike boundary vanishes by finite
propagation speed \cite{Laue1911,LandauLifshitz2}. Each $\Theta(\cdot,\hat n)$ is future causal (the dominant energy condition \cite{HawkingEllis73}), hence so
is the integral. (ii) $\Theta$ is quadratic and divergence free on every solution, so its polarisation is divergence free on pairs of solutions; the
interference charge is conserved and vanishes when the supports are disjoint. The mass formula is that of two future null vectors. (iii) The cross terms
of $u$ and of $S_{{\rm t}z}$ cancel pointwise; a right-moving pulse is multiplied by $e^{-\chi}$ and stretched by $e^\chi$ under the boost, and a left-moving one
the reverse.
\end{proof}

Item (i) is inherited (von Laue \cite{Laue1911}) and is the mass of a whole of \cref{en:whole}, restated; items (ii) and (iii) are fixed by
computation (r2-16: the interference part for two Maxwell plane waves with wave vectors $(0,0,3)$ and $(4,0,0)$, the null sum $(3;0,0,3)+(4;4,0,0)$ of
mass squared $24$ and the antiparallel sum $(3;0,0,3)+(4;0,0,-4)$ of mass squared $48$; r2-15, by a genuine Lorentz transformation of the fields and exact integration). With $a=b=1$, (iii) is item~(4) of \cref{en:whole}. The mass of the whole comes from the non-collinearity of the two parts' four-momenta (their three-momenta are antiparallel), not from their
interference, which is absent at $|t|\ge1$. The four-momentum
of the field is the integrated row $\int\Theta^{0\mu}d^3x$, which uses one hyperplane, that is one clock; the result does not depend on it.

\begin{remark}[The photon]
\label{r2:rem:photon}
Inherited, and open in the arena. Quantising the free field gives translation generators with $\widehat P^\mu|\bk,\lambda\rangle=\hbar k^\mu|\bk,\lambda\rangle$,
$k$ null and $\lambda=\pm1$ \cite{Wigner39,WeinbergQFT}; this quantisation is not derived in the arena. The arena supplies the classical counterparts: a null field such as $z=(1,i,0)e^{ik(x_3-t)}$ of \cref{ex:circular} has $E_6=|\mathbf S|$
and momentum density along its propagation (r2-17), and on positive-frequency modes the eigenvalue of $G$ is $i$ times the helicity (\cref{ex:Gcharge}), where the quantum of $\hbar$ per photon is put in by hand. Nor are the arena's own Fock modes photons: $a^\pm_i$ obey the anticommutation
relations on $\R^8$, and no finite matrices obey $[a,a^\dagger]=\id$, since the trace of a commutator vanishes (r2-0b, r2-17).
\end{remark}

\begin{remark}[Base and fibre]
\label{r2:rem:base}
The field's four-momentum is a vector of the base $E_{\hat n}$, the spacetime of the classical reading; the six components of $X_{\rm EM}$ are values in the fibre
$V_{\rm EM}\cong\Lambda^2E_{\hat n}^*$, with units of field strength, and its temporal triple carries $\bB$, not time. So $X_{\rm EM}$ is not a displacement and supplies
no extra time components of momentum, (constructed: the field is a function on the base $E_{\hat n}$, so its translation
charges carry indices of $E_{\hat n}$; covariance alone would allow more, since the invariants $|\bB|^2-|\bE|^2$ and $\bE\cdot\bB$ map covariantly to the
trivial module $U_{\hat n}$, r2-11b). A field that depended on the surplus coordinates would carry surplus
translation charges, the field analogue of $q_a$; these are what the gauge principle of \cref{r2:prop:shellII} removes (\cref{cor:weak}). The two
readings of $X$ thus meet in one place, which is a reading (interpretive): the classical reading supplies the base $E_{\hat n}$, and the field
reading the fibre over it.
\end{remark}

% --- source: two_readings/r2_fourmomentum#6 (The no-go in its scope)
\subsection{The no-go in its scope}
\label{r2:sec:nogo}

\begin{lemma}[Anti-isometry]
\label{r2:lem:anti}
If $G^\dagger=G$ and $G^2=-\id$ then $\eta(Gv,Gw)=-\eta(v,w)$.
\end{lemma}

\begin{proof}
$\eta(Gv,Gw)=\eta(v,G^\dagger Gw)=\eta(v,G^2w)$.
\end{proof}

Restated (\cref{prop:twoK} and \cref{thm:neutral}; r2-0a, r2-18).

\begin{proposition}[No Lorentzian plane is $G$-invariant]
\label{r2:prop:nogo}
Every nondegenerate $G$-invariant subspace of $V$ has neutral signature $(k,k)$. In particular no four-plane of signature $(1,3)$, and no $E_{\hat n}$, is
invariant under an $\eta$-self-adjoint $G$ with $G^2=-\id$.
\end{proposition}

\begin{proof}
If $GW\subset W$, then $G|_W$ is invertible and by \cref{r2:lem:anti} makes $\eta|_W$ congruent to $-\eta|_W$; by Sylvester's law of inertia
\cite{ONeill83} a signature $(p,q)$ then equals $(q,p)$, so it is $(k,k)$. For the binding root the Gram matrix of $\operatorname{span}(v,Gv)$ is
$\bigl(\begin{smallmatrix}a&b\\b&-a\end{smallmatrix}\bigr)$, of determinant $-(a^2+b^2)\le0$.
\end{proof}

Restated (r2-18: the symbolic identities, and 300 random $G$-invariant four-planes of signature $(2,2)$, in floating point
with tolerance $10^{-9}$).

\begin{proposition}[Leakage at every clock]
\label{r2:prop:leak}
For the interior Lagrangian $L_{\rm int}=\tfrac m2\,\eta(\dot X+\kappa GX,\dot X+\kappa GX)$ of equation~\cref{eq:splitL}, with $X$ length-valued and $\kappa=c/\ell$
as in \eqref{eq:kappa}:
\begin{enumerate}[label=(\roman*),leftmargin=2.2em]
\item $p=m\,\eta(\dot X+\kappa GX)$, the Euler--Lagrange equation is $\dot p=\kappa G^{\!\top}p$, and $\pi=\eta^{-1}p$ obeys $\dot\pi=\kappa G\pi$, so
$\pi(\lambda)=(\cos\kappa\lambda+\sin\kappa\lambda\,G)\pi(0)$. At $\hat n=e_1$, $\pi(0)=e_5$ gives $\pi(\lambda)=\cos(\kappa\lambda)e_5-\sin(\kappa\lambda)e_2$ and
$\pr_{e_1}\pi=\cos(\kappa\lambda)e_5$; along every solution $\eta(\pi,\pi)$ becomes $\cos2\kappa\lambda\,\eta(\pi,\pi)+\sin2\kappa\lambda\,\eta(\pi,G\pi)$.
\item For every clock, $\pr^{\perp}_{\hat n}G\hat n=0$ and $\pr^{\perp}_{\hat n}G\pr_{\hat n}=-(\id-nn^{\!\top})\circ\sigma$, where $\sigma:s_a\mapsto e_a$; hence
$\rank\pr^{\perp}_{\hat n}G\pr_{\hat n}=2$, with kernel $\operatorname{span}(\hat n,G\hat n)$ in $E_{\hat n}$.
\item A vector $v\in E_{\hat n}$ stays in $E_{\hat n}$ under $\exp(\theta G)$ for all $\theta$ iff $v\in\operatorname{span}(\hat n,G\hat n)$, a plane of signature $(1,1)$; at $\hat n=e_1$
it is $\langle e_1,e_4\rangle$.
\end{enumerate}
\end{proposition}

\begin{proof}
(i) Direct variation; $\eta^{-1}G^{\!\top}\eta=G$, and $e^{\theta G}=\cos\theta+\sin\theta\,G$ since $G^2=-\id$; $G$ is an anti-isometry, so $e^{\theta G}$ is not an
isometry. (ii) $G$ maps $V^+$ into $V^-\subset E_{\hat n}$, so $\pr^{\perp}_{\hat n}G\hat n=0$; on $V^-$, $G=-\sigma$ lands in $V^+$, of which $\pr^{\perp}_{\hat n}$ keeps the part orthogonal to
$\hat n$: an idempotent of trace two composed with an isomorphism. (iii) If $e^{\theta G}v\in E_{\hat n}$ for all $\theta$, the derivative at $0$ gives $\pr^{\perp}_{\hat n}Gv=0$, so
$v\in\operatorname{span}(\hat n,G\hat n)$; conversely this plane is $G$-invariant, and $\eta(G\hat n,G\hat n)=-1$.
\end{proof}

Fixed (r2-19a with sympy Euler--Lagrange equations; r2-19b symbolic in a general frame, and exact ranks at six rational clocks; r2-19c). At $\hat n=e_1$ the rank is two and only $\langle e_1,e_4\rangle$ stays; (ii) extends the rank to every clock, and the
mechanism is that $G$ maps the three spatial directions into $V^+$, and of $V^+$ only the $\hat n$-component returns to $E_{\hat n}$. If the interior
momentum $\pi$ is identified with $K^\sharp$ and $q\approx0$ is imposed on its flow, the momentum is confined to the $(1,1)$ plane $\operatorname{span}(\hat n,G\hat n)$ and
two of the three spatial momenta are lost: the identification is obstructed. Allowing $q\neq0$ does not help, since neither $\pr_{\hat n}\pi$ nor even
$\eta(\pi,\pi)$ is conserved.

\begin{proposition}[The two Lorentz algebras and the clock]
\label{r2:prop:twoalg}
$\lint(\hat n)$ fixes $U_{\hat n}$ pointwise and $\lG$ fixes no nonzero vector; $\lG\cap\lint(\hat n)=0$; and $\lG$ meets the stabiliser $\so{E_{\hat n}}\oplus\so{U_{\hat n}}$ of the
clock exactly in the line of the diagonal rotation $\sum_in_iL_i$. $G$ commutes with $\lG$ and not with $\lint(\hat n)$.
\end{proposition}

The first two statements restate \cref{prop:notconj}; the rest is fixed (r2-20 at $\hat n=e_1,(1,2,2)/3$, and $\widehat W$ for the fixed planes).
\Cref{r2:prop:leak} is the vector-module face of the exclusion of \cref{prop:emexclusive}: the operator that is the Hodge star for $\lG$ moves every
four-plane of $\lint(\hat n)$, and the only symmetry the two structures share is the rotation about the clock, whose axis plane $\operatorname{span}(\hat n,G\hat n)$ is the one
that survives in \cref{r2:prop:leak}(iii).

The no-go does not obstruct $G$ as the Hodge star. On the field module $\Lambda^2E_{\hat n}^*$ the star acts on the whole module, which is neutral of signature $(3,3)$
(\cref{r2:prop:dict}); \cref{r2:prop:nogo} forbids only invariant subspaces that are not neutral, and \cref{thm:neutral} requires exactly $P=Q$, which
$\Lambda^2E_{\hat n}^*$ satisfies. So $G_{\rm EM}=\star$ is electric--magnetic duality on field strengths, commuting with the Lorentz action, while $G_{\rm cl}$ on the vector
module $E_{\hat n}$ is obstructed. This is restated (\cref{thm:emfield} and~\cref{thm:neutral}; r2-12b with \cref{r2:prop:nogo}) .

% --- source: two_readings/r2_fourmomentum#7 (Answer)
\subsection{Answer}
\label{r2:sec:answer}

For the particle reading the four-momentum is constructed as follows. $X$ is taken with its lengths on the three directions of $V^-$; a clock $\hat n\in V^+$
selects the four-plane $E_{\hat n}=\R \hat n\oplus V^-$, and every four-plane containing the lengths is of this form. The six-momentum conjugate to $Y$ splits into the
four-momentum $P_{\hat n}=(E/c,\bp)$, conjugate to $(ct_{\hat n},x^a)$, and two surplus momenta $q_a$, Lorentz scalars. $P_{\hat n}$ is a four-vector and the Noether charge of
the translations of $E_{\hat n}$. What it needs beyond the arena is a dynamics: the arena's metric gives the ambient shell, on which $q$ lowers the mass;
the gauge principle of \cref{sec:constraints} gives the reduced shell $q\approx0$, on which the temporal momentum is parallel to the clock and $m_4=\mu$; a positive internal mass
needs a functional of the opposite sign, that of extra space dimensions. For the field reading the four-momentum is the integrated first row of the stress
tensor on the base $E_{\hat n}$, a four-vector for an isolated whole, with the mass of a whole from parts whose four-momenta are not collinear; no linear or quadratic covariant map
produces it pointwise, and $X_{\rm EM}$ supplies no time components. The two readings share the clock: it fixes $E_{\hat n}$, the complex structure $J(\hat n)$ of the
Dirac spinor and the hyperplane on which the field's momentum is integrated.

{\small
\renewcommand{\arraystretch}{1.12}
\begin{longtable}{@{}>{\raggedright\arraybackslash}p{9.0cm}>{\raggedright\arraybackslash}p{2.8cm}>{\raggedright\arraybackslash}p{2.6cm}@{}}
\toprule
statement & class & checks\\
\midrule
\endhead
four-planes containing $V^-$ are the $E_{\hat n}$ (\cref{r2:prop:planes}) & fixed & r2-1a,b\\
$\pr_{\hat n}$ covariant under $SO^+(1,3)\times SO(2)$ only (\cref{r2:prop:cov}) & restated & r2-2a--c\\
$J(\hat n)=I_T\hat n$, $J(\hat n)=u_1u_2$, $J(-\hat n)=-J(\hat n)$ (\cref{r2:prop:J}) & restated & r2-3a--d\\
\cref{thm:i} at every clock; mass form $I_T$ clock-independent (\cref{r2:thm:i}) & fixed & r2-4a--d\\
$P_{\hat n}$ four-vector, Noether charge of translations (\cref{r2:prop:cons}) & fixed & r2-5a--c\\
ambient shell, tachyonic for $|\mathbf q|>M_6c$ (\cref{r2:prop:shellI}) & fixed & r2-6\\
Option I is the momentum reading of the undivided energy (\cref{r2:prop:RE}) & restated & r2-7\\
first-class reduction $=$ \cref{thm:wlgauge} at $A=0$ (\cref{r2:prop:shellII}) & restated & r2-8a,b\\
internal mass functional; toy has signature $(1,5)$ (\cref{r2:prop:shellIII}) & constructed; sign fixed & r2-9a,b\\
the arena prefers Option II (\cref{r2:rem:prefer}) & interpretive & --\\
$m_4(\hat n)^2\le\mu^2$, equality iff $\hat n\parallel\boldsymbol{\mathsf E}$; octant maximum (\cref{r2:prop:rotate}) & fixed & r2-10a--d\\
mass as projection onto the clock (\cref{r2:rem:projection}) & interpretive & --\\
dictionary equivariant, $\star$ on a neutral module (\cref{r2:prop:dict}) & restated & r2-12a,b\\
no linear map to a four-vector (\cref{r2:prop:nolinear}(i)) & inherited & r2-11a,b\\
no quadratic map to a four-vector (\cref{r2:prop:nolinear}(ii)) & inherited & r2-11b\\
pointwise $(E_6,-\mathbf S)$ is a tensor row (\cref{r2:prop:stress}) & restated & r2-13, r2-14a,b\\
von Laue four-vector, mass of a whole (\cref{r2:prop:whole}(i)) & inherited & --\\
interference conserved; exact packets (\cref{r2:prop:whole}(ii),(iii)) & fixed & r2-15, r2-16\\
photon $P=\hbar k$, two helicities (\cref{r2:rem:photon}) & inherited; open in the arena & r2-17\\
$X_{\rm EM}$ supplies no time components (\cref{r2:rem:base}) & constructed & r2-11b\\
no $G$-invariant $(1,3)$ plane (\cref{r2:prop:nogo}) & restated & r2-18\\
$\rank\pr^{\perp}_{\hat n}G\pr_{\hat n}=2$ at every clock; $\pi\equiv K^\sharp$ with $q\approx0$ (\cref{r2:prop:leak}) & fixed; obstructed & r2-19a--c\\
$\lG\cap(\so{E_{\hat n}}\oplus\so{U_{\hat n}})=\R\sum_in_iL_i$ (\cref{r2:prop:twoalg}) & fixed & r2-20\\
$G=\star$ on $\Lambda^2E_{\hat n}^*$ not obstructed & restated & r2-12b\\
a principle making the normal translations gauge & open & --\\
\bottomrule
\end{longtable}}

What remains open, located precisely, is a principle that makes the surplus translations gauge rather than
chosen, a dynamics that selects the clock rather than aligning it with the momentum, and a derivation of the photon's Hilbert space from the arena. Nothing
in this section identifies $X_{\rm cl}$ with $X_{\rm EM}$; they share the clock, and through it the base on which both four-momenta live.

% --- source: two_readings/r3_octants#0 ((head))
% Track r3 -- the field-strength sector X_EM = (B;E) on the Clifford octants.

% --- source: two_readings/r3_octants#1 (The field-strength sector X\_EM = (B;E) on the Clifford octa)
\section{\texorpdfstring{The field-strength sector $X_{\rm EM}=(\bB;\bE)$ on the Clifford octants}{The field-strength sector X\_EM = (B;E) on the Clifford octants}}
\label{r3:sec:main}

The question asked in the introduction about the octants concerns the field reading $X_{\rm EM}=(\bB;\bE)$, where each of the six components
$E_x,E_y,E_z$, $B_x,B_y,B_z$ is to sit on an octant. As a $\lint$-module ($V_{\rm EM}\cong\Lambda^2E_{\hat n}^*$, \cref{r1:sec:types}), $X_{\rm EM}$ is not a sum
of two vectors and not a displacement: the pair $(\bB;\bE)$ labels the six components of the field two-form $F$, and the question is how
these six components, as six independent components of one field strength, sit on the eight octants of the spinor module. Everything below
is a classical field representation: placing $(\bB;\bE)$ on six lines of $S$ quantises nothing. A construction of this kind was proposed in an earlier, superseded version of the programme: the eight octants are the joint
eigenspaces of the pole operators $\omega_i$; the six mixed octants are field octants, those with one plus sign (temporal weight one
there, Fock weight $N=2$ here) carrying the electric components (as boosts) and those with two plus signs ($N=1$) the magnetic
components (as rotations); the two poles carry ``mass'' and
``statistics''; and Maxwell's equations are ``pooled'' across the octants. This section establishes in the conventions of \cref{sec:conventions} what of this is true. In brief: the octants are the eight lines of a Fock basis of the spinor module $S=\R^8$, which is the exterior
algebra of a null three-plane; there is a field-octant map $\Phi$ placing the electric member $\bE$ of the pair on the Fock weight $N=1$ and the magnetic member $\bB$ on $N=2$, unique up to
a scale and a quarter turn of duality once rotations and Maxwell's equations are required, the quarter turn giving the earlier
labelling; it carries the rotations, a parity and the
premetric $\mathfrak{gl}(3)$, but no even element of $\Cl(3,3)$ carries the boosts or the duality rotation; the Spin$(3,3)$-invariant mass
form pairs each octant with its antipode and $\bE$ with $\bB$; and Maxwell's four equations are one first-order equation on $S$, with the
two Gauss laws on the poles. In textbook variables the whole construction is Hestenes' Pauli-algebra electrodynamics, written with
the three temporal generators acting on the vacuum pole.

\paragraph{Conventions.} $V=\R^{3,3}$, $\eta=\operatorname{diag}(1,1,1,\allowbreak-1,-1,-1)$, the generators $e_1,\dots,e_6$ are those of eq.~\cref{eq:realisation}, $\omega:=e_1\cdots e_6=\Gamma$, $T:=e_1e_2e_3$ (the mass form, \cref{cor:fullmass}), and $\varpi=\sum_i\omega_i$,
$\omega_i:=e_ie_{i+3}$. The field point is $X=\sum_i(B_ie_i+E_ie_{i+3})$ (\cref{thm:emfield}), with $\bB$ minus the textbook magnetic
field, $\Btext=-\bB$; $z=\bB+i\bE$, $q=z\cdot z$, $E_6=\tfrac12|z|^2$, $E_0=\tfrac12|q|$ (eq.~\cref{en:defs}). The binding root is
$Ge_i=e_{i+3}$, $Ge_{i+3}=-e_i$; the grading root is $\cK=\tfrac12\ad_{\varpi}$, $\cK e_i=-e_{i+3}$, $\cK e_{i+3}=-e_i$
(\cref{prop:twoK} and \cref{lem:Gclifford}). On $V$ we use coordinates
$(\bB;\bE)\in\R^6$. Units are Heaviside--Lorentz with $c=1$ unless SI is stated. Dimensions: in SI $[\bE]=\mathrm{V\,m^{-1}}$,
$[\bB]=\mathrm T$, so that $\bE$ and $c\bB$ share a unit; $[t]=\mathrm s$, $[x^a]=\mathrm m$, $[\partial_a]=\mathrm m^{-1}$; $[\rho]=\mathrm{C\,m^{-3}}$,
$[\bj]=\mathrm{A\,m^{-2}}$, $[\varepsilon_0]=\mathrm{F\,m^{-1}}$; the generators, the Fock operators, the octant basis vectors and $T$ are
dimensionless, and $\psi=\Phi_{\rm oct}(X)$ carries the dimension of a field strength. Identifiers r3-$n$ refer to the script
\texttt{r3\_octants.py}, which checks every finite-dimensional and symbolic claim below exactly.

% --- source: two_readings/r3_octants#4 (The field-octant map)
\subsection{The field-octant map}
\label{r3:sec:phi}

The map should intertwine the rotations of the field. Which rotations act on both triples as the field's rotations is fixed first.

\begin{lemma}[The diagonal rotations]
\label{r3:lem:rot}
For cyclic $(i,j,k)$ the bivector $L_i=\tfrac12(e_{j+3}e_{k+3}-e_je_k)$ of \cref{sec:energy} acts on $V$ by $\delta\bB=e_i\times\bB$, $\delta\bE=e_i\times\bE$
and commutes with $G$ and $\cK$. The paired bivector $e_je_k+e_{j+3}e_{k+3}$ rotates the two triples in opposite senses and commutes with
neither. The diagonal rotations $\so3_\Delta=\operatorname{span}\{L_i\}$ are the compact part of $\lG$ and the antisymmetric part of
$\mathfrak{gl}(3)$, and $\so3_\Delta=\lG\cap\mathfrak{gl}(3)$. In Fock form
\[
L_i=a^+_ka^-_j-a^+_ja^-_k,\qquad K_i=\tfrac12(e_ke_{j+3}-e_je_{k+3})=a^-_ja^-_k-a^+_ja^+_k .
\]
The spin lift of $L_i$ (left multiplication) satisfies $L_i^3=-L_i$, so $\exp(2\pi L_i)=\id$ on $S$; it annihilates exactly
$\operatorname{span}\{|\mathrm{vac}\rangle,|\mathrm{full}\rangle\}$ and acts as $L_if_m=\epsilon_{imn}f_n$, $L_ig_m=\epsilon_{imn}\mathsf g_n$ on
\[
f_k:=a^+_k|\mathrm{vac}\rangle,\qquad \mathsf g_m:=\tfrac12\epsilon_{mjk}\,a^+_ja^+_k|\mathrm{vac}\rangle .
\]
\end{lemma}

\begin{proof}
$[e_je_k,e_k]=2e_j$ and $[e_je_k,e_j]=-2e_k$, while $[e_{j+3}e_{k+3},e_{k+3}]=-2e_{j+3}$ because $e_{k+3}^2=-\id$: the relative minus sign makes the
triples turn alike. $G$ and $\cK$ map $e_i$ to $\pm e_{i+3}$, so they commute with $R\oplus R'$ iff $R=R'$. The Fock forms follow from
$e=a^++a^-$, $e'=a^+-a^-$ and the CAR, and the action on $f,g$ from $[L_i,a^+_m]=\epsilon_{imn}a^+_n$ and $L_i|\mathrm{vac}\rangle=0$. The intersection is
computed in r3-9, which also checks $K_i=GL_i$ on $V$ (\cref{prop:spinmatrices}) and the brackets of \cref{thm:emfield}.
\end{proof}

Restated and fixed (r3-9). The paired combination $e_je_k+e_{j+3}e_{k+3}$, a natural first guess, is the wrong one; with the minus sign it is the diagonal
rotation of \cref{thm:GLorentz}, whose spin lift has integer spin on $S$, as \cref{prop:halfspinors} states.

\begin{theorem}[The field-octant map]
\label{r3:thm:phi}
\begin{enumerate}[label=(\roman*),leftmargin=2.2em]
\item The linear maps $V\to S$ intertwining $\so3_\Delta$ form a four-dimensional space. Those sending $\bE$ into weight one and $\bB$ into
weight two form the two-parameter family
\[
\Phi_{{\rm oct},\alpha,\beta}(X)=\alpha\sum_kE_kf_k+\beta\sum_mB_mg_m ,
\]
and $\Phi:=\Phi_{1,1}$ is an isometry of $\R^6$ onto the six field octants.
\item Under $a^+_I|\mathrm{vac}\rangle\leftrightarrow dx^I$, $\Phi_{\rm oct}(X)=E_k\,dx^k+\tfrac12\epsilon_{mjk}B_m\,dx^j\wedge dx^k$: $\bE$ is a one-form and $\bB$ a two-form
on $\R^3$.
\item The spin lift of $\mathfrak{gl}(3)$ preserves the image of $\Phi$ and induces $\bE\mapsto(A-\tfrac12\tr A)\bE$ and $\bB\mapsto(\tfrac12\tr A-A^{\!\top})\bB$: $\bE$
transforms in $L\otimes(\det L)^{-1/2}$ and $\bB$ in $L^*\otimes(\det L)^{1/2}$, contragredient up to the determinant.
\item The same $\mathfrak{gl}(3)\subset\so{3,3}$ acts on the field point through the light-cone components, $\bB+\bE\mapsto A(\bB+\bE)$,
$\bB-\bE\mapsto-A^{\!\top}(\bB-\bE)$. For symmetric $A$ this is $\bB\mapsto A\bE$, $\bE\mapsto A\bB$, a set of boosts, and $\Phi$ intertwines the
vector action of $\mathfrak{gl}(3)$ with its spin lift only on $\so3_\Delta$.
\end{enumerate}
\end{theorem}

\begin{proof}
(i) As $\so3_\Delta$-modules $V=2\cdot\R^3$ and, by \cref{r3:lem:rot}, $S=2\cdot\R^3\oplus2\cdot\R$ (the $f$'s, the $g$'s and the poles). The
vector representation is absolutely irreducible, so by Schur's lemma \cite{GoodmanWallach} $\Hom(\R^3,\R^3)=\R$ and $\Hom(\R^3,\R)=0$; the
weight condition leaves one scale on each copy. (ii) $a^+_k$ acts as $dx^k\wedge$, and $g_1=a^+_2a^+_3|\mathrm{vac}\rangle\leftrightarrow dx^2\wedge dx^3$
cyclically. (iii) follows from \cref{r3:prop:gl3}: on $\Lambda^1$ the action is $A-\tfrac12\tr A$, and on $\Lambda^2(L)\cong L^*\otimes\det L$, in the
basis $\mathsf g_m$, it is $\tr A-A^{\!\top}$, from which $\tfrac12\tr A$ is subtracted. (iv) is \cref{r3:prop:gl3} on $V$ with \cref{r3:eq:Xfock}
(r3-10, r3-11).
\end{proof}

Items (i), (iii), (iv) are fixed; $\Phi$ itself is constructed. Item (ii) is the premetric statement that the electric field strength is a
one-form and the magnetic field strength a two-form on space, $F=\bE\wedge d\sigma+\bB$ in the notation of Hehl and Obukhov \cite{HehlObukhov03}
(inherited). Item (iv) makes precise the statement that under $\mathfrak{gl}(3,\R)$ $\bE$ is a vector and $\bB$ a two-form: it is
true of the spin lift on the octants and false of the vector action on the field point, where the symmetric part of $\mathfrak{gl}(3)$
mixes $\bE$ with $\bB$. The two scales of (i) are tied together by Maxwell's equations in \cref{r3:prop:maxunique}.

\begin{proposition}[Parity on the octants]
\label{r3:prop:parity}
\begin{enumerate}[label=(\roman*),leftmargin=2.2em]
\item $\Ad_T|_V=\eta$ is the frame parity, fixing $\bB$ and reversing $\bE$ (\cref{en:parity}), and $T$ is, up to a factor, the only
element of $\Cl(3,3)$ implementing it. $T$ carries each octant to its antipode,
\[
T|\mathrm{vac}\rangle=|\mathrm{full}\rangle,\quad Tf_k=\mathsf g_k,\quad Tg_k=-f_k,\quad T|\mathrm{full}\rangle=-|\mathrm{vac}\rangle,
\]
and $T\Phi=-\Phi\circ G$: on the field octants the spinor lift of parity acts as minus the Hodge star, exchanging electric and magnetic octants.
\item $(-1)^N\Phi=\Phi\circ\eta$: the weight parity, which is the chirality $\gamma_5=-\Gamma$, acts as $-1$ on the image of $\bE$ and $+1$ on the image
of $\bB$; but $\Ad_{(-1)^N}=-\id$ on $V$.
\item The Dirac parity $P=\gamma^0=e_1$ of \cref{thm:PT} acts on the frame field by
$(\bB;\bE)\mapsto(B_1,-B_2,-B_3;-\bE)$, and $e_1f_1=|\mathrm{vac}\rangle$: it leaves the field octants.
\end{enumerate}
\end{proposition}

\begin{proof}
(i) $T$ commutes with $e_1,e_2,e_3$ and anticommutes with $e_4,e_5,e_6$; uniqueness is Skolem--Noether. $T=\prod_i(a^+_i+a^-_i)$ flips all three
signs, and the displayed values are r3-12. With them $T\Phi_{\rm oct}(X)=\sum_k(E_kg_k-B_kf_k)$, while $GX$ has $\bB'=-\bE$, $\bE'=\bB$. (ii) $(-1)^N$ is $-1$ on
$f$ and $+1$ on $g$; $\Gamma$ anticommutes with every vector. (iii) $e_1$ commutes with itself and anticommutes with the other generators,
and it flips the first sign (r3-12).
\end{proof}

Fixed. The polar and axial characters of $\bE$ and $\bB$ are present on the octants, but they are carried by the weight parity
$(-1)^N=-\Gamma$, whose Clifford action on $V$ is the total reflection $X\mapsto-X$, and not by the spinor lift $T$ of the parity, which acts on
the field octants as the duality rotation by $-\pi/2$. This is possible because $\Phi$ intertwines the rotations only, not the Clifford
action of $\mathrm{Pin}(3,3)$. The statement that parity acts as $-1$ on the electric and $+1$ on the magnetic octants is therefore false if
parity means the spinor lift of the spatial inversion, and true if it means the weight parity.

\begin{remark}[The earlier field-octant bijection]
\label{r3:rem:bijection}
An alternative labelling calls the octants with one plus sign electric and those with two plus signs magnetic, the octant with its single
plus sign in pair $k$ carrying $E_k$ and the octant with its single minus sign in pair $m$ carrying $B_m$. Its electric octants are
the lines of the $\mathsf g_k$, of weight $N=2$, and its magnetic ones those of the $f_m$, of weight $N=1$: up to signs, which it does not fix, its
labelling is that of $\Phi\circ G$ (\cref{r3:prop:maxunique}), not of $\Phi$. In both labellings $E_k$ and $B_k$ sit on antipodal octants,
which the mass form pairs (\cref{r3:prop:energies}). The earlier bijection mapped octants to a basis of interior bivectors, a bijection of labels
and not a linear map; it is admissible as a labelling and equivalent to $\Phi$ by a duality rotation of a quarter turn, and nothing in the
octant structure prefers it. Its derivation of $E=cB$ from a count of ``temporal kernels'' is not reproduced: in the present construction the
ratio of the two scales of $\Phi_{{\rm oct},\alpha,\beta}$ is $c$ in SI units, the conversion constant of \cref{en:onec}, fixed by Faraday's
law (\cref{r3:prop:maxunique} and \cref{r3:thm:maxwell}(iii)).
\end{remark}

% --- source: two_readings/r3_octants#5 (Clifford multiplication does not give the pair reading)
\subsection{Clifford multiplication does not give the pair reading}
\label{r3:sec:poles}

The field point is also a Clifford vector, and one may ask what it does to the octants by left multiplication. The answer shows why the
pair reading $X_{\rm EM}=(\bB;\bE)$ needs the map $\Phi$ and is not obtained by multiplying.

\begin{proposition}[Light-cone amplitudes]
\label{r3:prop:poles}
For every $X=\sum_i(B_ie_i+E_ie_{i+3})$
\begin{equation}
\label{r3:eq:poles}
X|\mathrm{vac}\rangle=\sum_i(B_i+E_i)\,f_i,\qquad X|\mathrm{full}\rangle=\sum_i(B_i-E_i)\,g_i=\sum_i(-1)^{i-1}(B_i-E_i)\,|111\oplus\delta_i\rangle ,
\end{equation}
and $X^2=\eta(X,X)\id=\sum_i(B_i+E_i)(B_i-E_i)\,\id$. The amplitudes $B_i\pm E_i$ are the coordinates of $X$ on $L$ and $L'$, the eigen-coordinates of
$\cK$, with $\bB+\bE=\sqrt2\,\Re(e^{-i\pi/4}z)$ and $\bB-\bE=-\sqrt2\,\Im(e^{-i\pi/4}z)$. Further
\begin{equation}
\label{r3:eq:eighth}
X\bigl(|\mathrm{vac}\rangle+|\mathrm{full}\rangle\bigr)=\sqrt2\,\Phi\bigl(e^{\pi G/4}X\bigr).
\end{equation}
No spinor $\chi$ satisfies $X\chi=\Phi_{\rm oct}(X)$ for all $X$; the Clifford multiplications $X\mapsto X\chi$ that intertwine the rotations are those with
$\chi\in\operatorname{span}\{|\mathrm{vac}\rangle,|\mathrm{full}\rangle\}$, and none of them places $\bE$ and $\bB$ on separate weights.
\end{proposition}

\begin{proof}
\Cref{r3:eq:poles} follows from \cref{r3:eq:Xfock} with $a^-_i|\mathrm{vac}\rangle=0$, $a^+_i|\mathrm{full}\rangle=0$ and $a^-_i|\mathrm{full}\rangle=g_i$; $X^2$ follows from the
CAR. The duality rotation $e^{\pi G/4}$ gives $\bE'=(\bE+\bB)/\sqrt2$, $\bB'=(\bB-\bE)/\sqrt2$, which is \cref{r3:eq:eighth}. If $X\chi=\Phi_{\rm oct}(X)$, then
$e_1\chi=g_1$ and $e_4\chi=f_1$, i.e.\ $\chi=e_1|011\rangle=|111\rangle$ and $\chi=-e_4|100\rangle=|000\rangle$, a contradiction. $L_i(X\chi)=[L_i,X]\chi+XL_i\chi$, so
equivariance requires $L_i\chi=0$, and then $X(a|\mathrm{vac}\rangle+b|\mathrm{full}\rangle)$ puts $a(\bB+\bE)$ on weight one (r3-13).
\end{proof}

Fixed; the reading is interpretive. Read as a Clifford vector acting on the vacuum pole, the field point does not keep the two members of
the pair apart: it places on the $i$-th electric octant the sum $E_i+B_i$, which in textbook variables is $E_i-B_{{\rm t}i}$, and on the
filled pole it places $B_i-E_i=-(E_i+B_{{\rm t}i})$ on the $i$-th magnetic octant. These sums are the coordinates of $X$ on the null planes
$L$ and $L'$, the light-cone components of the para-complex structure; they mix the electric and the magnetic member and are not the
field strength. A literal reading ``$X=\bE+\bB$'' would therefore be wrong: the pair $(\bB;\bE)$ is placed on the octants by $\Phi$, which is
not Clifford multiplication by any spinor, and the Clifford action differs from $\Phi$ by a duality rotation of an eighth of a turn,
\cref{r3:eq:eighth}. (An eighth of a turn also separates the Klein correspondence from the frame reading in \cref{thm:klein}; the
maps are different and we draw no conclusion from the coincidence.)

% --- source: two_readings/r3_octants#6 (Boosts and duality cannot act by even elements)
\subsection{Boosts and duality cannot act by even elements}
\label{r3:sec:boosts}

\begin{theorem}[The even obstruction]
\label{r3:thm:noeven}
Every even element of $\Cl(3,3)$ commutes with $(-1)^N$ and maps an octant of weight $w$ into the weights $w$ and $w\pm2$; this includes $\so{3,3}$,
$\mathrm{Spin}(3,3)$ and the spin lifts of $\lG$ and $\lint$. No even $Q$ satisfies $Q\Phi=\Phi K_k$ for a boost $K_k$ of $\lG$, or $Q\Phi=\Phi G$.
At the level of groups, no element of $\Cl(3,3)$ of definite parity intertwines $e^{sK_k}$ $(s\neq0)$ with $\Phi$, nor $e^{\vartheta G}$ unless
$\vartheta\in\tfrac\pi2\Z$.
\end{theorem}

\begin{proof}
$(-1)^N=-\omega$ commutes with even and anticommutes with odd elements, and weights of equal parity differ by $0$ or $2$. By \cref{ex:emboost}, $K_3e_1=e_5$; so $\Phi_{\rm oct}(e_1)=g_1$ has weight two and $\Phi_{\rm oct}(K_3e_1)=f_2$ weight one, and an even $Q$ cannot map the first to the
second. Likewise $\Phi_{\rm oct}(e_1)=g_1$ and $\Phi_{\rm oct}(Ge_1)=\Phi_{\rm oct}(e_4)=f_1$. For the groups, $e^{sK_3}e_1=\cosh s\,e_1+\sinh s\,e_5$, so $\Phi_{\rm oct}(e^{sK_3}e_1)$ has
components of both parities when $s\neq0$, while an element of definite parity maps $g_1$ to a vector of definite parity; the same holds for
$e^{\vartheta G}e_1=\cos\vartheta\,e_1+\sin\vartheta\,e_4$ unless $\cos\vartheta\sin\vartheta=0$ (r3-14, r3-15).
\end{proof}

Obstructed, under the hypothesis that the operator has definite parity; the hypothesis covers every element of the spin group and of its
Lie algebra. The scope is the realisation chosen here: the theorem says that the spin group of $\Cl(3,3)$, acting on $S$ by left multiplication,
does not implement the Lorentz action of $\lG$ on the field values placed in the octants by the map $\Phi$. It does not say that electromagnetic
fields fail to be Lorentz covariant: they transform covariantly as two-forms, through $\lG$ acting on $V$, with no reference to spinors. Nor does a
shared operator identify the two modules: that the same temporal volume element enters both the spin-$\tfrac12$ module $S$ and the octant
representation of the spin-one field strengths does not make these modules equivalent, and they are not (\cref{r3:prop:odd} and r4-13; X-06).
The intertwiners that do exist are odd:

\begin{proposition}[Odd intertwiners and the octant representation]
\label{r3:prop:odd}
\begin{enumerate}[label=(\roman*),leftmargin=2.2em]
\item $-T\Phi=\Phi G$ and $Q_k\Phi=\Phi K_k$ for
\[
Q_k:=-TL_k=-\tfrac12e_k(\id-\omega_j\omega_l)=-\tfrac12(e_k+e_{k+3}\omega),\qquad\{k,j,l\}=\{1,2,3\},
\]
an odd element with parts of grades one and five. The solutions of $Q\Phi=\Phi K_k$ form an affine space of dimension $16$, of dimension $8$
among odd elements, and contain no Clifford vector; no Clifford vector implements $G$ on the octants.
\item On $S=\Cl(V^+)|\mathrm{vac}\rangle$ (\cref{r3:prop:fock}), $Q_k(a|\mathrm{vac}\rangle)=-\tfrac12[e_k,a]|\mathrm{vac}\rangle$ and $-T$ is left multiplication by the
central element $-e_1e_2e_3$ of $\Cl(V^+)$.
\item $\rho_{\rm oct}(L_k)=L_k$, $\rho_{\rm oct}(K_k)=Q_k$ is a representation of $\lG$ on $S$. Its Casimir $\sum_k\rho_{\rm oct}(K_k)^2-\rho_{\rm oct}(L_k)^2$
is $4$ on the field octants, the field-strength value, and $0$ on the poles, and its intertwiners $V\to S$ are $\operatorname{span}\{\Phi,\Phi G\}$.
\item The spin lift of $\lG$ has Casimir $3$ on all of $S$ and admits no nonzero intertwiner $V\to S$. Its boosts carry $f_k$ to $-|\mathrm{full}\rangle$
and $|\mathrm{vac}\rangle$ to $-\mathsf g_k$: under the genuine Lorentz action the electric octants are the spatial components of the four-vector $S^+$
whose time component is the filled pole, and the magnetic octants those of the four-vector $S^-$ whose time component is the vacuum.
\end{enumerate}
\end{proposition}

\begin{proof}
(i) $T$ commutes with the same-sign bivectors, hence with $L_k$, so $-TL_k\Phi=-T\Phi L_k=\Phi GL_k=\Phi K_k$ by \cref{r3:prop:parity}(i) and $K_k=GL_k$.
The closed forms use $T(e_je_l)=-e_k$ for cyclic $(k,j,l)$ and $e_{k+3}\omega=-e_k\omega_j\omega_l$. The values of a solution on the two poles are free, eight
components each, four of them of the parity opposite to the pole (r3-15). (ii) $e_{k+3}\omega$ commutes with $e_1,e_2,e_3$ and sends $|\mathrm{vac}\rangle$ to $-e_k|\mathrm{vac}\rangle$,
so it acts on $a|\mathrm{vac}\rangle$ as $-a\,e_k$. (iii) $T$ commutes with $\rho_{\rm oct}(L)$ and $T^2=-\id$, so the brackets of $\lG$ follow from
those of $\so3$; $\rho_{\rm oct}(K)^2=-L^2$ gives the Casimir, and an intertwiner is a $T$-linear map $\C^3\to\C^3$ commuting with $\so3$, hence a
complex multiple of $\Phi$ (r3-16). (iv) The Casimir is \cref{prop:halfspinors}; $S=2(\tfrac12,\tfrac12)$ contains no
$(1,0)\oplus(0,1)$; the boost formulas are $K_k=a^-_ja^-_l-a^+_ja^+_l$ on $f_k$ and $|\mathrm{vac}\rangle$ (r3-16).
\end{proof}

Constructed: the operators $Q_k$ and the representation $\rho_{\rm oct}$; fixed: the solution counts of (i) and the Casimir values and the
intertwiner space of (iii); restated and fixed: (iv). Item (ii) identifies the construction: in textbook variables (\cref{r3:thm:maxwell}) the
field octants carry Hestenes' Pauli-algebra field $F=\bE+I\Btext$ \cite{Hestenes66}, written with $\sigma_k=e_k$ and $I=T$, and the boosts and the
duality rotation act as they do in the Pauli algebra, through the commutator with $\sigma_k$ and through multiplication by $I$. The
commutator is a left multiplication combined with a right multiplication, and on $S$ the right multiplication is the odd element
$e_{k+3}\omega$: this is the reason the boosts are odd. Which symmetries of the field the octant picture carries is now settled.

\begin{corollary}[What the octants carry]
\label{r3:cor:carries}
By even elements of $\Cl(3,3)$ the field octants carry the rotations $\so3_\Delta$ and the premetric $\mathfrak{gl}(3)$ of \cref{r3:thm:phi}(iii),
through their spin lifts, and the parity, through $(-1)^N=-\omega$. They do not carry the boosts or the duality
rotations, which act only through the odd operators $Q_k$ and $-T$ of the representation $\rho_{\rm oct}$, the Pauli algebra of the temporal triple.
The spin-lifted Lorentz algebra $\lG$ makes the field octants parts of two four-vectors, not of a field strength.
\end{corollary}

There is no nonzero linear Lorentz-equivariant map from the field-strength module to the four-vectors. Item (iv) gives the same fact for the spinor module, whose modules under $\lG$ are four-vectors (for $E_n$ itself
the statement is \cref{r2:prop:nolinear}); on the octants the Hodge star of the field is represented, but by the odd element $-T$, with $T=I_T=e_1e_2e_3$, not by an element of $\mathrm{Spin}(3,3)$.

% --- source: two_readings/r3_octants#7 (Energies on the octants)
\subsection{Energies on the octants}
\label{r3:sec:energies}

\begin{proposition}[Energies and the Chevalley--Mukai pairing]
\label{r3:prop:energies}
Let $\psi=\Phi_{\rm oct}(X)=\psi_{\bE}+\psi_{\bB}$ with $\psi_{\bE}=\sum_kE_kf_k$, $\psi_{\bB}=\sum_mB_mg_m$.
\begin{enumerate}[label=(\roman*),leftmargin=2.2em]
\item $|\psi|^2=|\bE|^2+|\bB|^2=2E_6$ and $\psi^{\!\top}(-1)^N\psi=|\bB|^2-|\bE|^2=\Re q$.
\item $\psi_{\bE}^{\!\top}T\psi_{\bB}=-\bE\cdot\bB=-\tfrac12\Im q$, while $\psi^{\!\top}T\psi=0$ and $\psi^{\!\top}T\Gamma\psi=2\,\bE\cdot\bB=\Im q$. Hence
\[
E_0=\tfrac12|q|=\tfrac12\Bigl[\bigl(\psi^{\!\top}(-1)^N\psi\bigr)^2+4\bigl(\psi_{\bE}^{\!\top}T\psi_{\bB}\bigr)^2\Bigr]^{1/2}.
\]
\item The $\mathrm{Spin}(3,3)$-invariant bilinear forms on $S$ are $\operatorname{span}\{T,T\Gamma\}$, $T$ antisymmetric and $T\Gamma$ symmetric. Under
$\Lambda^*(L)\cong S$ the Chevalley--Mukai pairing $(\varphi,\chi)=[\tau(\varphi)\wedge\chi]_{\rm top}$, with $\tau$ the reversal, equals
$-\varphi^{\!\top}T\chi$. Two octants pair to a nonzero value iff their null planes $N_\sigma,N_\tau$ are transverse, iff they are antipodal.
\item $(-1)^N$ is invariant under the six compact bivectors and under $\rho_{\rm oct}$, and under none of the nine boosts. The Euclidean norm is
invariant under the compact bivectors and the $\rho_{\rm oct}(L_k)$, but not under the boosts $\rho_{\rm oct}(K_k)=Q_k$, which are symmetric:
$|\psi|^2=2E_6$ changes under boosts, $|\Phi_{\rm oct}(e^{sK_3}e_1)|^2=\cosh2s$.
\item In SI units, with $\psi=\Phi_{\rm oct}(c\bB;\bE)$ $[\mathrm{V\,m^{-1}}]$, $u=\tfrac{\varepsilon_0}2|\psi|^2$ is the energy density and
$\psi^{\!\top}(-1)^N\psi=\Re q_c$, $q_c=z_c\cdot z_c$, $z_c=c\bB+i\bE$ (\cref{en:SI}).
\end{enumerate}
\end{proposition}

\begin{proof}
The $f_k,\mathsf g_m$ are orthonormal, $(-1)^N=\mp1$ on them, $Tg_m=-f_m$ and $\Gamma=-(-1)^N$, which gives (i), (ii) and the formula for $E_0$ by
$|q|^2=(\Re q)^2+(\Im q)^2$. (iii) A form $M$ is invariant iff $\beta^{\!\top}M+M\beta=0$ for the fifteen bivectors, a linear system with a
two-dimensional solution space (r3-18, the count of \cref{thm:klein}); $T$ commutes with the antisymmetric (same-sign) and
anticommutes with the symmetric (mixed) bivectors (\cref{en:parity}), and $\Gamma$, which anticommutes with every vector, commutes
with all of them. The Mukai identity is checked on all $64$ pairs of basis forms, and $\dim(N_\sigma\cap N_\tau)$ is the number of agreeing signs.
(iv) $e_k$ and $e_{k+3}\omega$ are symmetric matrices, so $Q_k^{\!\top}+Q_k=2Q_k\neq0$; (iv), (v): r3-17, r3-18.
\end{proof}

Items (i), (ii), (iv) are fixed, (iii) is inherited in substance and verified here, (v) restates \cref{sec:energy}. By (i), the Maxwell
Lagrangian $L_M=-\tfrac12\Re q$ of \cref{r1:prop:maxwell} is $-\tfrac12\psi^{\!\top}(-1)^N\psi=\tfrac12(|\psi_{\bE}|^2-|\psi_{\bB}|^2)$ on the octants, half the
squared electric minus the squared magnetic octant amplitudes. The Chevalley form on the spinors
of $L\oplus L'$, antisymmetric for $\dim L=3$ since $(-1)^{3\cdot2/2}=-1$, and its non-vanishing on pure spinors exactly for transverse null
planes are Chevalley's \cite{Chevalley54}; Gualtieri identifies the form with the Mukai pairing (Section~2.4 of \cite{Gualtieri03}, whose
Proposition~2.21 restates the transversality criterion). So the Spin$(3,3)$-invariant mass form $T$ is, on the octants, the Chevalley--Mukai pairing of each octant with its
antipode: of the vacuum with the filled pole and of $f_k$ with $\mathsf g_k$, that is, of $E_k$ with $B_k$, with the sign fixed by (ii):
$\psi_{\bE}^{\!\top}T\psi_{\bB}=-\bE\cdot\bB=-\tfrac12\Im q$. The Lorentz invariant $\Re q$ appears on
the octants through the form $(-1)^N$, which is not $\mathrm{Spin}(3,3)$-invariant, while the invariant symmetric form $T\Gamma$ gives $\Im q$;
both facts are consistent with \cref{r3:thm:noeven} and with the representation $\rho_{\rm oct}$, which preserves $(-1)^N$.

% --- source: two_readings/r3_octants#8 (Maxwell's equations on the octants)
\subsection{Maxwell's equations on the octants}
\label{r3:sec:maxwell}

\begin{proposition}[Exterior derivative and codifferential]
\label{r3:prop:dk}
On $S$-valued fields put $d:=\sum_ia^+_i\partial_i$ and $\delta:=-\sum_ia^-_i\partial_i$. Then $d^2=\delta^2=0$,
\[
d-\delta=\sum_ie_i\partial_i=:D,\qquad d+\delta=\sum_ie_{i+3}\partial_i=:D',\qquad D^2=\Delta,\qquad D'^2=-\Delta,
\]
and under $f_k\leftrightarrow dx^k$, $\mathsf g_m\leftrightarrow\star dx^m$, $|\mathrm{vac}\rangle\leftrightarrow1$, $|\mathrm{full}\rangle\leftrightarrow dx^1\wedge dx^2\wedge dx^3$, $d$ and $\delta$
are the exterior derivative and the codifferential of the Euclidean $\R^3$: $d(\bE\cdot f)=(\curl\bE)\cdot g$, $d(\bB\cdot g)=(\nabla\cdot\bB)|\mathrm{full}\rangle$,
$\delta(\bE\cdot f)=-(\nabla\cdot\bE)|\mathrm{vac}\rangle$, $\delta(\bB\cdot g)=(\curl\bB)\cdot f$.
\end{proposition}

\begin{proof}
$d^2=\sum_{ij}a^+_ia^+_j\partial_i\partial_j=0$ because $a^+_ia^+_j$ is antisymmetric, and likewise for $\delta$; $a^+_i\pm a^-_i$ are $e_i$ and $e_{i+3}$, and
$D^2=\sum_ie_i^2\partial_i^2$, $D'^2=\sum_ie_{i+3}^2\partial_i^2$. The four formulas use $a^-_ig_m=\epsilon_{mik}f_k$ and $a^+_mg_m=|\mathrm{full}\rangle$ (r3-19).
\end{proof}

Inherited, and verified here: $D=d-\delta$ is K\"ahler's operator on inhomogeneous forms \cite{Kahler62}, the forms read as spinors \cite{Graf78}
and the Dirac--K\"ahler equation of Becher and Joos \cite{BecherJoos82}, whose Minkowski solutions Benn and Tucker read as a
fourfold-degenerate Dirac field \cite{BennTucker83}. Here the operator acts on space, and the forms carry the field strength, not a fermion.

\begin{theorem}[Maxwell's equations as one equation on the octants]
\label{r3:thm:maxwell}
For every smooth field $X(t,\mathbf x)$,
\begin{equation}
\label{r3:eq:maxwell}
(\partial_t-D)\,\Phi_{\rm oct}(X)=-(\nabla\cdot\bE)\,|\mathrm{vac}\rangle+(\partial_t\bE+\curl\bB)\cdot f+(\partial_t\bB-\curl\bE)\cdot g-(\nabla\cdot\bB)\,|\mathrm{full}\rangle .
\end{equation}
Consequently:
\begin{enumerate}[label=(\roman*),leftmargin=2.2em]
\item Maxwell's equations with sources in the frame variables, $\partial_tz=-i\curl z-i\bj$ and $\nabla\cdot z=i\rho$ (eq.~\cref{eq:RSsource}), are equivalent to
\[
(\partial_t-D)\,\Phi_{\rm oct}(X)=-\bigl(\rho\,|\mathrm{vac}\rangle+\bj\cdot f\bigr):
\]
the vacuum component is Gauss's law $\nabla\cdot\bE=\rho$, the filled component the absence of magnetic charge, the electric octants carry the
Amp\`ere--Maxwell law $\partial_t\bE=-\curl\bB-\bj$ and the magnetic octants Faraday's law $\partial_t\bB=\curl\bE$.
\item $(\partial_t+D)(\partial_t-D)=\partial_t^2-\Delta$ on all eight components, while $(\partial_t+D')(\partial_t-D')=\partial_t^2+\Delta$. In $D=\sum_ie_i\partial_i$ the
derivatives $\partial_i$ are taken along the three physical spatial coordinates $x^i$, the lengths of the classical reading, which sit on $V^-$;
the Clifford generators multiplying them are the temporal $e_1,e_2,e_3$, of square $+\id$, the generators that carry $\bB$ in the field reading,
and not the spatial ones.
\item In SI units $\bigl(\tfrac1c\partial_t-D\bigr)\Phi_{\rm oct}(c\bB;\bE)=-(c\rho\,|\mathrm{vac}\rangle+\bj\cdot f)/(\varepsilon_0c)$.
\item In textbook variables $\Phi_{\rm t}(X):=\bE\cdot f+\Btext\cdot g=-(-1)^N\Phi_{\rm oct}(X)$ satisfies $T\Phi_{\rm t}=\Phi_{\rm t}\circ G$ and
\[
(\partial_t+D)\,\Phi_{\rm t}(X)=\rho\,|\mathrm{vac}\rangle-\bj\cdot f,\qquad \Phi_{\rm t}(X)=\Bigl(\sum_kE_ke_k+\sum_kB_{{\rm t}k}\,Te_k\Bigr)|\mathrm{vac}\rangle ,
\]
which, under $S\cong\Cl(V^+)$ with $\sigma_k=e_k$ and $I=T$, is the equation $(\partial_t+\nabla)(\bE+I\Btext)=\rho-\bj$ of the Pauli algebra.
\item With pole components $\psi=\Phi_{\rm oct}(X)+\phi_0|\mathrm{vac}\rangle+\phi_3|\mathrm{full}\rangle$, $(\partial_t-D)\psi=0$ reads $\partial_t\phi_0=\nabla\cdot\bE$, $\partial_t\phi_3=\nabla\cdot\bB$,
$\partial_t\bE=-\curl\bB+\nabla\phi_0$, $\partial_t\bB=\curl\bE+\nabla\phi_3$; every component of a solution obeys the wave equation, so the Gauss constraints
hold for all times if $\phi_0,\phi_3$ and the two divergences vanish initially.
\end{enumerate}
\end{theorem}

\begin{proof}
\Cref{r3:eq:maxwell} is $\partial_t\Phi_{\rm oct}(X)-(d-\delta)\Phi_{\rm oct}(X)$ with the four formulas of \cref{r3:prop:dk}; it is checked on general symbolic fields
$E_k(t,\mathbf x)$, $B_k(t,\mathbf x)$ (r3-20), and the real and imaginary parts of the residuals of eq.~\cref{eq:RSsource} are exactly the
components of $(\partial_t-D)\Phi_{\rm oct}(X)+\rho|\mathrm{vac}\rangle+\bj\cdot f$. (ii) $(\partial_t+D)(\partial_t-D)=\partial_t^2-D^2$ (r3-21). (iii) is (i) with $\Phi_{\rm oct}(\bE,c\bB)$ and
$\partial_t\bE=-c^2\curl\bB-\bj/\varepsilon_0$, $\nabla\cdot\bE=\rho/\varepsilon_0$ (r3-24, on first-order jets). (iv) $\Phi_{\rm t}(X)=\bE\cdot f-\bB\cdot g$; $(-1)^N$
anticommutes with $D$; $Te_k=e_je_l$ cyclically gives the second form; under $a|\mathrm{vac}\rangle\leftrightarrow a$ left multiplication by $e_k$ is that by $\sigma_k$ (r3-24). (v) is
\cref{r3:eq:maxwell} together with $D(\phi_0|\mathrm{vac}\rangle)=\nabla\phi_0\cdot f$ and $D(\phi_3|\mathrm{full}\rangle)=\nabla\phi_3\cdot g$ (r3-21); the wave equation follows from (ii), and
$\phi_0$ with $\phi_0=\partial_t\phi_0=0$ at $t=0$ vanishes identically.
\end{proof}

Fixed in the realisation and inherited in substance. Item (ii) is a statement about the representation and not about spacetime: the
octant form of Maxwell's equations uses the Pauli triple $e_1,e_2,e_3$ because it is the triple of square $+\id$ that makes the operator hyperbolic,
and this choice is constructed. It does not show that spatial differentiation follows from three temporal directions, and the spatial
coordinates $x^i$ remain the lengths of $V^-$; the Dirac operator of the particle reading couples the same derivatives to the spatial
generators through $\gamma^0\gamma^a$ instead (\cref{r4:sec:main}). Item (iv) is the Pauli-algebra form, obtained by multiplying with $\gamma_0$, of the
spacetime-algebra equation $\nabla F=J$ of Hestenes \cite{Hestenes66}, which \cref{thm:nablaF} realises on the interior four-plane; here
it is realised on the temporal triple acting on the vacuum. The element $T$ that plays the pseudoscalar $I$ is $I_T$,
with $I_T^2=-\id$ and $[I_T,e_i]=0$ for the temporal $e_i$ (r3-24). The circularly polarised wave of \cref{ex:circular} solves $(\partial_t-D)\Phi_{\rm oct}(X)=0$, and the same wave with $\bB$ reversed does
not (r3-23). The textbook signs, $\partial_t\bE=\curl\Btext-\bj$ and $\partial_t\Btext=-\curl\bE$, become those of (i) in the frame
variable $\bB=-\Btext$.

\begin{proposition}[Maxwell fixes the scales and leaves a duality angle]
\label{r3:prop:maxunique}
Let $\Phi_M(X)=(a\bE+b\bB)\cdot f+(c\bE+d\bB)\cdot g$ be the general rotation-equivariant map into the field octants. Then $(\partial_t-D)\Phi_M(X)$ vanishes
for every source-free Maxwell field iff $c=-b$ and $d=a$, that is iff $\Phi_M=r\,\Phi\circ e^{\vartheta G}$ for some $r\ge0$ and $\vartheta$. The map puts
$\bE$ and $\bB$ on single weights iff $\vartheta\in\tfrac\pi2\Z$: $\Phi$ for $\vartheta=0$ and $\Phi\circ G$, with $\bE$ on weight two, for $\vartheta=\tfrac\pi2$.
\end{proposition}

\begin{proof}
Substituting $\partial_t\bE=-\curl\bB$, $\partial_t\bB=\curl\bE$ and the vanishing divergences, the residual has $f$-component $(d-a)\curl\bB+(b+c)\curl\bE$,
$g$-component $(d-a)\curl\bE-(b+c)\curl\bB$ and no pole component; it vanishes on all first-order Maxwell jets iff $d=a$, $c=-b$ (r3-22). Conversely
$\Phi\circ e^{\vartheta G}$ maps solutions to solutions, since $e^{\vartheta G}$, the duality rotation $z\mapsto e^{i\vartheta}z$, preserves the source-free system.
\end{proof}

Fixed. The two scales of \cref{r3:thm:phi}(i) are tied by Faraday's and Amp\`ere's laws, $\beta=\alpha$ in Heaviside--Lorentz units and $\beta=c\,\alpha$
in SI, and the energy normalisation of \cref{r3:prop:energies}(i) fixes $\alpha$ up to sign. What remains is the duality angle. With electric
sources, $\Phi$ places the charge density on the vacuum pole and the current on the electric octants, and $\Phi\circ G$ places them on the
filled pole and on the weight-two octants; the convention that all charges are electric (\cref{sec:spin}, \cite{Jackson}) fixes which
octants carry the current once the labelling is chosen, not the labelling.

\begin{remark}[Octant pooling]
\label{r3:rem:pooling}
Each $e_i$ maps the octant $\sigma$ to the octant with $\sigma_i$ reversed, the image of $\sigma$ under the octant generator $R_i$, and
$e_{i+3}=e_i\omega_i$ acts on the same pair of octants (r3-25). Hence the $\sigma$-component of $D\psi$ collects $\partial_i$ of the three neighbouring
octants only: the Gauss law on the vacuum collects $\sum_i\partial_iE_i$ from the three adjacent electric octants, and each evolution equation
collects a curl from two octants of the other weight. This is the exact content of the proposed ``pooling'' of Maxwell's equations across
octants, and it holds as $D=\sum_ie_i\partial_i$. The proposed mechanism, a per-pair exchange of $e_i$ with $e_{i+3}$ converging on the pole
$(-,-,-)$, is not confirmed: the exchange replaces $D$ by $D'$ and the wave operator by the four-dimensional Laplacian
(\cref{r3:thm:maxwell}(ii)), and nothing converges on one pole, the two poles carrying one Gauss law each.
\end{remark}

% --- source: two_readings/r3_octants#9 (The answer: X\_EM = (B;E) with each component on an octant)
\subsection{\texorpdfstring{The answer: $X_{\rm EM}=(\bB;\bE)$ with each component on an octant}{The answer: X\_EM = (B;E) with each component on an octant}}
\label{r3:sec:answer}

The pair reading has one exact realisation. The field-octant map (\cref{r3:thm:phi,r3:prop:maxunique}),
$\Phi_{\rm oct}(X)=\sum_k(E_k+B_kT)e_k|\mathrm{vac}\rangle$, puts each $E_k$ on the octant with a single minus sign in pair $k$ and each $B_k$ on its antipode; it
is unique up to a scale once rotation equivariance, Maxwell's equations and the separation of $\bE$ from $\bB$ are required, apart from the
alternative $\Phi\circ G$ that the earlier labelling uses; in textbook variables the same octants carry, through $\Phi_{\rm t}=-(-1)^N\Phi$
(\cref{r3:thm:maxwell}(iv)), Hestenes' $\bE+I\Btext$ on the temporal Pauli algebra, with $I=T$. Clifford multiplication by the field point does
not realise the pair: it places the light-cone sums $\bB\pm\bE$ on the octants (\cref{r3:prop:poles}), and differs from $\Phi$ by an eighth of a turn
of duality. The construction is a representation of the classical field strength; it is not a quantisation of the field.

\paragraph{The main obstruction.} On the octants the rotations act through the spin representation, but the boosts and the duality rotation
of the field act only through the odd operators $Q_k$ and $-T$ of $\rho_{\rm oct}$ (\cref{r3:thm:noeven,r3:prop:odd}), whereas a Dirac field transforms
under the spin lift of the interior Lorentz algebra. A covariant theory containing both the Dirac field and the field octants must say how
$\rho_{\rm oct}$, the spin lift acting on Dirac spinors and the map $\Phi$ coexist. This is the principal question the octant picture leaves, and the
element $T$, which is at once the mass form of the Dirac reading and the duality operator on the octants, is where the two meet
(\cref{r4:sec:main}).

\paragraph{What the poles carry.} With $\Phi$, the vacuum pole carries Gauss's law and the charge density, the filled pole the absence of
magnetic charge (\cref{r3:thm:maxwell}); they are the Fock vacuum and the filled state, the rotation scalars, the eigenlines of $\mathfrak{gl}(3)$ of
weights $\Lambda^0$ and $\Lambda^3$, and, under the spin-lifted $\lG$, the time components of the two four-vectors $S^\mp$ (\cref{r3:prop:odd}(iv)).
The mass form pairs them with each other as it pairs $E_k$ with $B_k$ (\cref{r3:prop:energies}). Of the earlier labels, ``statistics'' survives as
the constructed correspondence between the Fisher--Rao octant of clocks and the vacuum pole $(+,+,+)$ (\cref{r3:prop:clock}). ``Mass'' does not
survive as a property of the pole $(-,-,-)$: the mass form treats both poles alike, and the interior Lorentz algebra distinguishes no octant,
since the stabiliser in $\lint$ of the octant $\sigma$ is the three-dimensional span of $e_1e_4=\omega_1$ and the null rotations
$(e_1-\sigma_1e_4)e_5$, $(e_1-\sigma_1e_4)e_6$, depending on $\sigma_1$ only, while $g_2,g_3$ commute with $\lint$ and $g_1$ normalises it (r3-26). The
octant Fock structure is the canonical anticommutation algebra of one constituent; it is not a photon Fock space, and the quantisation of the
field stays open.

\paragraph{Open.} Whether a dynamics, rather than a convention, selects $\Phi$ or $\Phi\circ G$; what the two scalar fields $\phi_0,\phi_3$ of the
extended system of \cref{r3:thm:maxwell}(v) mean, if anything; and how $\rho_{\rm oct}$, which is not part of the spin representation, enters a covariant
dynamics with the Dirac sector of \cref{ch:ch02_lorentz}; its covariant couplings linear in the field are determined in \cref{r4:thm:couplings}.

\Cref{r3:tab:claims} lists every claim of the section with its check and its class.

{\small\renewcommand{\arraystretch}{1.12}
\begin{longtable}{@{}>{\raggedright\arraybackslash}p{0.09\textwidth}>{\raggedright\arraybackslash}p{0.56\textwidth}>{\raggedright\arraybackslash}p{0.27\textwidth}@{}}
\caption{The claims of this section, with their checks and classes.}\label{r3:tab:claims}\\
\toprule
id & claim & class\\
\midrule
\endfirsthead
\toprule
id & claim & class\\
\midrule
\endhead
r3-1 & octants $=$ basis lines, sign triples, CAR, $N=\#$ minus signs, vacuum $(+,+,+)$ & restated (\cref{prop:grading}), fixed\\
r3-2 & octant $\sigma$ is the pure spinor of $N_\sigma=\operatorname{span}\{e_i-\sigma_ie_{i+3}\}$ & fixed; inherited (Cartan--Chevalley)\\
r3-3 & $\Gamma=-\omega_1\omega_2\omega_3$, $(-1)^N=-\Gamma=\gamma_5$; the earlier $I=-\omega_1\omega_2\omega_3$ confirmed & restated, fixed\\
r3-4, 5 & reflections lift to $e_a\omega$, unique; octant group simply transitive; lifted group of order $16$, centre $\{\pm\id,\pm e_4e_5e_6\}$ & fixed\\
r3-6 & top eigenline of $B_{\operatorname{diag}x}$ is the octant $\operatorname{sign}x$; clock octant $\leftrightarrow$ vacuum & fixed; identification constructed\\
r3-7 & $S\cong\Lambda^*(L)\cong\Cl(V^+)|\mathrm{vac}\rangle$, weights $1,3,3,1$; $X$ acts by $2X_L\wedge$ and contraction with $\eta(X_{L'},\cdot)$ & inherited (Chevalley), verified\\
r3-8 & commutant of $\cK$ $=\mathfrak{gl}(3)$, acting as $\Lambda^w(A)-\tfrac12\tr A$ & inherited (Gualtieri), fixed here\\
r3-9 & diagonal rotations $L_i=\tfrac12(e_{j+3}e_{k+3}-e_je_k)$, Fock forms of $L_i,K_i$ & restated, fixed\\
r3-10 & equivariant maps: dimension $4$; $\Phi_{{\rm oct},\alpha,\beta}$ & fixed; $\Phi$ constructed\\
r3-11 & premetric $\mathfrak{gl}(3)$ on the image; vector action mixes $\bE,\bB$ & fixed; premetric reading inherited\\
r3-12 & $T$ lifts parity and acts as $-G$ on the octants; $(-1)^N$ acts as parity & fixed\\
r3-13 & $X|\mathrm{vac}\rangle$, $X|\mathrm{full}\rangle$ light-cone amplitudes; \cref{r3:eq:eighth}; $\Phi$ is no Clifford map & fixed; reading interpretive\\
r3-14 & no even intertwiner of boosts or of $G$ & obstructed\\
r3-15, 16 & odd intertwiners $Q_k,-T$; $\rho_{\rm oct}$; spin lift: four-vectors, no intertwiner & $Q_k$, $\rho_{\rm oct}$ constructed; counts fixed; restated (iv)\\
r3-17, 18 & energies; invariant forms; $T=$ Chevalley--Mukai pairing of antipodes; $|\psi|^2$ not $\rho_{\rm oct}$-invariant & fixed; pairing inherited\\
r3-19 & $d,\delta$ on the octants, $d\mp\delta=D,D'$ & inherited (K\"ahler), verified\\
r3-20, 21, 23, 24 & Maxwell $\Leftrightarrow(\partial_t-D)\Phi_{\rm oct}(X)=-(\rho|\mathrm{vac}\rangle+\bj\cdot f)$; wave operator; SI; Hestenes form & fixed; inherited in substance\\
r3-22 & Maxwell-compatible equivariant maps $=r\Phi\circ e^{\vartheta G}$ & fixed\\
r3-25 & pooling $=$ nearest-neighbour structure of $D$ & fixed; mechanism refuted\\
r3-26 & $\lint$ distinguishes no octant & fixed; ``mass pole'' refuted\\
\bottomrule
\end{longtable}}

% --- source: two_readings/r4_bridge#0 ((head))
% Track r4 -- the bridge T = nJ(n) between the two readings.
% New references: sections/bib_r4.tex.

% --- source: two_readings/r4_bridge#1 (The bridge T = nJ(n) between the two readings)
\section{\texorpdfstring{The bridge $T=\hat nJ(\hat n)$ between the two readings}{The bridge T = nJ(n) between the two readings}}
\label{r4:sec:main}

Two sections of this chapter meet in one element of $\Cl(3,3)$. The clock $\hat n\in V^+$ of \cref{r2:sec:main} determines the complex structure
$J(\hat n)=T\hat n$ of the Dirac reading (\cref{r2:prop:J,r2:thm:i}), while the octant representation of the field in \cref{r3:sec:main} uses $T=e_1e_2e_3$
itself, through $-T$, as the duality operator (\cref{r3:prop:parity,r3:prop:odd}). Since $\hat n$ commutes with $T$ and $\hat n^2=\id$,
$T=\hat nJ(\hat n)=J(\hat n)\hat n$ for every clock. This section asks whether this identity can be built into both readings at once while their distinct
Lorentz representations are kept apart, also when the clock varies or the temporal axes are permuted, and it answers the points left open
in \cref{r3:sec:main}: the remark on the spatial derivatives of the octant Maxwell operator after \cref{r3:thm:maxwell}, the paragraph ``The
main obstruction'' of \cref{r3:sec:answer}, and the question of its paragraph ``Open'' whether $\rho_{\rm oct}$ can enter a coupling to the Dirac
sector. Three expectations are tested on the way: (E1) that an odd permutation of the temporal axes reverses the mass form and the duality
operator together; (E2) that the couplings of the Dirac field to the field, which contain $T$, change sign with it under such a permutation;
(E3) that such a coupling can be written through $T$ and $J(\hat n)$. In brief:
(1) $T$ is invariant and $J(\hat n)$ covariant under every rotation of the temporal triple and under the diagonal rotations; an odd permutation of the
temporal axes reverses $T$ and $J(\hat n)$ together, hence the mass form; paired with the same permutation of the spatial axes, the only form that
respects $G$, it also reverses the duality operator and turns the octant map into its textbook-orientation form; reversing the clock is a
half-turn, which fixes $T$, and the charge conjugation of \cref{sec:discrete} is a half-turn composed with $T$ (\cref{r4:sec:identity});
(2) $T$ is at once the Dirac mass form, the Dirac parity and the lift of the frame parity, and the duality on the octants; $J(\hat n)$ acts on the
octants as a rotation of the transverse field about $\hat n$ and maps the longitudinal field to the two poles that carry the Gauss laws
(\cref{r4:sec:roles});
(3) the spin lift of $\lint(\hat n)$ and the octant representation $\rho_{\rm oct}$ are inequivalent on the same module, their joint commutant is $\R\id$,
no covariant bilinear coupling exists, and the covariant trilinear couplings of a Grassmann Dirac field to the field octants form a
two-dimensional space, the Pauli term and its dual, written through $T$ and $J(\hat n)$ (\cref{r4:sec:reps});
(4) the two first-order operators use conjugate Pauli triples, related only by elements of mixed parity, and neither derives spatial
differentiation from temporal directions (\cref{r4:sec:operators});
(5) the field reading does not depend on the clock and the particle reading does, and $T$ is, up to scale, the only element that is the clock
times a Dirac complex structure for every clock (\cref{r4:sec:clock}). \Cref{r4:sec:answer} collects the verdicts.

\paragraph{Conventions and units.} Those of \cref{r2:sec:main,r3:sec:main}: $V=\R^{3,3}$, $\eta=\operatorname{diag}(1,1,1,\allowbreak-1,-1,-1)$, the
realisation of eq.~\cref{eq:realisation}, $s_a:=e_{a+3}$, $T:=e_1e_2e_3$, $\Gamma:=\omega=e_1\cdots e_6$, the binding root $Ge_i=e_{i+3}$,
$Ge_{i+3}=-e_i$, the diagonal rotations $L_i=\tfrac12(e_{j+3}e_{k+3}-e_je_k)$ and boosts $K_i=\tfrac12(e_ke_{j+3}-e_je_{k+3})$ of $\lG$ for cyclic
$(i,j,k)$, acting on $V$ by the commutator, and the octant notation $|\mathrm{vac}\rangle$, $|\mathrm{full}\rangle$, $f_k$, $\mathsf g_m$, $\Phi_{\rm oct}(X)=\bE\cdot f+\bB\cdot g$ of
\cref{r3:prop:octants,r3:lem:rot,r3:thm:phi}; for $v\in\R^3$ put $f(v)=\sum_kv_kf_k$ and $\mathsf g(v)=\sum_mv_mg_m$. The field point is
$X=\sum_i(B_ie_i+E_ie_{i+3})$ with $\bB$ minus the textbook magnetic field, $\Btext=-\bB$. A \emph{clock} is a unit $\hat n\in V^+$ with an oriented
orthonormal frame $(\hat n,u_1,u_2)$ of $V^+$, so $nu_1u_2=T$ (\cref{r2:prop:J}); $\gamma^\mu_{\hat n}=(\hat n,s_1,s_2,s_3)$, $\lint(\hat n)$ is spanned by the
bivectors of $E_{\hat n}=\R \hat n\oplus V^-$, $\Omega_{\hat n}:=ns_1s_2s_3$, and $\phi_{\hat n}:\lG\to\lint(\hat n)$ is the dictionary of \cref{r2:prop:dict},
$L_i\mapsto s_j\wedge s_k$, $K_i\mapsto s_i\wedge \hat n$. A bivector $a\wedge b$ of orthogonal vectors is the Clifford element $\tfrac12ab$, and its spin
lift is left multiplication by that element, so the Dirac reading's Lorentz representation of $\lG$ on $S$ is
\[
\rho_{\hat n}:=\text{spin}\circ\phi_{\hat n},\qquad \rho_{\hat n}(L_i)=\tfrac12s_js_k,\qquad \rho_{\hat n}(K_i)=\tfrac12s_in ,
\]
and the field reading's is $\rho_{\rm oct}(L_i)=L_i$, $\rho_{\rm oct}(K_i)=Q_i=-TL_i$ (\cref{r3:prop:odd}). Further $t_i:=\tfrac12e_je_k$ spans the
rotations $\so3_+$ of the temporal triple alone, and $A_i:=\tfrac12(e_je_k+e_{j+3}e_{k+3})$ rotates the two triples in opposite senses
(\cref{r3:lem:rot}); $L(\hat n):=\sum_in_iL_i$, $A(\hat n):=\sum_in_iA_i$. Dimensions: $\hat n$, $J(\hat n)$, $T$, $\Gamma$, all Clifford elements and all
intertwiners are dimensionless; $[x^a]=[x^0]=[ct]=\mathrm m$ and $[\partial_\mu]=\mathrm m^{-1}$; in SI units a Dirac field has $[\psi]=\mathrm m^{-3/2}$ and
a Lagrangian density has the unit $\mathrm{J\,m^{-3}}$, as in \cref{r1:eq:LD}; the field enters as $X_{\rm EM}=(c\bB;\bE)$ with $[E_i]=[cB_i]=\mathrm{V\,m^{-1}}$; a
magnetic moment $\mu$ has $[\mu]=\mathrm{J\,T^{-1}}$ and an electric dipole moment $d$ has $[d]=\mathrm{C\,m}$. Computations use Heaviside--Lorentz units
with $\hbar=c=1$. Identifiers r4-$\hat n$ refer to \texttt{r4\_bridge.py}, which checks every finite-dimensional and symbolic claim below
in exact arithmetic (r4-1 re-verifies the conventions).

% --- source: two_readings/r4_bridge#2 (The identity and its covariance)
\subsection{The identity and its covariance}
\label{r4:sec:identity}

\begin{proposition}[The identity]
\label{r4:prop:identity}
For every $\hat n\in V^+$: $\hat n^2=|\hat n|^2\id$, $J(\hat n):=T\hat n=\hat nT=\hat n_1e_2e_3+\hat n_2e_3e_1+\hat n_3e_1e_2$, $J(\hat n)^2=-|\hat n|^2\id$, and $\hat nJ(\hat n)=J(\hat n)\hat n=|\hat n|^2\,T$. Hence for every clock
\[
T=\hat nJ(\hat n)=J(\hat n)\hat n,\qquad J(\hat n)=T\hat n=\hat nT,\qquad \hat n^2=\id,\qquad J(\hat n)^2=T^2=-\id .
\]
$J(\hat n)$ commutes with $\hat n,s_1,s_2,s_3$, and $\hat n\mapsto J(\hat n)$ is linear and injective.
\end{proposition}

\begin{proof}
By the blade rule (\cref{lem:blade}) $T$ commutes with each $e_i$ and anticommutes with each $s_a$. So $J(\hat n)^2=T\hat nT\hat n=T^2\hat n^2$ and
$\hat nJ(\hat n)=nnT=|\hat n|^2T$; $T$ and $\hat n$ each anticommute with $s_a$, so $J(\hat n)$ commutes with $s_a$ (r4-2, symbolic in $\hat n$ and on the rational parametrisation
$\hat n=(2a,2b,1-a^2-b^2)/(1+a^2+b^2)$ of the clock sphere).
\end{proof}

Restated from \cref{r2:prop:J} and from \cref{rem:ambient} ($J=\gamma^0T$), re-verified (r4-2). 

\begin{proposition}[Rotations]
\label{r4:prop:rot}
\begin{enumerate}[label=(\roman*),leftmargin=2.2em]
\item For $Y\in\so3_+\oplus\so3_\Delta$, $[Y,T]=0$ and $[Y,J(\hat n)]=J([Y,\hat n])$, with $[Y,\hat n]\in V^+$. Hence for the spin lift $R$ of a rotation of the temporal
triple alone, or of a diagonal rotation, $RTR^{-1}=T$ and $RJ(\hat n)R^{-1}=J(R\hat n)$, where $R\hat n:=R\hat nR^{-1}$.
\item The rotations $\so3_+$ do not commute with $G$; the diagonal rotations do, and are symmetries of the field reading: $\rho_{\rm oct}(L_i)=L_i$
and $\Phi$ is equivariant (\cref{r3:thm:phi}).
\item A rotor $q\in\mathrm{Spin}(V^+)$ with $qe_1q^{-1}=\hat n$ fixes $T$, the $s_a$ and $e_4e_5e_6$, and it carries $J(e_1)$, $\Omega_{e_1}$, $\lint(e_1)$,
$\phi_{e_1}$ to $J(\hat n)$, $\Omega_{\hat n}$, $\lint(\hat n)$, $\phi_{\hat n}$.
\end{enumerate}
\end{proposition}

\begin{proof}
$T$ commutes with every same-sign bivector, hence with $t_i$ and $L_i$, and $[Y,T\hat n]=T[Y,\hat n]$. The spatial part of $L_i$ commutes with $\hat n$, so
$[L_i,\hat n]=-[t_i,\hat n]\in V^+$. A temporal bivector commutes with every $s_a$ and with $e_4e_5e_6$; with $\phi_{\hat n}(L_i)=s_j\wedge s_k$, $\phi_{\hat n}(K_i)=s_i\wedge \hat n$ this
gives (iii). (ii) is \cref{r3:lem:rot}. The finite statements are checked with the rational rotors $q=u'(e_1-\hat n)$, $u'\perp \hat n$, at
$\hat n=(2,3,6)/7,(1,2,2)/3,(-1,4,8)/9$ and with an exact Cayley transform of $L_3$ (r4-3, r4-4).
\end{proof}

Fixed. Rotations of the temporal triple move the clock and with it the whole Dirac structure, but they are not symmetries of the field reading;
the diagonal rotations move the clock as well and are symmetries of both.

\begin{proposition}[Permutations of the temporal axes and the reversal of the clock]
\label{r4:prop:perm}
For $\sigma\in S_3$ let $A_\sigma$ be the automorphism of $\Cl(3,3)$ with $e_i\mapsto e_{\sigma(i)}$, $s_a\mapsto s_a$, and $A'_\sigma$ the one with
$e_i\mapsto e_{\sigma(i)}$, $s_a\mapsto s_{\sigma(a)}$; write $\sigma \hat n=\sum_in_ie_{\sigma(i)}$.
\begin{enumerate}[label=(\roman*),leftmargin=2.2em]
\item Both are inner, $A=\Ad_U$ with $U$ unique up to scale and orthogonal up to scale, and
\[
A(T)=\operatorname{sgn}(\sigma)\,T,\qquad A(J(\hat n))=\operatorname{sgn}(\sigma)\,J(\sigma \hat n),\qquad A(\hat n)=\sigma \hat n .
\]
\item $A'_\sigma$ commutes with $G$, permutes the pole operators, $A'(\omega_i)=\omega_{\sigma(i)}$, and with $U|\mathrm{vac}\rangle=|\mathrm{vac}\rangle$ it acts on
the octants by $U\Phi_{\rm oct}(X)=\Phi_{\rm oct}(\tilde X)$, $\tilde X=(\operatorname{sgn}(\sigma)P_\sigma\bB;P_\sigma\bE)$, where $P_\sigma$ permutes the components: $\bB$ is
transformed as an axial vector. The transported octant map $\Phi':=U\Phi_{\rm oct}(P_\sigma\oplus P_\sigma)^{-1}$ is $\Phi$ for even $\sigma$ and
$\Phi\circ(-\eta)=\Phi_{\rm t}=-(-1)^N\Phi$, the textbook-orientation map of \cref{r3:thm:maxwell}(iv), for odd $\sigma$, and its duality operator is
$A'(-T)$: $A'(-T)\Phi'=\Phi'G$.
\item No nontrivial $A_\sigma$ commutes with $G$.
\item Let $R=nu_2$, the rotor of the half-turn about $u_1$. Then $R\hat nR^{-1}=-\hat n$, $RTR^{-1}=T$, $RJ(\hat n)R^{-1}=J(-\hat n)=-J(\hat n)$, and $RT=u_1$: every unit vector
$C$ of the surplus plane, the charge conjugation of \cref{thm:C} at the clock $\hat n$, is the half-turn about $C$, which reverses the
clock, composed with $T$.
\end{enumerate}
\end{proposition}

\begin{proof}
(i) $\Cl(3,3)\cong M_8(\R)$ is central simple (Skolem--Noether); $e_{\sigma(1)}e_{\sigma(2)}e_{\sigma(3)}=\operatorname{sgn}(\sigma)T$ by reordering anticommuting factors,
and $A(J(\hat n))=A(T)A(\hat n)$. (ii) $P_\sigma\oplus P_\sigma$ commutes with the block matrix of $G$. $A'(a^+_k)=a^+_{\sigma(k)}$, so $Uf_k=f_{\sigma(k)}$ and
$Ug_m=\tfrac12\epsilon_{mjk}a^+_{\sigma(j)}a^+_{\sigma(k)}|\mathrm{vac}\rangle=\operatorname{sgn}(\sigma)g_{\sigma(m)}$, since $\epsilon_{\sigma(m)\sigma(j)\sigma(k)}=\operatorname{sgn}(\sigma)\epsilon_{mjk}$.
Hence $\Phi'(Y)=\Phi_{\rm oct}(\operatorname{sgn}(\sigma)\bB_Y;\bE_Y)$, and $\Phi\circ(-\eta)=-(-1)^N\Phi$ because $(-1)^N$ is $-1$ on the $f$'s and $+1$ on the $g$'s; the duality
relation is the transport of $-T\Phi=\Phi G$. (iii) A temporal permutation alone moves $e_i$ and fixes its partner $e_{i+3}$. (iv) $R$ is an even
element of $\Cl(V^+)$, so it commutes with $T$; $nu_2nu_1u_2=-u_2u_1u_2=u_1$. All of this is checked for the six permutations, alone and paired, and
for three clocks (r4-5, r4-6, r4-7).
\end{proof}

Fixed. Expectation (E1) is confirmed in the form (i): an odd permutation reverses $T$ and $J(\hat n)$ together, and the paired one, which is
the only one that respects $G$ and the octants, reverses the mass form $T$ and the duality operator $-T$ together. The reason is (ii): an odd
permutation of the axes is an improper rotation, under which $\bB$ is axial; on the octants it changes the orientation convention $\Btext=-\bB$ of \cref{sec:conventions} into the textbook one, for which the duality operator is $+T$ (\cref{r3:thm:maxwell}(iv)). Reversing $\hat n$ reverses $J(\hat n)$, and (iv) sharpens this: the reversal is a rotation in $\mathrm{Spin}(V^+)$, which leaves $T$, the product of the two reversed objects, fixed.

\begin{table}[t]
\centering
\footnotesize
\setlength{\tabcolsep}{3pt}
\renewcommand{\arraystretch}{1.15}
\begin{tabular}{@{}>{\raggedright\arraybackslash}p{2.6cm}>{\raggedright\arraybackslash}p{1.9cm}>{\raggedright\arraybackslash}p{5.2cm}>{\raggedright\arraybackslash}p{5.0cm}@{}}
\toprule
operation & $T$; $J(\hat n)$ & Dirac reading: rest flow $e^{-m\lambda T}$, senses & field reading: $-T$, $G$, $\Phi$\\
\midrule
rotation $R$ of $V^+$ alone & $T$; $J(R\hat n)$ & flow invariant; eigenspaces $\hat n=\pm1$ go to $R\hat n=\pm1$, senses kept & not a symmetry: $[\so3_+,G]\neq0$\\
diagonal rotation & $T$; $J(R\hat n)$ & as above, with $V^-$ rotated alike & symmetry: $\rho_{\rm oct}(L)=L$, $-T$, $G$, $\Phi$ preserved\\
clock reversal (half-turn) & $T$; $-J(\hat n)$ & flow invariant; the eigenspaces $\hat n=\pm1$ are exchanged, so the particle and antiparticle labels relative to $\hat n$ are exchanged & the octant structure does not involve $\hat n$; the half-turn itself is not a field symmetry\\
odd $\sigma$, temporal alone & $-T$; $-J(\sigma \hat n)$ & $e^{-m\lambda T}\mapsto e^{+m\lambda T}$: sense reversed on each eigenspace; $(m,q)\mapsto(-m,-q)$ & not a symmetry: breaks $G$\\
odd $\sigma$, paired & $-T$; $-J(\sigma \hat n)$ & as for $\sigma$ alone, spatial axes permuted & $G$ kept, octants permuted, $\Phi\mapsto\Phi_{\rm t}$, duality $-T\mapsto+T$\\
$C=RT$, $C\in U_{\hat n}$ & $T$; $-J(\hat n)$ & same flow; particle $\leftrightarrow$ antiparticle (\cref{mb:senses}) & not a field symmetry ($C$ is a vector)\\
\bottomrule
\end{tabular}
\caption{What the operations of \cref{r4:prop:rot,r4:prop:perm} do to $T$, $J(\hat n)$ and the two readings (r4-3 to r4-8).}
\label{r4:tab:ops}
\end{table}

\begin{corollary}[The two senses at every clock]
\label{r4:cor:senses}
On the eigenspaces $\hat n=\pm1$ of any clock, $T=\pm J(\hat n)$, so the rest flow $e^{-m\lambda T}$ of \cref{mb:restflow} is $e^{\mp m\lambda J(\hat n)}$, the two
senses of the charge rotation that are particle and antiparticle there (\cref{mb:senses}). An odd temporal permutation carries the
Dirac data $(K^\mu_{\hat n},T,J(\hat n))$, with $K^\mu_{\hat n}=-\hat n\gamma^\mu_{\hat n}$, to $(K^\mu_{\sigma \hat n},-T,-J(\sigma \hat n))$: the transported Lagrangian \cref{r1:eq:LD} has mass $-m$ and
charge $-q$. The axial rotation by $\Omega_{\sigma \hat n}$, which preserves every $K^\mu$ and $J$ and reverses $T$ (\cref{mb:massless}(iv)), restores the
mass, leaving $q\mapsto-q$, the effect of $C$.
\end{corollary}

\begin{proof}
With $\Pi_\pm=\tfrac12(\id\pm \hat n)$, $T\Pi_\pm=\tfrac12(T\pm J(\hat n))=\pm J(\hat n)\Pi_\pm$ because $J(\hat n)\hat n=T$. The rest follows from \cref{r4:prop:perm}(i) and from
$\Omega^{\!\top}T\Omega=-T$, $\Omega^{\!\top}K^\mu\Omega=K^\mu$, $[\Omega,J]=0$ (r4-8).
\end{proof}

\cref{mb:restflow} is restated at every clock (fixed); the effect of the permutations on the Dirac reading is fixed; that a Dirac mass
sign is not observable alone is inherited \cite{ItzyksonZuber}. In the decomposition $C=RT$ of \cref{r4:prop:perm}(iv), $T$ preserves the flow,
the eigenspaces of $\hat n$ and $J(\hat n)$ (it is a parity, \cref{r4:thm:roles}), so the reversal of the sense by $C$ in \cref{mb:senses} comes entirely
from the half-turn that reverses the clock (fixed; the reading of particle and antiparticle as the two clock orientations is interpretive).

% --- source: two_readings/r4_bridge#3 (The three roles of T, and J(n) on the octants)
\subsection{\texorpdfstring{The three roles of $T$, and $J(\hat n)$ on the octants}{The three roles of T, and J(n) on the octants}}
\label{r4:sec:roles}

\begin{theorem}[The three roles of $T$]
\label{r4:thm:roles}
\begin{enumerate}[label=(\alph*),leftmargin=2.2em]
\item \emph{Mass form.} $\gamma^0_{\hat n}J(\hat n)=T$ for every clock: the mass kernel of the Dirac reading is the same at every clock, while the kinetic kernels
$K^\mu_{\hat n}=TJ(\hat n)\gamma^\mu_{\hat n}=-\hat n\gamma^\mu_{\hat n}$ move with it.
\item \emph{Parity.} $\Ad_T|_V=\eta$, the frame parity fixing $\bB$ and reversing $\bE$, and $T$ is up to scale the only element implementing it. At every
clock $T$ commutes with $J(\hat n)$ and $TH_{\hat n}(\bp)T^{-1}=H_{\hat n}(-\bp)$ for $H_{\hat n}(\bp)=\sum_ans_ap_a+mn$: $T=\gamma^0J$ is the admissible Dirac parity with $P^2=-\id$
of \cref{thm:PT}, at every clock.
\item \emph{Duality.} $-T\Phi=\Phi G$, $Tf_k=\mathsf g_k$, $Tg_k=-f_k$, $T|\mathrm{vac}\rangle=|\mathrm{full}\rangle$: on the field octants $-T$ is the Hodge star.
\end{enumerate}
The complex structure $J(\hat n)$ is, for the clock $\hat n$, the selected imaginary unit and the charge generator of the Dirac reading (\cref{thm:i} at every clock, \cref{r2:thm:i}); in \cref{r1:eq:LD} it enters as $D_\mu=\partial_\mu+qA_\mu J(\hat n)$.
\end{theorem}

\begin{proof}
(a) is \cref{r4:prop:identity} with $\gamma^0_{\hat n}=\hat n$, and $TJ(\hat n)=-\hat n$. (b) $T$ commutes with $\hat n$ and anticommutes with $s_a$, so it commutes with $J(\hat n)$ and
reverses the momentum term of $H_{\hat n}$; uniqueness: the solutions of $Ye_a=\eta_{aa}e_aY$ form a line (Skolem--Noether). (c) is
\cref{r3:prop:parity}(i) (r4-9).
\end{proof}

Restated and fixed: (a) is \cref{rem:ambient} and \cref{prop:kernels1} and \cref{r2:thm:i}(iii) at every clock, (b) is \cref{en:parity} and
\cref{r3:prop:parity} together with \cref{thm:PT} at every clock, (c) is \cref{r3:prop:odd}(i); proved together here (r4-9). The three are
compatible: (b) and (a) say that the mass is parity-even, (b) and (c) that the parity lift acts on the field octants as minus the Hodge star, which is
\cref{r3:prop:parity}. Role (a) belongs to the particle reading and uses $T$ as the product $\hat nJ(\hat n)$; (c) belongs to the field reading and uses $T$
alone; (b) belongs to both, since $\Ad_T|_V=\eta$ is the parity of the frame $(\bB;\bE)$ and of the three lengths, and $T$ is the Dirac parity.

\begin{proposition}[$J(\hat n)$ relative to the field structure]
\label{r4:prop:Jdec}
With $L(\hat n)$ and $A(\hat n)$ as in the conventions:
\begin{enumerate}[label=(\roman*),leftmargin=2.2em]
\item $J(\hat n)=-L(\hat n)+A(\hat n)$, with $L(\hat n)A(\hat n)=A(\hat n)L(\hat n)=0$; for a clock, $-L(\hat n)^2$ and $-A(\hat n)^2$ are complementary projectors of rank four, $L(\hat n)^3=-L(\hat n)$ and
$A(\hat n)^3=-A(\hat n)$.
\item $L(\hat n)=\rho_{\hat n}(L(\hat n))-\tfrac12J(\hat n)$, in which $\rho_{\hat n}(L(\hat n))=\tfrac12\sum_in_is_js_k$ is the Dirac spin rotation of $V^-$ about the spatial axis
$\sum_in_is_i$ paired with $\hat n$ (it acts trivially on $V^+$); and $\rho_{\rm oct}(L_i)=\rho_{\hat n}(L_i)-\tfrac12J(e_i)$ for
every $i$ and every clock. Hence $e^{2\pi L(\hat n)}=e^{2\pi\rho_{\hat n}(L(\hat n))}e^{-\pi J(\hat n)}=(-\id)(-\id)=\id$.
\item For every $\hat n$,
\[
J(\hat n)\Phi_{\rm oct}(X)=-\Phi_{\rm oct}(\hat n\times\bB;\,\hat n\times\bE)+(\bE\cdot \hat n)\,|\mathrm{full}\rangle-(\bB\cdot \hat n)\,|\mathrm{vac}\rangle ,
\]
and $J(\hat n)|\mathrm{vac}\rangle=g(\hat n)$, $J(\hat n)|\mathrm{full}\rangle=-f(\hat n)$. The kernel of $L(\hat n)$ is spanned by the two poles and the longitudinal octants $f(\hat n)$, $g(\hat n)$;
on it $J(\hat n)=A(\hat n)$. On the transverse octants $f(u),g(u)$, $u\perp \hat n$, $J(\hat n)=-L(\hat n)$, the rotation about $\hat n$: $L(\hat n)f(u)=f(\hat n\times u)$, $L(\hat n)g(u)=g(\hat n\times u)$.
\item $J(\hat n)$ does not commute with $\rho_{\rm oct}$ and does not preserve the field octants $\Phi_{\rm oct}(V)$.
\end{enumerate}
\end{proposition}

\begin{proof}
(i) With $t(\hat n)=\sum_in_ie_je_k=J(\hat n)$ and $s(\hat n)=\sum_in_is_js_k$: $L(\hat n)=\tfrac12(s(\hat n)-t(\hat n))$, $A(\hat n)=\tfrac12(s(\hat n)+t(\hat n))$, $s(\hat n)$ and $t(\hat n)$ commute and
$s(\hat n)^2=t(\hat n)^2=-|\hat n|^2\id$, so $L(\hat n)A(\hat n)=\tfrac14(s^2-t^2)=0$ and $-L(\hat n)^2=\tfrac12(\id+st)$, $-A(\hat n)^2=\tfrac12(\id-st)$ with $(st)^2=\id$. (ii) $\rho_{\hat n}(L(\hat n))=\tfrac12s(\hat n)$,
and $L_i=\tfrac12s_js_k-\tfrac12e_je_k$ with $e_je_k=J(e_i)$; $\tfrac12s(\hat n)$ and $\tfrac12J(\hat n)$ commute and square to $-\tfrac14\id$. (iii) $J(\hat n)=\sum_in_ie_je_k$ and $e_I|\mathrm{vac}\rangle=a^+_I|\mathrm{vac}\rangle$
(\cref{r3:prop:fock}) give $J(\hat n)|\mathrm{vac}\rangle=g(\hat n)$; $L(\hat n)$ acts on the $f$'s and $g$'s as the vector rotation (\cref{r3:lem:rot}) and kills the poles; $A(\hat n)f(u)=(\hat n\cdot u)|\mathrm{full}\rangle$
and $A(\hat n)g(u)=-(\hat n\cdot u)|\mathrm{vac}\rangle$. (iv) follows from (iii), since $\rho_{\rm oct}$ preserves $\Phi_{\rm oct}(V)$ (r4-10, r4-11, symbolic in $\hat n$, $\bE$, $\bB$, $u$).
\end{proof}

Fixed. What this means is the following, and the reading is interpretive. The Dirac charge rotation $e^{\theta J(\hat n)}$, read on the octant module, is not
a transformation of the field: on transverse fields it is the rotation by $-\theta$ about the axis paired with the clock, and on the four-dimensional
complement it exchanges the longitudinal components with the pole components, $\bE\cdot \hat n$ with the slot of $\nabla\cdot\bB$ and $\bB\cdot \hat n$ with the slot of
$\nabla\cdot\bE$ (\cref{r3:thm:maxwell}). So the complex structure of the particle reading mixes the field with the Gauss constraints, and it is not in
the commutant of the field reading's Lorentz representation. Item (ii) shows where integer and half-integer spin part on $S$: the field rotation
$L(\hat n)$, which turns $\bB$ about $\hat n$ and $\bE$ about $\sum_in_is_i$, is the Dirac rotation about $\sum_in_is_i$ combined with half a turn of the charge,
and the two signs $-1$ cancel. $J(\hat n)$ must be kept apart from $G_{\rm EM}=\star$: on the octants $\star$ is $-T$ (\cref{r3:prop:odd}), $J(\hat n)=\hat nT$ differs from it by
the clock, and $J(\hat n)$ commutes with $\lint(\hat n)$ while $G$ commutes with $\lG$.

% --- source: two_readings/r4_bridge#4 (Two Lorentz representations on one module, and their couplin)
\subsection{Two Lorentz representations on one module, and their couplings}
\label{r4:sec:reps}

\begin{proposition}[Commutants]
\label{r4:prop:comm}
\begin{enumerate}[label=(\roman*),leftmargin=2.2em]
\item The commutant of $\rho_{\hat n}(\lG)$, the spin lift of $\lint(\hat n)$, in $\End_\R(S)$ is $\operatorname{span}\{\id,u_1,u_2,J(\hat n)\}\otimes\operatorname{span}\{\id,\Omega_{\hat n}\}\cong M_2(\C)$,
of dimension $8$.
\item The commutant of $\rho_{\rm oct}(\lG)$ is $\operatorname{span}\{\Pi_F,T\Pi_F\}\oplus\End(\operatorname{span}\{|\mathrm{vac}\rangle,|\mathrm{full}\rangle\})$, of dimension $6$, with $\Pi_F$
the projector onto the field octants.
\item The joint commutant of $\rho_{\hat n}(\lG)$ and $\rho_{\rm oct}(\lG)$ is $\R\id$, for every clock.
\end{enumerate}
\end{proposition}

\begin{proof}
(i) $u_a$ anticommutes with each vector of $E_{\hat n}$ and so commutes with its bivectors; $\Omega_{\hat n}$ commutes with the even elements of $\Cl(E_{\hat n})$; the eight
products are independent, and the commutant has dimension $8$ because $S$ complexifies to $2(\tfrac12,0)\oplus2(0,\tfrac12)$, whose commutant is
$M_2(\C)\oplus M_2(\C)$ \cite{GoodmanWallach}. (ii) $\Phi_{\rm oct}(V)$ is irreducible of complex type with commutant $\operatorname{span}\{\id,T\}$ (\cref{r3:prop:odd}(iii)), the poles are two
trivial summands, and there is no map between them. (iii) Exact computation at $\hat n=e_1,(2,3,6)/7,(1,2,2)/3$ (r4-12); structurally, $\rho_{\rm oct}(L_i)-\rho_{\hat n}(L_i)=-\tfrac12J(e_i)$
spans $\so3_+$, which with $\lint(\hat n)$ generates $\so{3,3}$, whose commutant $\operatorname{span}\{\id,\Gamma\}$ meets that of the odd $Q_i$ in $\R\id$.
\end{proof}

Fixed; (i) contains the commutant $\HH$ of the four $\gamma^\mu_{\hat n}$ of \cref{r2:thm:i}(ii). No complex structure, and no operator other than a scalar,
commutes with both Lorentz actions.

\begin{proposition}[Three inequivalent representations on $S$]
\label{r4:prop:ineq}
On the same module $S$ the Lie algebra $\lG\cong\so{1,3}$ acts in three ways: through $\rho_{\hat n}$ (complexification $2(\tfrac12,0)\oplus2(0,\tfrac12)$, two
Dirac spinors), through $\rho_{\rm oct}$ ($(1,0)\oplus(0,1)\oplus2(0,0)$, the field strength and two singlets) and through the spin lift of $\lG$ itself
($2(\tfrac12,\tfrac12)$). The Casimirs $\sum_k(\rho(K_k)^2-\rho(L_k)^2)$ and $(\sum_k\rho(K_k)\rho(L_k))^2$ are $\tfrac32$ and $-\tfrac9{16}$ for $\rho_{\hat n}$; $4$ and $-4$ on the field
octants and $0$ on the poles for $\rho_{\rm oct}$; $3$ and $0$ for the spin lift of $\lG$. In $\rho_{\hat n}$, $\rho_{\hat n}(L_3)^2=-\tfrac14\id$ and $e^{2\pi\rho_{\hat n}(L_3)}=-\id$; in the other two
$\rho(L_3)^3=-\rho(L_3)$ and $e^{2\pi\rho(L_3)}=\id$. The three are pairwise inequivalent, and no invertible $U$ satisfies $U\rho_{\hat n}(Y)U^{-1}=\rho_{\rm oct}(Y)$
for all $Y$.
\end{proposition}

\begin{proof}
Equivalent representations have conjugate Casimirs; the values are computed exactly (r4-13). The decompositions are \cref{prop:halfspinors} and \cref{r3:prop:odd}.
\end{proof}

Restated (the decompositions) and fixed (the Casimirs and the inequivalence). That $\rho_{\rm oct}$ does not commute with the spin lift of
$\lint(e_1)$ (r4-13) is the weakest form of this; the proposition says that the two are not even conjugate. They therefore cannot both be the Lorentz
action of one field on $S$. A theory containing a Dirac field and the field octants must treat them as two fields, $\psi$ in a copy of $S$ with
$\rho_{\hat n}$ and $\Phi_{\rm oct}(X)$ in another with $\rho_{\rm oct}$, and the joint action of $Y\in\lint(\hat n)$ is its spin lift $\rho_{\hat n}(\phi_{\hat n}^{-1}Y)$ on $\psi$ and
$\rho_{\rm oct}(\phi_{\hat n}^{-1}Y)$ on the octants; equivalently $Y'\in\lG$ acts by $\rho_{\hat n}(Y')$ and by $\rho_{\rm oct}(Y')$.
Since $\Phi$ intertwines $\lG$ on $V$ with $\rho_{\rm oct}$ (\cref{r3:prop:odd}(iii)), a coupling to $\Phi_{\rm oct}(X)$ is a coupling to $X$ with its own action
$Y\cdot X=[Y,X]$, and conversely.

For $X\in V$ let $\ell(X):=\sum_i(B_iL_i+E_iK_i)\in\lG$, an equivariant isomorphism $V\to\lG$ by the brackets and eq.~\cref{eq:emboost} of \cref{thm:emfield}, and
\[
S_{\hat n}(X):=\rho_{\hat n}(\ell(X))=\tfrac12\sum_i\bigl(B_i\,s_js_k+E_i\,s_in\bigr),
\]
the field strength as a bivector of $E_{\hat n}$; under the dictionary of \cref{r2:prop:dict} it is the Clifford image of the bivector
$\sum_{\mu<\nu}(F_X)_{\mu\nu}\,\gamma_\mu\wedge\gamma_\nu=\sum_i(E_i\,s_i\wedge \hat n+B_i\,s_j\wedge s_k)$, $\gamma_\mu:=\eta_{\mu\nu}\gamma^\nu_{\hat n}$, while the same sum
with $F_X^{\mu\nu}$ in place of $(F_X)_{\mu\nu}$ is $S_{\hat n}(\eta X)$ (r4-16).

\begin{theorem}[Covariant couplings of the Dirac field to the field octants: the derivative-free, field-strength-linear class]
\label{r4:thm:couplings}
Let $Y\in\lG$ act by $\rho_{\hat n}(Y)$ on $\psi$ and by $\rho_{\rm oct}(Y)$ on $\Phi_{\rm oct}(X)$, that is, by $Y\cdot X=[Y,X]$ on $X$. The class classified is that of
couplings linear in the field strength $X$, bilinear in $\psi$, without derivatives and invariant under $\lG$; minimal coupling to a potential
$A_\mu$ through a covariant derivative, couplings with derivatives and couplings nonlinear in $X$ lie outside it (X-07).
\begin{enumerate}[label=(\roman*),leftmargin=2.2em]
\item No nonzero bilinear form $\psi^{\!\top}M\Phi_{\rm oct}(X)$ is invariant, nor one on all of $S$ with $\rho_{\rm oct}$.
\item The invariant trilinear couplings $\psi^{\!\top}M(X)\psi$, $M:V\to\End(S)$ linear with $\rho_{\hat n}(Y)^{\!\top}M(X)+M(X)\rho_{\hat n}(Y)+M(Y\cdot X)=0$, are exactly
$M(X)=\hat nC\,S_{\hat n}(X)$ with $C$ in the commutant of \cref{r4:prop:comm}(i): an eight-dimensional space. For a real Grassmann $\psi$ only antisymmetric
kernels contribute, and these form the two-dimensional space spanned by
\[
M_{\hat n}(X)=T\,J(\hat n)\,S_{\hat n}(X)=-\hat n\,S_{\hat n}(X),\qquad M^\star_{\hat n}(X)=T\,\Gamma\,S_{\hat n}(X)=M_{\hat n}(GX),
\]
and $T\Gamma S_{\hat n}(X)=TJ(\hat n)\Omega_{\hat n}S_{\hat n}(X)$.
For a commuting $\psi$ the symmetric kernels form the six-dimensional space given by
$C\in\operatorname{span}\{u_1,u_2,J(\hat n)\}\otimes\operatorname{span}\{\id,\Omega_{\hat n}\}$. These dimensions hold at every clock.
\item Explicitly $M_{\hat n}(X)=\tfrac12\sum_aE_as_a-\tfrac12\hat n\sum_iB_is_js_k$ and $M^\star_{\hat n}(X)=\tfrac12\sum_aB_as_a+\tfrac12\hat n\sum_iE_is_js_k$. Both commute with $J(\hat n)$ (they carry no
charge); under the parity $\psi\mapsto \hat n\psi$ (or $\psi\mapsto T\psi$), $X\mapsto\eta X$, $M_{\hat n}$ is even, like the mass form $T$, and $M^\star_{\hat n}$ is odd.
\end{enumerate}
\end{theorem}

\begin{proof}
(i) The rotation $e^{2\pi\rho_{\hat n}(L_3)}$ is $-\id$ on $\psi$ and $e^{2\pi\rho_{\rm oct}(L_3)}=\id$ on the octants (\cref{r4:prop:ineq}), so an invariant form satisfies
$\beta(-\psi,\chi)=\beta(\psi,\chi)$; the exact solution spaces are zero at $\hat n=e_1,(2,3,6)/7$ (r4-14). (ii) For every $Y$, $\rho_{\hat n}(Y)^{\!\top}=-\hat n\rho_{\hat n}(Y)\hat n$: the rotations
$\tfrac12s_js_k$ are antisymmetric and commute with $\hat n$, the boosts $\tfrac12s_in$ are symmetric and anticommute with it. Hence, for $C$ commuting with $\rho_{\hat n}$,
\[
\rho_{\hat n}(Y)^{\!\top}\hat nCS_{\hat n}(X)+\hat nCS_{\hat n}(X)\rho_{\hat n}(Y)=\hat nC\bigl(S_{\hat n}(X)\rho_{\hat n}(Y)-\rho_{\hat n}(Y)S_{\hat n}(X)\bigr)=-\hat nC\,S_{\hat n}(Y\cdot X),
\]
using $[\rho_{\hat n}(Y),S_{\hat n}(X)]=\rho_{\hat n}([Y,\ell(X)])=\rho_{\hat n}(\ell(Y\cdot X))$. The map $C\mapsto \hat nCS_{\hat n}$ is injective, and the space of invariant maps $V\to S^*\otimes S^*$ has
dimension $8$: complexified, $V=(1,0)\oplus(0,1)$, $\Lambda^2S=(1,0)\oplus(0,1)\oplus6(0,0)\oplus4(\tfrac12,\tfrac12)$ and
$\mathrm{Sym}^2S=3(1,0)\oplus3(0,1)\oplus2(0,0)\oplus4(\tfrac12,\tfrac12)$, so the dimensions are $2$ and $6$ \cite{GoodmanWallach}. The symmetry of $\hat nCS$ follows from the transpose
rule and $C^{\!\top}$: $(\hat nCS)^{\!\top}=-\hat nS\hat nC^{\!\top}\hat n$, which is $-\hat nCS$ for $C=\id,\Omega_{\hat n}$ and $+\hat nCS$ for $C=u_a,J(\hat n)$ and their products with $\Omega_{\hat n}$. Finally
$TJ(\hat n)=-\hat n$, $TJ(\hat n)\Omega_{\hat n}=T\Gamma$, $T\Gamma=-e_4e_5e_6$, and $M_{\hat n}(GX)=M^\star_{\hat n}(X)$ by (iii). The exact row reduction of the $2304\times168$ and $2304\times216$ linear
systems at $\hat n=e_1$ and $\hat n=(2,3,6)/7$ gives $2$ and $6$, and the displayed kernels span them (r4-15, r4-16). By \cref{r4:prop:coupclock}(i) the dimensions are
the same at every clock. (iii) $\hat nS_{\hat n}(X)\hat n$ reverses the electric part; $\hat n\,e_4e_5e_6\,\hat n=-e_4e_5e_6$ (r4-17).
\end{proof}

Fixed: the dimensions, the explicit kernels and their properties. The textbook
comparison is inherited. In complex Dirac theory the Lorentz- and gauge-invariant couplings of $\psi$ to the field strength that are bilinear in
$\psi$ and contain no derivative are the Pauli term $\bar\psi\sigma^{\mu\nu}\psi F_{\mu\nu}$, the coupling of an anomalous magnetic moment introduced by
Pauli \cite{Pauli41}, and the electric-dipole term $i\bar\psi\sigma^{\mu\nu}\gamma_5\psi F_{\mu\nu}$, odd under parity and time reversal
\cite{PeskinSchroeder,ItzyksonZuber}. With $J(\hat n)$ in the place of $i$, $\bar\psi\cdots\psi\leftrightarrow\tfrac12\psi^{\!\top}T\cdots\psi$ as in \cref{r1:eq:LD}, and the real
form $\Omega_{\hat n}$ of $i\gamma_5$ (\cref{prop:massangle}), $\sigma^{\mu\nu}F_{\mu\nu}$ becomes a multiple of $J(\hat n)S_{\hat n}(X)$: $M_{\hat n}$ is the real form of the Pauli term
and $M^\star_{\hat n}$ that of the electric-dipole term, up to normalisation and the sign convention for $\gamma^a$; their parities in (iii) agree with the
textbook ones. The Lagrangian density
\begin{equation}
\label{r4:eq:LP}
\mathcal L_{\rm P}=\tfrac12\,\psi^{\!\top}\Bigl(\tfrac{\mu}{c}\,M_{\hat n}(X_{\rm EM})+d\,M^\star_{\hat n}(X_{\rm EM})\Bigr)\psi,\qquad X_{\rm EM}=(c\bB;\bE),
\end{equation}
has dimension $\mathrm{J\,m^{-3}}$ in SI units for $[\mu]=\mathrm{J\,T^{-1}}$ and $[d]=\mathrm{C\,m}$, and its two couplings have mass dimension $-1$ in natural units: they
are not renormalisable \cite{PeskinSchroeder}. The arena fixes neither constant (open). The kernel $M_{\hat n}$ is written through $T$ and $J(\hat n)$, as (E3) proposed:
$T$ is the adjoint form, $J(\hat n)$ the imaginary unit of $\sigma^{\mu\nu}$, and $S_{\hat n}(X)$ the field bivector; the dual uses the other $\mathrm{Spin}(3,3)$-invariant
form $T\Gamma$ (\cref{r3:prop:energies}(iii)) and no $J(\hat n)$. With these couplings the Dirac sector reads, at every clock,
\[
\mathcal L_D+\mathcal L_{\rm P}=\tfrac12\,\psi^{\!\top}\,T\,\Bigl[J(\hat n)\Bigl(\gamma^\mu_{\hat n}D_\mu+\tfrac{\mu}{c}S_{\hat n}(X_{\rm EM})+d\,\Omega_{\hat n}S_{\hat n}(X_{\rm EM})\Bigr)-m\Bigr]\psi,
\]
in the units of \cref{r1:eq:LD}: $T$ is the common left factor of the kinetic, mass and Pauli kernels.

\begin{proposition}[The couplings when the clock moves or the axes are permuted]
\label{r4:prop:coupclock}
The couplings of \cref{r4:thm:couplings} transform as follows.
\begin{enumerate}[label=(\roman*),leftmargin=2.2em]
\item Let $q\in\mathrm{Spin}(V^+)$ be a rotor with $qe_1q^{-1}=\hat n$. Then, with $X$ fixed,
\[
M_{\hat n}(X)=qM_{e_1}(X)q^{-1},\qquad M^\star_{\hat n}(X)=qM^\star_{e_1}(X)q^{-1},\qquad qTq^{-1}=T,\qquad qT\Gamma q^{-1}=T\Gamma :
\]
under $\psi\mapsto q\psi$ the couplings at $e_1$ become those at $\hat n$. The factors $J(\hat n)$, $\Omega_{\hat n}$, $\gamma^\mu_{\hat n}$, $S_{\hat n}$ are conjugated, and $T$, the
clock-independent part of every kernel in the display above, is fixed. The couplings themselves depend on the clock: $M_{e_1}\neq M_{e_2}$.
\item For a diagonal rotation $R$, $M_{R\hat n}(RX)=RM_{\hat n}(X)R^{-1}$ and likewise for $M^\star$: clock, spinor and field turn together.
\item For an odd temporal permutation alone, $A_\sigma(M_{\hat n}(X))=M_{\sigma \hat n}(X)$ and $A_\sigma(M^\star_{\hat n}(X))=M^\star_{\sigma \hat n}(X)$, while $A_\sigma(T)=-T$. For the paired
permutation, $A'_\sigma(M_{\hat n}(X))=M_{\sigma \hat n}(\tilde X)$ and $A'_\sigma(M^\star_{\hat n}(X))=\operatorname{sgn}(\sigma)M^\star_{\sigma \hat n}(\tilde X)$, with $\tilde X$ of \cref{r4:prop:perm}(ii).
\end{enumerate}
\end{proposition}

\begin{proof}
(i) $q$ is orthogonal and conjugates $\rho_{e_1}$ to $\rho_{\hat n}$ (\cref{r4:prop:rot}(iii)); $S_{\hat n}$, $J(\hat n)$, $\Omega_{\hat n}$ follow, and $T$, $\Gamma$ are fixed. (ii) The same with the
spin lift of a diagonal rotation, which also rotates $s_a$ and $X$. (iii) $A(TJ(\hat n))=TJ(\sigma \hat n)$ because the two signs cancel; $A$ fixes $s_a$ and maps $\hat n$ to
$\sigma \hat n$, so $A(S_{\hat n}(X))=S_{\sigma \hat n}(X)$; $A(\Gamma)=\operatorname{sgn}(\sigma)\Gamma$ for $A_\sigma$ and $\Gamma$ for $A'_\sigma$, and $A'(S_{\hat n}(X))=S_{\sigma \hat n}(\tilde X)$ because
$s_{\sigma(j)}s_{\sigma(k)}$ is $\operatorname{sgn}(\sigma)$ times the cyclically ordered pair (r4-18, r4-19).
\end{proof}

Fixed. Expectation (E2), that the couplings change sign with $T$ under an odd permutation, is refuted for the Pauli term under both kinds of
permutation and for the dual term under the temporal permutation alone; it holds only for the dual term under the paired permutation:
$M_{\hat n}=TJ(\hat n)S_{\hat n}$ contains the product $TJ(\hat n)=-\hat n$, in which the two reversals cancel. What changes sign is the
Pauli term \emph{relative to the mass term}, since $T$ reverses and $M_{\hat n}$ does not; after the axial rotation that restores the mass
(\cref{r4:cor:senses}), $M_{\hat n}$ is reversed. Physically the paired odd permutation is an improper rotation, and the electric-dipole coupling is the
parity-odd one (\cref{r4:thm:couplings}(iii)).

% --- source: two_readings/r4_bridge#5 (The two first-order operators and the spatial derivatives)
\subsection{The two first-order operators and the spatial derivatives}
\label{r4:sec:operators}

The Dirac operator of the clock $\hat n$ is $\gamma^\mu_{\hat n}\partial_\mu=\hat n\partial_0+\sum_as_a\partial_a$, and $\gamma^0\gamma^\mu\partial_\mu=\partial_0+\sum_a\alpha_a(\hat n)\partial_a$ with
$\alpha_a(\hat n):=ns_a$; the octant Maxwell operator is $\partial_t-D$ with $D=\sum_ae_a\partial_a$ (\cref{r3:thm:maxwell}). In both, $\partial_a=\partial/\partial x^a$ differentiates
along the three spatial coordinates, the lengths on $V^-$, in $\mathrm m^{-1}$.

\begin{proposition}[Two Pauli triples]
\label{r4:prop:triples}
\begin{enumerate}[label=(\roman*),leftmargin=2.2em]
\item $\alpha_1(\hat n),\alpha_2(\hat n),\alpha_3(\hat n)$ are symmetric, square to $\id$ and anticommute; so do $e_1,e_2,e_3$. For any such triple $\pi_a$,
$(\sum_ak_a\pi_a)^2=|k|^2\id$, so both $\partial_0\pm\sum_a\pi_a\partial_a$ factor the wave operator.
\item $J(\hat n)\alpha_a(\hat n)=Ts_a$ and $e_a=-TJ(e_a)$: the Dirac triple is one clock times the spatial generators, $\alpha_a(\hat n)=-TJ(\hat n)s_a$; the octant triple consists of
the three coordinate clocks themselves, $e_a=-TJ(e_a)$, one for each spatial axis.
\item No fixed element $Z$ gives $\alpha_a(\hat n)=Ze_a$ for all $a$, nor $\alpha_a(\hat n)=e_aZ$; in particular none of $T$, $\Gamma$, $\omega_i$, $J(\hat n)$ does.
\end{enumerate}
\end{proposition}

\begin{proof}
(i) $\hat n$ and $s_a$ are orthogonal matrices, $\hat n$ symmetric, $s_a$ antisymmetric, anticommuting: $(ns_a)^2=-\hat n^2s_a^2=\id$ and $ns_ans_b=-s_as_b$. (ii)
$J(\hat n)\hat n=T$, and $TJ(e_a)=T^2e_a$. (iii) $\alpha_a(\hat n)=Ze_a$ for all $a$ would give $Z=ns_ae_a$ for $a=1,2$, hence $s_1e_1=s_2e_2$, which is false;
$\alpha_a(\hat n)=e_aZ$ would give $e_1ns_1=e_2ns_2$, whose grade-one part $\hat n_1s_1=\hat n_2s_2$ forces $\hat n=\pm e_3$, and $e_1e_3s_1\neq e_2e_3s_2$ (r4-20, symbolic
in $\hat n$).
\end{proof}

\begin{theorem}[The triples are conjugate only by elements of mixed parity]
\label{r4:thm:conj}
For each clock the $U\in M_8(\R)$ with $U\alpha_a(\hat n)=e_aU$ $(a=1,2,3)$ form an eight-dimensional space, $\mathcal C\,U_0$ for any invertible
solution $U_0$, where $\mathcal C=\operatorname{span}\{\id,T\}\otimes\Cl^0(V^-)\cong M_2(\C)$ is the commutant of the $e_a$; its invertible elements are dense. No nonzero solution is even and none is odd, so no element of the Clifford group, of
$\mathrm{Spin}(3,3)$ or of $\so{3,3}$ relates the two triples. Every invertible solution carries the Dirac representation to the left-multiplication
representation of the temporal Pauli algebra,
\[
U\rho_{\hat n}(L_i)U^{-1}=-\tfrac12e_je_k=-\tfrac12J(e_i),\qquad U\rho_{\hat n}(K_i)U^{-1}=-\tfrac12e_i ,
\]
whereas $\rho_{\rm oct}(L_i)=-\tfrac12e_je_k+\tfrac12s_js_k$ and $\rho_{\rm oct}(K_i)=-\tfrac12e_i-\tfrac12s_i\omega$; by \cref{r4:prop:ineq} no $U$ relates $\rho_{\hat n}$ to $\rho_{\rm oct}$.
\end{theorem}

\begin{proof}
Both triples generate copies of $\Cl(3,0)\cong M_2(\C)$, which has a unique simple module, $\R^4$; $S$ is twice it for each, so intertwiners exist and form
a torsor of the commutant (inherited: Wedderburn's theorem). If $U$ had definite parity, $U\alpha_a$ would have the parity of $U$ and $e_aU$ the
opposite one, so both would vanish, and $U=0$ since $e_a$ is invertible. The exact solution
spaces, their even and odd parts, invertibility of a generic solution and the conjugated generators are computed at $\hat n=e_1$ and $(2,3,6)/7$ (r4-21);
$\rho_{\hat n}(L_i)=-\tfrac12\alpha_j\alpha_k$ and $\rho_{\hat n}(K_i)=-\tfrac12\alpha_i$.
\end{proof}

\begin{proposition}[Left and right in the Fock picture]
\label{r4:prop:fockLR}
On $S=\Cl(V^+)|\mathrm{vac}\rangle$ (\cref{r3:prop:fock}) the temporal generators act by left multiplication, $e_a(x|\mathrm{vac}\rangle)=(e_ax)|\mathrm{vac}\rangle$, and the spatial ones by twisted
right multiplication, $s_a(x|\mathrm{vac}\rangle)=(\hat x e_a)|\mathrm{vac}\rangle$, $\hat x=(-1)^{|x|}x$. Hence $\alpha_a(\hat n)(x|\mathrm{vac}\rangle)=(\hat n\hat xe_a)|\mathrm{vac}\rangle$.
\end{proposition}

\begin{proof}
On the eight blades $x=e_I$, $I\subset\{1,2,3\}$ (r4-22).
\end{proof}

Restated: the right action is that of \cref{r3:prop:odd}(ii), where $e_{a+3}\omega$ acts as $x\mapsto-xe_a$, on the Fock model of Chevalley
\cite{Chevalley54} (\cref{r3:prop:fock}); verified here on all blades. Inherited: in the Pauli and space-time algebras a Dirac spinor transforms
by one-sided and the field bivector by two-sided multiplication \cite{Hestenes66,Hestenes67}, which is the difference between $\rho_{\hat n}$ and $\rho_{\rm oct}$
in \cref{r4:thm:conj}.

These results answer the remark after \cref{r3:thm:maxwell}. The Dirac reading couples $\partial_a$ to the spatial generators, through $\gamma^0\gamma^a=\alpha_a(\hat n)=ns_a$;
in the Fock picture this is the clock $\hat n$ on the left and the temporal Pauli generator $e_a$ on the right. The octant field reading couples $\partial_a$ to
the temporal generators $e_a$ on the left. In both the index $a$ of the derivative is matched to the index $a$ of a generator by the pairing
$a\leftrightarrow a+3$ of the arena, which is background structure, and the derivative is along the length $x^a$ on $V^-$ in both. The two triples are
conjugate (fixed), by elements of mixed parity only, which are no symmetry of the arena (fixed), and any Pauli triple would make the operator
hyperbolic (fixed). So neither operator shows that spatial differentiation follows from temporal directions: the octant operator uses the three
coordinate clocks $e_a$ as spatial Pauli matrices, one per axis, not one clock, and the statement that space is differentiated ``through time'' is a
reading of the choice of representation (interpretive), not a consequence of either equation; as an inference it is refuted, since conjugating by $U$
rewrites Maxwell's equations with the Dirac triple for the octant map $U^{-1}\Phi$.

% --- source: two_readings/r4_bridge#6 (Varying and permuting the clock)
\subsection{Varying and permuting the clock}
\label{r4:sec:clock}

\begin{table}[t]
\centering
\footnotesize
\renewcommand{\arraystretch}{1.15}
\begin{tabular}{@{}>{\raggedright\arraybackslash}p{7.3cm}>{\raggedright\arraybackslash}p{7.3cm}@{}}
\toprule
depends on the clock $\hat n$ & independent of the clock\\
\midrule
$\gamma^0_{\hat n}=\hat n$, $J(\hat n)$, $\Omega_{\hat n}$; $\lint(\hat n)$, $\phi_{\hat n}$, the boosts $\rho_{\hat n}(K_i)$; $\alpha_a(\hat n)$; $K^\mu_{\hat n}$, $H_{\hat n}(\bp)$; the Dirac complex structure $\pm J(\hat n)$ and charge;
the eigenspaces $\hat n=\pm1$ that carry the two senses; the couplings $M_{\hat n}$, $M^\star_{\hat n}$; the base time $t_{\hat n}$ &
$T$, $T\Gamma=-e_4e_5e_6=-\hat n\Omega_{\hat n}$, $\Gamma=J(\hat n)\Omega_{\hat n}$; the rest flow $e^{-m\lambda T}$; the rotations $\rho_{\hat n}(L_i)=\tfrac12s_js_k$; $\rho_{\rm oct}$, $\Phi$, $G$, $\lG$, the octants and
pole operators; $D=\sum_ae_a\partial_a$ and the coefficients of the octant Maxwell equation \cref{r3:eq:maxwell}\\
\bottomrule
\end{tabular}
\caption{Clock dependence of the objects of the two readings (r4-2, r4-3, r4-18, r4-23).}
\label{r4:tab:clock}
\end{table}

\begin{proposition}[The field reading does not see the clock; the particle reading does]
\label{r4:prop:clockdep}
The objects of the right column of \cref{r4:tab:clock} are defined without $\hat n$; the octant Maxwell equation at clock $\hat n$ is
$(\partial_{t_{\hat n}}-D)\Phi_{\rm oct}(X)=-(\rho|\mathrm{vac}\rangle+\bj\cdot f)$, with the same matrices for every clock, the clock entering only through the time coordinate of the
base $E_{\hat n}$. The objects of the left column depend on $\hat n$: $\hat n\mapsto J(\hat n)$ is injective, $\lint(\hat n)=\lint(\hat n')$ only for $\hat n'=\pm \hat n$, and $M_{e_1}\neq M_{e_2}$. All of
them are carried from one clock to another by the rotors of $\mathrm{Spin}(V^+)$, which are not symmetries of the field reading, and, together with the field,
by the diagonal rotations, which are.
\end{proposition}

\begin{proof}
By inspection of the definitions in \cref{r3:sec:main}; \cref{r4:prop:identity,r4:prop:rot,r4:prop:coupclock}; $\lint(\hat n)$ determines $E_{\hat n}$ and hence $\pm \hat n$ (\cref{r2:prop:planes}).
\end{proof}

Fixed. This answers the question whether different clocks are physically different sectors, gauge-equivalent
choices or an order parameter: the diagonal rotations map the coupled data at $(\hat n,\psi,X)$ to those at $(R\hat n,R\psi,RX)$ (\cref{r4:prop:coupclock}(ii)), so for
theories built from the structures of this chapter a change of clock is a change of frame (fixed as a statement of equivalence; its reading as ``no
sector'' is interpretive); whether a dynamics singles out a clock is open, as in \cref{r2:rem:projection}. Equivalence under a global symmetry is
not gauge redundancy: symmetry-related configurations can be distinct physical states, and a gauge reading would need an action, a
redundancy of its solutions and gauge-invariant observables. A position-dependent clock rotor adds $R^{-1}\partial_\mu R$ to the Dirac operator and
so needs a connection that is not constructed here (X-02, TIME-08).

\begin{theorem}[Uniqueness of $T$]
\label{r4:thm:unique}
Let $Y\in\Cl(3,3)$.
\begin{enumerate}[label=(\roman*),leftmargin=2.2em]
\item If $\hat n^{-1}Y$ commutes with $\gamma^0_{\hat n},\dots,\gamma^3_{\hat n}$ for three linearly independent clocks, then $Y\in\R T$. Three clocks of the Fisher--Rao octant suffice,
e.g.\ $e_1$, $(1,2,2)/3$, $(2,3,6)/7$. For two clocks $\hat n,\hat n'$ the solutions are $\operatorname{span}\{T,\,(\hat n\times \hat n')\,e_4e_5e_6\}$.
\item If $\hat n^{-1}Y$ lies, for two independent clocks, in the joint commutant $\operatorname{span}\{\id,\Omega_{\hat n},J(\hat n),\Omega_{\hat n}J(\hat n)\}$ of $\lint(\hat n)$ and $J(\hat n)$ in which \cref{thm:i} selects the imaginary unit, then $Y\in\operatorname{span}\{T,e_4e_5e_6\}=\operatorname{span}\{T,T\Gamma\}$, the two $\mathrm{Spin}(3,3)$-invariant forms. Indeed $\hat nJ(\hat n)=T$ and
$\hat n\Omega_{\hat n}=e_4e_5e_6$ for every clock.
\item Consequently $T$ is, up to scale, the only element of $\Cl(3,3)$ that is the product of the clock with a Dirac complex structure (a square root of $-\id$ in
the commutant of the $\gamma^\mu_{\hat n}$) for every clock and that does not depend on the clock: the products are $\pm \hat nJ(\hat n)=\pm T$.
\end{enumerate}
\end{theorem}

\begin{proof}
(i) The commutant of the $\gamma^\mu_{\hat n}$ is $\operatorname{span}\{\id,J(\hat n),u_1\Omega_{\hat n},u_2\Omega_{\hat n}\}$ (\cref{r2:thm:i}(ii)), so $Y\in \hat n\cdot(\ldots)=\operatorname{span}\{n,T,u_1e_4e_5e_6,u_2e_4e_5e_6\}$, using
$nu_a\Omega_{\hat n}=-u_ae_4e_5e_6$. The four have grades $1,3,4,4$, so for $Y$ fixed the grade-one part lies in $\bigcap\R \hat n=0$, the grade-four part in $(\bigcap U_{\hat n})\,e_4e_5e_6$,
which is $0$ for three independent clocks and the line of $\hat n\times \hat n'$ for two. (ii) $\hat n\cdot\operatorname{span}\{\id,\Omega_{\hat n},J(\hat n),\Omega_{\hat n}J(\hat n)\}=\operatorname{span}\{n,e_4e_5e_6,T,\hat n\Gamma\}$, of grades
$1,3,3,5$; the grades one and five vanish for two independent clocks; $T\Gamma=-e_4e_5e_6$ (r4-1) and $\operatorname{span}\{T,T\Gamma\}$ is the space of invariant forms
(\cref{r3:prop:energies}(iii)). (iii) If $\hat n^{-1}Y$ is a square root of $-\id$ commuting with the $\gamma^\mu_{\hat n}$ for every clock, then $Y\in\R T$ by (i), so
$\hat n^{-1}Y=aJ(\hat n)$, and $(aJ(\hat n))^2=-a^2\id$ gives $a=\pm1$; with the selection of \cref{thm:i} at every clock (\cref{r2:thm:i}) the Dirac
complex structures are $\pm J(\hat n)$ and the products are $\pm T$ directly. For two clocks the elements $\cos\vartheta\,T+\sin\vartheta\,\hat w\,e_4e_5e_6$,
$\hat w$ the unit vector along $\hat n\times \hat n'$, meet the conditions at both, so (iii), like (i), needs three clocks. Exact solutions in r4-23.
\end{proof}

Fixed; by the two-clock exception in (i) two clocks do not single out $T$, and three are needed. At a single fixed clock the theorem gives no
uniqueness. The commutant $\operatorname{span}\{\id,J(\hat n),u_1\Omega_{\hat n},u_2\Omega_{\hat n}\}$ is a quaternion algebra, and its square roots of $-\id$ form the
two-sphere $aJ(\hat n)+b\,u_1\Omega_{\hat n}+c\,u_2\Omega_{\hat n}$, $a^2+b^2+c^2=1$; for instance $e_2\Omega_{e_1}$ is a Dirac complex structure at $\hat n=e_1$ other than
$\pm J(e_1)$ (X-01; z1-7c). So ``a square root of $-\id$ in the commutant'' selects $\pm J(\hat n)$ only with an extra datum: pointwise, the
joint commutant with $\lint(\hat n)$ and $J(\hat n)$ of \cref{thm:i} (or commutation with $T$); globally, the two hypotheses of (iii),
that the product with the clock be the same element for three independent clocks. The second element of (ii), $e_4e_5e_6$, is the
clock times the other complex structure $\Omega_{\hat n}$, which \cref{thm:i} excludes because it anticommutes with the $\gamma^\mu_{\hat n}$.

\begin{remark}[The ``covariant mass'' $e_4e_5e_6$]
\label{r4:rem:V}
A possible assignment is that, of the two trivectors anticommuting with the three Dirac matrices
$\alpha_a$, the spatial volume $e_4e_5e_6$ is the scalar mass and $e_1e_2e_3$ its pseudoscalar partner. In the present conventions the first half is
confirmed at every clock: the trivectors anticommuting with $\alpha_1(\hat n),\alpha_2(\hat n),\alpha_3(\hat n)$ are $\operatorname{span}\{T,e_4e_5e_6\}$, the clock-independent space of
\cref{r4:thm:unique}(ii) (r4-24, at four clocks; at all clocks by \cref{r4:prop:rot}(iii), since conjugation by a rotor preserves the grade). The assignment is corrected. As the matrix $\beta$ of $H=\sum_a\alpha_ap_a+m\beta$, $e_4e_5e_6=\hat n\Omega_{\hat n}$ is admissible
(symmetric, square $\id$, commuting with $J(\hat n)$, $H^2=(|\bp|^2+m^2)\id$), but relative to the kinetic kernels with $\gamma^0=\hat n$ it gives the pseudoscalar mass
of \cref{prop:massangle}, whose Grassmann kernel $J(\hat n)e_4e_5e_6=T\Omega_{\hat n}$ depends on the clock; the scalar mass has $\beta=\hat n$ and the
clock-independent kernel $J(\hat n)\hat n=T$. The element $e_1e_2e_3$ is not a $\beta$ at all (it is antisymmetric with square $-\id$); it is the kernel.
\end{remark}

Interpretation. In the particle reading the clock and the complex structure both move, and their product $T$ does not; in the field reading $T$ is
the duality on the octants and, as the Chevalley--Mukai pairing, the form pairing $E_k$ with $B_k$ (\cref{r3:prop:energies}). The mass form is the part of
the particle reading that the field reading shares. This is interpretive; what is fixed is \cref{r4:thm:roles,r4:thm:unique}.

% --- source: two_readings/r4_bridge#7 (Answer)
\subsection{Answer}
\label{r4:sec:answer}

\begin{enumerate}[label=(\arabic*),leftmargin=2.2em]
\item \emph{Consistency.} $T=\hat nJ(\hat n)$ holds at every clock and is compatible with all operations considered: $T$ is invariant under all rotations of the
temporal triple and the diagonal rotations, reversed together with $J(\hat n)$ by odd permutations, and with the duality operator by the paired ones,
and fixed by the
reversal of the clock (\cref{r4:prop:rot,r4:prop:perm}; fixed). Its three roles, mass form, parity and duality, are compatible (\cref{r4:thm:roles};
restated and fixed). It incorporates consistently into both readings.
\item \emph{The Lorentz representations stay apart.} On $S$ the Dirac representation $\rho_{\hat n}$ and the octant representation $\rho_{\rm oct}$ are inequivalent,
their joint commutant is $\R\id$, and no operator, $T$ included, relates them (\cref{r4:prop:comm,r4:prop:ineq,r4:thm:conj}; fixed). A single module
carrying both as the Lorentz action of one field is obstructed; the Dirac field and the field octants are two fields.
\item \emph{The bridge as a coupling.} The answer to ``The main obstruction'' of \cref{r3:sec:answer}: $\rho_{\rm oct}$, $\rho_{\hat n}$ and $\Phi$ coexist as the
actions on two different fields, and the covariant local couplings without derivatives that are linear in the field strength are trilinear,
quadratic in $\psi$: the Pauli kernel $TJ(\hat n)S_{\hat n}(X)$ and the electric-dipole kernel $T\Gamma S_{\hat n}(X)$, a two-dimensional space at every clock,
bilinear couplings being excluded by the $2\pi$ rotation (\cref{r4:thm:couplings}; fixed, the textbook terms inherited). Couplings of higher
order in $X$ exist, for instance $(\psi^{\!\top}T\psi)\,\eta(X,X)$, since $T$ is $\rho_{\hat n}$-invariant and $\eta(X,X)$ is $\lG$-invariant (r4-25). Since $\Phi$ is an equivalence onto its image, the octants add no coupling that $X$ itself
does not have. $T$ is the clock-independent factor of every Dirac kernel; the couplings are conjugated by the rotation that moves the clock, and an odd
permutation paired with the spatial one reverses the electric-dipole term with $T$ but not the Pauli term, and the temporal one alone reverses
neither (\cref{r4:prop:coupclock}; fixed, (E2) refuted for the Pauli term).
\item \emph{$J(\hat n)$ on the field.} $J(\hat n)=-L(\hat n)+A(\hat n)$ rotates the transverse octants about the axis with the components of $\hat n$ and maps the
longitudinal octants to the Gauss-law poles; this field rotation $L(\hat n)$ is the Dirac rotation about the spatial axis $\sum_in_is_i$ combined with
half a turn of the charge (\cref{r4:prop:Jdec}; fixed; the reading interpretive).
\item \emph{Spatial derivatives.} The remark after \cref{r3:thm:maxwell} is answered by \cref{r4:prop:triples,r4:thm:conj,r4:prop:fockLR}: both operators
differentiate along the lengths of $V^-$, the Dirac reading through the spatial generators and one clock, the octant reading through the three temporal
generators; the triples are conjugate only by mixed-parity elements, and nothing derives spatial differentiation from temporal directions (fixed; the
inference refuted).
\item \emph{The clock.} The field reading is clock-independent and the particle reading is not; $T$ is the unique clock-independent product of the
clock with a Dirac complex structure, and with the second complex structure $\Omega_{\hat n}$ the space becomes $\operatorname{span}\{T,T\Gamma\}$ (\cref{r4:prop:clockdep,r4:thm:unique};
fixed). Clocks are related by the diagonal rotations, a symmetry of both readings and of the couplings, so a change of clock is a change of frame
(fixed as an equivalence); whether a dynamics selects a clock is open.
\item \emph{What remains for a covariant Dirac--Maxwell theory in the arena.} The minimal coupling through the potential and $J(\hat n)$, $D_\mu=\partial_\mu+qA_\mu J(\hat n)$
(\cref{r1:prop:total}, with the Lagrangian \cref{r1:eq:Ltot}; inherited, restated), which uses $A_\mu$ on the base and not the octants; the Pauli and electric-dipole couplings
\cref{r4:eq:LP}, with two constants the arena does not fix (open); the absence of a common Lorentz representation on one $S$ (obstructed at one constituent;
the two readings can be housed together at the mixed second rung of \cref{sec:mixed}, at the price stated there, restated); and the
quantisation of either reading (open).
\end{enumerate}

\paragraph{Open items.} Reversing $\hat n$ reverses $J(\hat n)$ (\cref{r4:prop:perm}(iv)); the temporal permutations and rotations are
\cref{r4:prop:rot,r4:prop:perm}; $J(\hat n)$ is not the duality, a common Lorentz representation on $S$ does not exist, and the cross-sector
coupling has the units, covariance and parity stated in \cref{r4:thm:couplings} and \cref{r4:eq:LP} (\cref{r4:prop:Jdec}, \cref{r4:sec:reps});
the clock dependence is \cref{r4:prop:clockdep}. The projector vacuum, the mismatch between the mixed and the spectator embeddings, and the
symplectic and quantum consistency of a cross-sector coupling are not treated here (open).

{\small\renewcommand{\arraystretch}{1.12}
\begin{longtable}{@{}>{\raggedright\arraybackslash}p{0.09\textwidth}>{\raggedright\arraybackslash}p{0.58\textwidth}>{\raggedright\arraybackslash}p{0.27\textwidth}@{}}
\caption{The claims of this section, with their checks and classes.}\label{r4:tab:claims}\\
\toprule
id & claim & class\\
\midrule
\endfirsthead
\toprule
id & claim & class\\
\midrule
\endhead
r4-1 & conventions: Clifford relations, octants, $T$, $\Gamma$, $T\Gamma=-e_4e_5e_6$ & restated\\
r4-2 & $T=\hat nJ(\hat n)=J(\hat n)\hat n$, $J(\hat n)^2=-\id$ at every clock (\cref{r4:prop:identity}) & restated (\cref{r2:prop:J})\\
r4-3, 4 & $T$ invariant, $J(\hat n)$ covariant under $\so3_+$ and $\so3_\Delta$; only $\so3_\Delta$ commutes with $G$ (\cref{r4:prop:rot}) & fixed\\
r4-5, 6 & odd permutations: $T\mapsto-T$, $J(\hat n)\mapsto-J(\sigma \hat n)$; paired: $\Phi\mapsto\Phi_{\rm t}$, duality $-T\mapsto T$ (\cref{r4:prop:perm}) & fixed\\
r4-7 & clock reversal is a half-turn fixing $T$; $C=RT$ (\cref{r4:prop:perm}(iv)) & fixed\\
r4-8 & $T=\pm J(\hat n)$ on $\hat n=\pm1$; odd permutation $\Rightarrow(m,q)\mapsto(-m,-q)$ (\cref{r4:cor:senses}) & restated (\cref{mb:senses}); fixed\\
r4-9 & three roles of $T$; $T$ is the Dirac parity at every clock (\cref{r4:thm:roles}) & restated, fixed\\
r4-10 & $J(\hat n)=-L(\hat n)+A(\hat n)$, $LA=0$; $L(\hat n)=\rho_{\hat n}(L(\hat n))-\tfrac12J(\hat n)$ (\cref{r4:prop:Jdec}) & fixed\\
r4-11 & $J(\hat n)$ rotates transverse octants, maps longitudinal ones to the poles (\cref{r4:prop:Jdec}) & fixed; reading interpretive\\
r4-12 & commutants $M_2(\C)$, $\C\oplus M_2(\R)$, joint $\R\id$ (\cref{r4:prop:comm}) & fixed\\
r4-13 & three inequivalent Lorentz representations on $S$ (\cref{r4:prop:ineq}) & restated (decompositions), fixed\\
r4-14 & no invariant bilinear coupling $\psi^{\!\top}M\Phi_{\rm oct}(X)$ (\cref{r4:thm:couplings}(i)) & obstructed\\
r4-15, 16 & trilinear couplings: dimension $2$ (Grassmann), $6$ (commuting); Pauli $TJ(\hat n)S_{\hat n}$ and dual $T\Gamma S_{\hat n}$ & fixed; textbook terms inherited\\
r4-17 & Pauli term parity-even, dual parity-odd (\cref{r4:thm:couplings}(iii)) & fixed; inherited in substance\\
r4-18 & couplings conjugated by the clock rotor, $T$ fixed; diagonal covariance (\cref{r4:prop:coupclock}) & fixed\\
r4-19 & odd permutation: Pauli term does not change sign, dual does when paired (\cref{r4:prop:coupclock}(iii)) & fixed; (E2) refuted for the Pauli term\\
r4-20 & two Pauli triples, $J(\hat n)\alpha_a=Ts_a$, $e_a=-TJ(e_a)$, no one-sided relation (\cref{r4:prop:triples}) & fixed\\
r4-21 & intertwiners: dimension $8$, none of definite parity; Dirac $\mapsto$ left Pauli representation (\cref{r4:thm:conj}) & fixed; existence inherited\\
r4-22 & temporal $=$ left, spatial $=$ twisted right multiplication (\cref{r4:prop:fockLR}) & restated (\cref{r3:prop:odd}), verified; one-/two-sided inherited\\
r4-23 & uniqueness of $T$; two-clock exception; $\operatorname{span}\{T,T\Gamma\}$ (\cref{r4:thm:unique}) & fixed\\
r4-24 & trivectors anticommuting with $\alpha_a(\hat n)$; $e_4e_5e_6$ as $\beta$ is the pseudoscalar mass (\cref{r4:rem:V}) & fixed; earlier assignment corrected\\
r4-25 & a covariant quartic coupling $(\psi^{\!\top}T\psi)\,\eta(X,X)$: the trilinear couplings are the only ones linear in $X$, not the only ones (\cref{r4:sec:answer}) & fixed\\
-- & spatial derivatives ``from temporal directions'' (\cref{r4:sec:operators}) & interpretive; inference refuted\\
-- & a change of clock is a change of frame (\cref{r4:prop:clockdep}) & fixed (equivalence); interpretive\\
-- & mass form $=$ the shared part of the two readings & interpretive\\
-- & values of $\mu$, $d$; a dynamics selecting the clock; quantisation & open\\
\bottomrule
\end{longtable}}

% --- source: tex/s13_spin#0 ((head))


% --- source: tex/s13_spin#1 (Spin one in both readings)
\section{Spin one in both readings}
\label{sec:spin}

A photon is a massless particle of spin one with two helicities. Its covariant description is a field strength, a two-form, and
its potential is a four-vector defined up to gauge. The arena offers two places for it. In the \emph{interior reading} the
potential is a grade-one element of the four-plane $E$ on which the Dirac sector is built and the field strength is the grade-two
module of \cref{prop:bivectorF}. In the \emph{frame reading} the field strength occupies the six directions themselves, under the
algebra $\lG$ of \cref{thm:GLorentz}. This section shows that each reading produces a free field with the spin content of the photon,
that the two are related by the Klein correspondence, and that at the second rung both can be housed under one Lorentz action and
coupled to a Dirac block. \Cref{fig:tworoutes} summarises the comparison.

% --- source: tex/s13_spin#2 (The univectorial photon)
\subsection{The univectorial photon}

\begin{proposition}[Grade one under the interior action]
\label{prop:gradeone}
Under $\lint\oplus\uu1_J$, with $J=e_2e_3$ acting by $\ad(J/2)$,
\[
V=\bigl(\tfrac12,\tfrac12\bigr)_0\oplus(0,0)_{+1}\oplus(0,0)_{-1}.
\]
The Lorentz Casimir is $\operatorname{diag}(3,0,0,3,3,3)$ on $(e_1,\dots,e_6)$, and $\ad(J/2)$ vanishes on the four-plane and acts on
$(e_2,e_3)$ by $\bigl(\begin{smallmatrix}0&1\\-1&0\end{smallmatrix}\bigr)$, with eigenvalues $\pm i$.
\end{proposition}

The potential $A=A^\mu\gamma_\mu$ is therefore neutral under the charge, as a photon must be, while the surplus pair is a Lorentz
scalar of charge $\pm1$. The photon potential is the four-plane part of grade one, not the whole of it.

\begin{proposition}[Spin one and spin one-half from the same rotors]
\label{prop:samerotors}
The interior rotations $S_k$ satisfy $[S_i,S_j]=\epsilon_{ijk}S_k$. On grade one $\ad S_k$ acts on $(e_4,e_5,e_6)$ as the standard
rotation generator $r_k$ and annihilates $e_1,e_2,e_3$, so that $-\sum_k(\ad S_k)^2=\operatorname{diag}(0,0,0,2,2,2)$ and
$\spec\ad S_3=\{\pm i,0,0,0,0\}$. On $S=\R^8$, by left multiplication, $-\sum_kS_k^2=\tfrac34\id$ and $S_3$ has eigenvalues $\pm\tfrac i2$.
Consequently $\exp(2\pi\ad S_3)=\id$ on grade one while $\exp(2\pi S_3)=-\id$ on $S$.
\end{proposition}

A vector transforms by conjugation, $A\mapsto RAR^{-1}$ with $R=e^{\theta S_3}$, and a spinor by left multiplication,
$\psi\mapsto R\psi$. At $\theta=2\pi$ the rotor is $-\id$; the sign cancels in the first case and survives in the second. This is the
space-time-algebra account of the difference between vectors and spinors \cite{Hestenes66}, realised here by the rotors of the Dirac
sector: the same $S_k$ give spin one to the potential and spin one-half to matter.

\begin{theorem}[Maxwell's equations as one Clifford equation]
\label{thm:nablaF}
Put $\nabla=\gamma^\mu\partial_\mu$ and $\sigma_k=\gamma_k\gamma_0$.
\begin{enumerate}[label=(\roman*),leftmargin=2.2em]
\item For $A=A^\mu\gamma_\mu$, $\nabla A=\partial\cdot A+\nabla\wedge A$, and for $F:=\nabla\wedge A$ the grade-one part of $\nabla F$ is
$(\Box A^\nu-\partial^\nu\partial\cdot A)\gamma_\nu$, while its grade-three part vanishes identically.
\item $F=\bE+I_4\,\bB$ with $I_4=\gamma_0\gamma_1\gamma_2\gamma_3=-\Omega$ (\cref{prop:bivectorF}), $\bE=\sum_kE_k\sigma_k$,
$E_k=-\partial_kA^0-\partial_tA^k$ and $\bB=\curl\bA$; in terms of $\Omega$, $F=\bE-\Omega\,\bB$.
\item For every $F=\sum_k(E_k\sigma_k+B_kI_4\sigma_k)$,
\[
\nabla F=(\nabla\cdot\bE)\gamma_0+(\curl\bB-\partial_t\bE)^k\gamma_k+I_4\bigl[-(\nabla\cdot\bB)\gamma_0+(\partial_t\bB+\curl\bE)^k\gamma_k\bigr],
\]
with no other grades.
\end{enumerate}
Hence $\nabla F=\mathcal J$, with $\mathcal J=\rho\gamma_0+\mathcal J^k\gamma_k$, is Gauss's law, the Amp\`ere--Maxwell law, the absence of
magnetic charge and Faraday's law, and for $F=\nabla\wedge A$ the last two hold identically.
\end{theorem}

This is Hestenes' single Maxwell equation \cite{Hestenes66}, here on the four-plane of the Dirac sector inside the $8\times8$ matrices
of $\Cl(3,3)$. Plane waves show how the gauge freedom removes the unphysical polarisations and how helicity appears.

\begin{example}[A plane wave of the univectorial photon]
\label{ex:univectorwave}
Let $k=\gamma_0+\gamma_3$, so $k^2=0$, and $\varepsilon=\varepsilon^\mu\gamma_\mu$ with $k\cdot\varepsilon=0$, that is $\varepsilon^0=\varepsilon^3$. Then
\[
F(\varepsilon)=k\varepsilon=\varepsilon^1F_1+\varepsilon^2F_2,\qquad F_j=k\gamma_j ,
\]
because the component of $\varepsilon$ along $k$ drops out, $k^2=0$. The condition $k\cdot\varepsilon=0$ has removed the component along the
other null direction $\gamma_0-\gamma_3$, so only the two transverse polarisations remain, and the gauge freedom
$\varepsilon\mapsto\varepsilon+ck$ is exactly the helicity-zero direction. Moreover $F_j^2=0$ and $F_1F_2+F_2F_1=0$, so every $F(\varepsilon)$ squares
to zero: in the space-time algebra $F^2=\bE^2-\bB^2+2I_4\,\bE\cdot\bB$, so this is the null condition $|\bE|=|\bB|$,
$\bE\perp\bB$. On $\operatorname{span}\{F_1,F_2\}$,
\[
\ad S_3=\begin{pmatrix}0&-1\\1&0\end{pmatrix},\qquad I_4=\begin{pmatrix}0&1\\-1&0\end{pmatrix}=-\ad S_3=-\Omega ,
\]
both with eigenvalues $\pm i$. For positive-frequency waves the eigenvectors of $\ad S_3$ with eigenvalues $\mp i$ carry helicity
$\pm1$, and on them the Hodge star is $i$ times the helicity: circular polarisation is self-duality or anti-self-duality of $F$.
\end{example}

\begin{proposition}[Every field-strength module under the interior action]
\label{prop:censusint}
The Lorentz Casimir on the grades of $\Cl(3,3)$, the number of field-strength modules (Casimir value $4$) and their charges under
$\ad(J/2)$ are
\begin{center}
\begin{tabular}{@{}clcc@{}}
\toprule
grade & spectrum of the Casimir (multiplicity) & field-strength modules & charges\\
\midrule
$0,6$ & $0\,(1)$ & $0$ & ---\\
$1,5$ & $3\,(4)$, $0\,(2)$ & $0$ & ---\\
$2,4$ & $4\,(6)$, $3\,(8)$, $0\,(1)$ & $1$ & $0$\\
$3$ & $4\,(12)$, $3\,(8)$ & $2$ & $\pm1$\\
\bottomrule
\end{tabular}
\end{center}
and on $S$ the Casimir is $\tfrac32$ throughout.
\end{proposition}

Because $J$ commutes with the four-plane, the grade-four module is $J$ times the grade-two one, and $J$ is the complex structure at
one constituent (\cref{thm:i}). Up to that complex structure the neutral field strength is unique and sits at grade two; the two
grade-three copies carry charge $\pm1$ and are not candidates for the photon.

% --- source: tex/s13_spin#3 (The frame photon)
\subsection{The frame photon}

In the frame reading the field strength is the vector $X$ itself, and the boosts of $\lG$ turn out to be the spin matrices of the
photon.

\begin{proposition}[Spin matrices on the frame]
\label{prop:spinmatrices}
$K_k=GL_k$ for $k=1,2,3$, both families commute with $G$ and preserve $\eta$, and the Casimir $\sum K_k^2-\sum L_k^2$ equals $4\,\id$ on
$V$, the field-strength value. On $z=\bB+i\bE$ the rotation $L_k$ acts as $r_k$ and the boost $K_k$ as
\[
s_k=i\,r_k,\qquad (s_k)_{ab}=-i\epsilon_{kab},
\]
which are Hermitian with $[s_i,s_j]=i\epsilon_{ijk}s_k$, $\sum_ks_k^2=2\,\id$ and $\spec s_3=\{-1,0,1\}$. With $p=-i\nabla$ one has
$(s\cdot p)z=\curl z$ identically, so the frame Maxwell system \eqref{eq:RS} is the photon wave equation
\[
i\partial_tz=(s\cdot p)\,z,\qquad p\cdot z=0 ,
\]
and the kernel of $s_3$ is spanned by $\hat e_3$: transversality removes exactly the helicity-zero component of a wave moving along
$\hat e_3$.
\end{proposition}

These are the spin-one matrices that Oppenheimer and Bialynicki-Birula use to write the photon wave function
\cite{Oppenheimer1931,BialynickiBirula96}. The rotations act on $z$ as $r_k$, and their Hermitian generators $s_k=ir_k$ are the action
of the boosts, because on this module a boost is $G$ times a rotation.

\begin{proposition}[Orthogonality on the frame]
\label{prop:orthframe}
For $X=\sum_i(B_ie_i+E_ie_{i+3})$,
\[
z\cdot z=\eta(X,X)-i\,\eta(X,GX).
\]
Hence $\bE\perp\bB$ exactly when $X\perp_\eta GX$, $|\bE|=|\bB|$ exactly when $X$ is $\eta$-null, and both hold exactly when the plane
$\operatorname{span}\{X,GX\}$ is totally $\eta$-null.
\end{proposition}

\begin{theorem}[The axis of a null field]
\label{thm:axis}
For real $a,b\in\R^3$,
\[
\bigl(s\cdot(a\times b)\bigr)(a+ib)=|b|^2a+i|a|^2b-(a\cdot b)(b+ia).
\]
Consequently a nonzero $z=\bB+i\bE$ satisfies $z\cdot z=0$ if and only if $(s\cdot n)z=z$ for some real unit vector $n$, and then
$n=\bB\times\bE/|\bB\times\bE|$ uniquely.
\end{theorem}

\begin{proof}
If $z\cdot z=0$ then $\bB\cdot\bE=0$ and $|\bB|=|\bE|\neq0$; the identity gives $(s\cdot(\bB\times\bE))z=|\bB|^2z$ and
$|\bB\times\bE|=|\bB|^2$. Conversely $n\cdot((s\cdot n)z)=0$ and $z\cdot((s\cdot n)z)=0$ identically, and the characteristic polynomial
of $s\cdot n$ is $-\lambda(\lambda^2-|n|^2)$, so the eigenvalue $1$ forces $n\cdot z=0$ and $z\cdot z=0$. For uniqueness,
$(s\cdot n)z=in\times z=z$ splits into $\bB=-n\times\bE$ and $\bE=n\times\bB$, whence $n\perp\bB$ and $\bB\times\bE=|\bB|^2n$.
\end{proof}

Since $\bB\times\bE=\bE\times\Btext$ is the Poynting vector, a null field is the $m=+1$ state of the photon spin along its own energy
flux, the frame form of the classical fact that a null field determines a unique direction of propagation
\cite{Synge56,Robinson61}. Orthogonality with equal magnitudes therefore singles out the null fields and their axis. It does not
single out a helicity.

\begin{example}[Three null waves with one axis]
\label{ex:threewaves}
The fields
\[
z_\pm=(1,i,0)\,e^{\pm ik(x_3-t)},\qquad z_{\rm lin}=(1,i,0)\cos k(x_3-t)
\]
all solve \eqref{eq:RS}. In each, $\bB\cdot\bE=0$, $|\bB|=|\bE|$ and $\bB\times\bE$ points along $+\hat e_3$ wherever $z\neq0$, so each is an
$m=+1$ state along $\hat e_3$. At $x_3=0$ the electric field turns according to $(\bE\times\dot\bE)\cdot\hat e_3=+k,\,-k,\,0$: helicity
$+1$, helicity $-1$, and a linear superposition. The difference is the sign of the frequency, $z_+\propto e^{-ikt}$ against
$z_-\propto e^{+ikt}$, which is the Bialynicki-Birula assignment for the Riemann--Silberstein vector \cite{BialynickiBirula96}. Not every
radiation field is null either: $z=(1,i,0)\cos k(x_3-t)+(0,1,i)\cos k(x_1-t)$ solves the system and has $|\bB|=|\bE|$, but
$\bB\cdot\bE=\cos k(x_1-t)\cos k(x_3-t)$.
\end{example}

Orthogonality alone is weaker still. The field $X=(\bB;\bE)=(1,0,0;\,0,2,0)$ has $\eta(X,GX)=0$, so $\bE\perp\bB$, but
$\eta(X,X)=1-4=-3$; the boost $e^{\theta K_3}$ with $\theta=-\tfrac12\log3$ carries it to $(0,0,0;\,0,\sqrt3,0)$, a purely electric field
with the same invariant. A field with $\bE\perp\bB$ and $|\bE|\neq|\bB|$ is purely electric or purely magnetic in a suitable frame
\cite{LandauLifshitz2}. On the null line $\C(1,i,0)$, by contrast, $e^{\theta K_3}$ acts by multiplication by $e^\theta$; with
$\theta=-\operatorname{artanh}\beta$ for a frame moving with velocity $\beta\hat e_3$ this is the Doppler factor $\sqrt{(1-\beta)/(1+\beta)}$.
The stabiliser of the axis acts on a null field by the helicity phase and by the Doppler factor.

% --- source: tex/s13_spin#4 (Potentials on the frame, and helicity as a charge)
\subsection{Potentials on the frame, and helicity as a charge}
\label{sec:potentials}

Place a vector potential on each triple, $\bA$ on $e_1,e_2,e_3$ and $\mathbf C$ on $e_4,e_5,e_6$, write $X_A=(\bA;\mathbf C)$ and
$X_F=\curl X_A$, the curl acting on each triple, and take
\begin{equation}
\label{eq:SSL}
\mathcal L=\tfrac12\,(\curl X_A)\cdot G\dot X_A-\tfrac12\,|\curl X_A|^2
\end{equation}
with the Euclidean dot product of $\R^6$, for which $G$ is orthogonal.

\begin{theorem}[Duality-symmetric Maxwell theory on the frame]
\label{thm:SS}
\begin{enumerate}[label=(\roman*),leftmargin=2.2em]
\item The Euler--Lagrange expression of \eqref{eq:SSL} is $G\curl\bigl(\dot X_A+G\curl X_A\bigr)$. Every solution of the first-order
equation $\dot X_A=-G\curl X_A$ solves it, and on a contractible region every solution satisfies the first-order equation up to a
time-dependent gauge transformation.
\item On solutions, $X_F$ obeys the frame Maxwell system \eqref{eq:RS} with $\bB=\curl\bA$ and $\bE=\curl\mathbf C$.
\item Under $X_A\mapsto X_A+(\nabla\lambda_1;\nabla\lambda_2)$, $X_F$ is unchanged and $\mathcal L$ changes by a divergence.
\item The Hamiltonian is $\tfrac12|\curl X_A|^2=\tfrac12(|\bB|^2+|\bE|^2)$, the energy density.
\item $\mathcal L(e^{\theta G}X_A)=\mathcal L(X_A)$ exactly, and the Noether current of $\delta X_A=GX_A$ has density
$j^0=-\tfrac12X_A\cdot\curl X_A$. For fields that decay or are periodic the conserved charge is $-\mathcal K$, where
\[
\mathcal K=\tfrac12\int X_A\cdot\curl X_A\,d^3x=\tfrac12\int(\bA\cdot\bB+\mathbf C\cdot\bE)\,d^3x
\]
is the optical helicity.
\end{enumerate}
\end{theorem}

This is the duality-symmetric first-order formulation of Schwarz and Sen \cite{SchwarzSen94}, with two vector potentials in the
temporal gauge and with the binding root $G$ in the place of their $SO(2)$ generator; the conserved charge of the duality rotation is,
up to sign, the optical helicity \cite{Calkin65,DeserTeitelboim76,CameronBarnettYao12}.

\begin{example}[The $G$-charge is the helicity]
\label{ex:Gcharge}
For waves $e^{i(k x_3-\omega t)}$ the first-order equation reads $\omega a=G(k\times a)$, a $6\times6$ matrix with eigenvalues $\pm k$, each
twice, and $0$ twice (the longitudinal, pure-gauge pair). The eigenspace $\omega=+k$ is spanned by
\[
\begin{aligned}
a_1&=(1,-i)\otimes(1,i,0), & Ga_1&=ia_1, & s_3a_1&=a_1,\\
a_2&=(1,i)\otimes(1,-i,0), & Ga_2&=-ia_2, & s_3a_2&=-a_2,
\end{aligned}
\]
with $s_3$ acting on each triple. For the real waves $X_A=\Re(a\,e^{ik(x_3-t)})$ the helicity density $\tfrac12X_A\cdot\curl X_A$ and the
energy density are $(k,k^2)$ for $a_1$, $(-k,k^2)$ for $a_2$, and $(0,4k^2\sin^2k(x_3-t))$ for $a_1+a_2$. On positive-frequency modes the
eigenvalue of $G$ is $i$ times the helicity. The ratio of helicity to energy is $\pm1/\omega$, which for a quantised field of energy
$N\hbar\omega$ is $\pm N\hbar$: one unit of $\hbar$ per photon, the spin-one value.
\end{example}

The frame potentials are not components of a Lorentz tensor. Neither the constant boost $\delta X_A=K_3X_A$ nor the boost of the
coordinates is a symmetry of \eqref{eq:SSL}; the boost that is one is the Schwarz--Sen transformation $\delta X_A=t\,\partial_3X_A-x_3\,G\curl
X_A$, whose curl is, identically,
\[
\curl\bigl(t\,\partial_3X_A-x_3\,G\curl X_A\bigr)=t\,\partial_3X_F-x_3\,G\curl X_F-K_3X_F ,
\]
which on solutions is the boost of the field strength along $\hat e_3$ with its orbital part and its spin part $K_3$. Boosts reach the
potentials only through their curl and only on shell \cite{SchwarzSen94}. Duality rotations, for their part, are not Lorentz
transformations: $e^{\theta G}$ preserves the energy density and rotates $\eta(X,X)$ into $\eta(X,GX)$.

% --- source: tex/s13_spin#5 (The bridge between the readings)
\subsection{The bridge between the readings}
\label{sec:bridge}

\begin{proposition}[Half-spinors under the two algebras]
\label{prop:halfspinors}
Let $S^\pm$ be the $\pm1$ eigenspaces of the volume element $e_1e_2\cdots e_6$, whose square is $\id$. Under $\lint$ the Casimir on $S^\pm$ is $\tfrac32$, and $S^+$ carries a two-parameter family of invariant antisymmetric forms and no
invariant symmetric one: it is a Weyl module. Under $\lG$ the Casimir on all of $S$ is $3$, the invariant symmetric forms on $S^+$ form a
one-parameter family of Lorentzian forms, and there is no invariant antisymmetric form: $S^+$ is a Minkowski four-vector, and
$\exp(2\pi L_3)=+\id$ on $S$. On $\C^8$ the two algebras give $2\,(\tfrac12,0)\oplus2\,(0,\tfrac12)$ and $2\,(\tfrac12,\tfrac12)$
respectively, and they meet in zero.
\end{proposition}

So under the binding algebra every spinor of $\Cl(3,3)$ has integer spin, in accordance with $\mathrm{Spin}(3,3)\cong SL(4,\R)$, whose
half-spinor modules are $\R^4$ and its dual \cite{Porteous}. This sharpens \cref{prop:emexclusive}: with $\lG$ as the Lorentz algebra
one constituent carries no spin one-half at all.

\begin{theorem}[The Klein correspondence]
\label{thm:klein}
The $\mathrm{Spin}(3,3)$-invariant bilinear forms on $S$ form a two-dimensional space. For one of them, $\beta$, the map
$W(X)(\psi,\chi)=\beta(\psi,X\chi)$ on $S^+$ is antisymmetric and equivariant, and
\[
\Pf W(X)=\eta(X,X).
\]
Hence $V\cong\Lambda^2(S^+)^*$ as $\mathrm{Spin}(3,3)$-modules, and the null cone of $\eta$ is the set of decomposable two-forms.
\end{theorem}

This is the Klein correspondence between the points of the Klein quadric and the lines of projective three-space
\cite{WardWells}. It can be compared with the frame reading. Let $T:\R^{1,3}\to S^+$ be the $\lG$-intertwiner, unique up to scale. Then
\[
T^*W(X)=-F_X+\star F_X=\sqrt2\,F_{e^{3\pi G/4}X}\qquad\text{for all }X,
\]
so the canonical identification of $V$ with two-forms on the half-spinor and the frame reading of the field strength differ by a
duality rotation of an eighth of a turn, $-\pi/4$ modulo $\pi$. In the Klein frame $\eta(X,X)$ is the Pfaffian, the invariant of
$\bE\cdot\bB$ type, and $\eta(X,GX)$ the Hodge pairing; in the frame reading the roles are the other way round. For free fields the two
readings are equivalent, since duality rotations map solutions to solutions; a statement that singles out one invariant, such as which
of them is the Lagrangian, holds in one reading and not in the other. A plane wave $\kappa\wedge\alpha$ built on $S^+$, with $\kappa$ the
lowered null momentum $(1,0,0,1)$ and $\alpha$ the polarisation $(0,1,0,0)$, lands on the doubly null vector
$X=(-\tfrac12,-\tfrac12,0;\,\tfrac12,-\tfrac12,0)$, whose Poynting axis $\bB\times\bE=(0,0,\tfrac12)$ is the spatial momentum; the Klein
planes of $X$ and of $GX$ meet in the null line $[\kappa]$, as two lines of projective space joined by a line of the quadric must.

\begin{figure}[t]
\centering
\begin{tikzpicture}[>=Stealth,font=\small]
\node[draw=tcol,thick,rounded corners=4pt,fill=tcol!6,text width=6.0cm,align=left,inner sep=6pt] (int) at (0,0) {%
\textbf{\textcolor{tcol}{interior reading}}, Lorentz algebra $\lint$\\[3pt]
grade 1: potential $A$, $(\tfrac12,\tfrac12)_0$; surplus pair $(0,0)_{\pm1}$\\
grade 2: field strength $F$, $(1,0)\oplus(0,1)$; $\Omega=\star$\\
spinors: Weyl, $2(\tfrac12,0)\oplus2(0,\tfrac12)$\\
equation: $\nabla F=\mathcal J$ (\cref{thm:nablaF})};
\node[draw=scol,thick,rounded corners=4pt,fill=scol!6,text width=6.0cm,align=left,inner sep=6pt] (bin) at (7.8,0) {%
\textbf{\textcolor{scol}{frame reading}}, Lorentz algebra $\lG$\\[3pt]
grade 1: field strength $X=(\bB;\bE)$; $G=\star$\\
boosts: spin matrices $s_k=ir_k$\\
spinors: four-vectors, $2(\tfrac12,\tfrac12)$\\
equation: $i\partial_tz=(s\cdot p)z$ (\cref{prop:spinmatrices})};
\draw[<->,thick,gray] (int.east) -- node[below,align=center,font=\scriptsize] {not\\ conj.} (bin.west);
\node[draw=gcol,thick,rounded corners=4pt,fill=gcol!6,text width=13.6cm,align=left,inner sep=6pt] (mix) at (3.9,-3.25) {%
\textbf{\textcolor{gcol}{mixed rung $n=2$}}: $\lint$ on constituent 1, $\lG$ on constituent 2. Grade one carries
$V\oplus\Lambda^2V\oplus2\cdot\mathbf1$; the spinor module has a rank-$16$ Dirac block; one grade-three element identifies the
frame field of constituent 2 with $dA$ of constituent 1. Cost: the internal algebra shrinks to one dimension
(\cref{prop:mixedcost}).};
\draw[->,thick,gray] (int.south) -- (int.south |- mix.north);
\draw[->,thick,gray] (bin.south) -- (bin.south |- mix.north);
\node[gray,anchor=south] at (3.9,1.7) {Klein: $V\cong\Lambda^2(S^+)^*$, duality angle $-\pi/4$};
\end{tikzpicture}
\caption{Two places for the photon in the arena (\cref{sec:spin}). Each reading gives a free field with the spin content of the photon;
on one constituent they use non-conjugate Lorentz algebras and do not merge, and at the second rung they can be housed together at
the price of the internal algebra.}
\label{fig:tworoutes}
\end{figure}

% --- source: tex/s14_orth#0 ((head))


% --- source: tex/s14_orth#1 (The orthogonality of the electric and magnetic fields)
\section{The orthogonality of the electric and magnetic fields}
\label{sec:orth}

In the frame reading the orthogonality of $\bE$ and $\bB$ is the reality of the complex invariant, $\Im q=2\,\bE\cdot\bB=0$, and
by \cref{prop:planes} it is also the condition for the spacetime two-form to be a single plane. Since the arena puts $\bE$ and $\bB$
on two orthogonal triples, it is natural to ask what orthogonality means physically: what breaks it, whether its failure is a
symmetry breaking, and whether it can produce mass. The answers are classical once the right questions are asked, and they come out
sharply negative on the last point.

% --- source: tex/s14_orth#2 (What orthogonality is)
\subsection{What orthogonality is}
\label{sec:orth:what}

The field strength $F^\mu{}_\nu$ has eigenvalues $\pm\epsilon_E$ and $\pm i\epsilon_M$ with
\[
\epsilon_E=|\Im\sqrt q|,\qquad \epsilon_M=|\Re\sqrt q|,
\]
and $\bE\cdot\bB=\Pf F=0$ exactly when $F$ is one plane: a pure boost, a pure rotation or a null rotation. Otherwise $F$ generates a
screw, a boost and a rotation about a common axis. Two simple facts locate orthogonality among all fields.

\begin{proposition}[Single modes are orthogonal]
\label{prop:modes}
Let $F=k\wedge b$ with $k,b$ real four-vectors. Then $\Pf F=0$ and $q=(k\cdot k)(b\cdot b)-(k\cdot b)^2$ is real. The field of a mode
$A_\mu=\Re(a_\mu e^{-ik\cdot x})$ with real $k$ and $k\cdot a=0$ is $k\wedge b(x)$ with $k\cdot b=0$, for every polarisation. A Maxwell mode,
$k\cdot k=0$, is therefore null at every point, and a Proca mode, $k\cdot k=m^2$, has $q=m^2(b\cdot b)\le0$: it is orthogonal and of
electric type everywhere.
\end{proposition}

A mass does not show in the angle between the fields of a single mode; it shows in the other invariant, which it makes nonzero.
The null cone $q=0$ is the massless locus, and orthogonality is a much larger set, containing every single mode of every mass.

\begin{proposition}[Superposition breaks it]
\label{prop:twomodes}
$\Pf(k_1\wedge b_1+k_2\wedge b_2)=\det(k_1^\flat,b_1^\flat,k_2^\flat,b_2^\flat)$. At a point, two modes give $\bE\cdot\bB\neq0$ exactly when
their two planes together span spacetime. Waves propagating along one null direction $\ell$ never do: the total field is
$\ell\wedge\sum_j\omega_jb_j$, a single null plane.
\end{proposition}

\begin{example}[Counter-propagating waves]
\label{ex:counter}
In textbook variables, with $\omega=k$, let a wave along $+z$ polarised along $x$ meet a wave along $-z$. With crossed polarisations,
$\bE=(\cos(kz-kt),\cos(kz+kt),0)$ and $\Btext=(\cos(kz+kt),\cos(kz-kt),0)$; with parallel ones,
$\bE=(\cos(kz-kt)+\cos(kz+kt),0,0)$ and $\Btext=(0,\cos(kz-kt)-\cos(kz+kt),0)$. Both solve Maxwell's equations, and
\[
\begin{aligned}
&\text{crossed:}&& \bE\cdot\Btext=\cos2kz+\cos2kt, && |\Btext|^2-|\bE|^2=0,\\
&\text{parallel:}&& \bE\cdot\Btext=0, && |\Btext|^2-|\bE|^2=-2(\cos2kz+\cos2kt).
\end{aligned}
\]
The two configurations have the same energy density and the same Poynting vector at every point. The reason is structural: the crossed
configuration is the parallel one turned by $\pi/4$ about the $z$ axis and then duality-rotated by $\pi/4$; the rotation preserves the
energy density and the Poynting vector along $z$, the duality rotation preserves the entire energy--momentum tensor, and only the phase
of $q$ moves.
\end{example}

Nothing exotic is needed to break orthogonality. Near the ground the fair-weather electric field of the atmosphere, about $100$~V/m
and vertical \cite{FeynmanII}, meets the geomagnetic field, which is inclined to the horizontal everywhere except on the magnetic
equator.

% --- source: tex/s14_orth#3 (How it changes in time)
\subsection{How it changes in time}

\begin{proposition}[The local law]
\label{prop:EBlocal}
With $\partial_t\bB=\curl\bE$ and $\partial_t\bE=-\curl\bB-\bj$,
\[
\partial_tq=-2i\,z\cdot(\curl z)-2i\,z\cdot\bj,\qquad \partial_t(\bE\cdot\bB)=\bE\cdot\curl\bE-\bB\cdot\curl\bB-\bj\cdot\bB ,
\]
in textbook variables $\partial_t(\bE\cdot\Btext)=\Btext\cdot\curl\Btext-\bE\cdot\curl\bE-\bj\cdot\Btext$. With potentials,
$\Btext=\curl\bA$, $\bE=-\partial_t\bA-\nabla\phi$, and in a source-free region $\bE=-\curl\mathbf C$, $\Btext=-\partial_t\mathbf C+\nabla\chi$,
\[
\partial_t(\bA\cdot\Btext)+\nabla\cdot(\phi\Btext+\bE\times\bA)=-2\,\bE\cdot\Btext,\qquad
\partial_t(\mathbf C\cdot\bE)-\nabla\cdot(\Btext\times\mathbf C+\chi\bE)=-2\,\bE\cdot\Btext .
\]
\end{proposition}

In vacuum, $\bE\cdot\bB$ changes at a point only through $\bE\cdot\curl\bE$ and $\bB\cdot\curl\bB$, which measure how much each field
twists about itself. Integrated over a region with vanishing boundary fluxes, $\int\bE\cdot\Btext$ is minus half the rate of change of
the magnetic helicity $\int\bA\cdot\Btext$ \cite{Woltjer58,BergerField84}, and equally of the electric helicity, so that their
difference, the optical helicity of \cref{thm:SS}, is conserved in free space \cite{Calkin65,CameronBarnettYao12}. Non-orthogonality
is the agent that changes the helicity of each field while their combination is kept. In matter, charges moving along the magnetic field
change $\bE\cdot\bB$ through $-\bj\cdot\bB$.

Conductors screen it. For the generalised Ohm law
$\bE+\bv\times\Btext=\eta_{\rm R}\bj+(\bj\times\Btext)/(ne)-(\nabla\cdot\mathsf P_e)/(ne)+(m_e/ne^2)\,d\bj/dt$ \cite{PriestForbes},
\[
\bE\cdot\Btext=\eta_{\rm R}\,\bj\cdot\Btext-\frac{\Btext\cdot(\nabla\cdot\mathsf P_e)}{ne}+\frac{m_e}{ne^2}\,\Btext\cdot\frac{d\bj}{dt} :
\]
the flow term and the Hall term drop out, ideal magnetohydrodynamics enforces orthogonality everywhere, and only resistivity, electron
pressure and electron inertia allow $\bE\cdot\Btext\neq0$. This is the criterion of Schindler, Hesse and Birn for magnetic reconnection,
a nonvanishing integral of $E_\parallel$ along field lines \cite{SchindlerHesseBirn88}. Where charges are scarce the screening fails: a
rotating magnetised neutron star in vacuum induces $\bE\cdot\Btext\neq0$ at its surface, charges drawn from the surface fill the
magnetosphere until $\bE\cdot\Btext=0$ \cite{GoldreichJulian69}, and the regions left without enough charge keep $\bE\cdot\Btext\neq0$ and
accelerate particles \cite{RudermanSutherland75}.

% --- source: tex/s14_orth#4 (Non-orthogonality breaks parity and time reversal, and is ne)
\subsection{Non-orthogonality breaks parity and time reversal, and is never spontaneous}

\begin{theorem}[Orthogonality is achirality]
\label{thm:achiral}
For a field strength $F$ at a point the following are equivalent: (a) $\bE\cdot\bB=0$; (b) the mirror image $PF$ equals $\Lambda F$ for some
proper orthochronous Lorentz transformation $\Lambda$; (c) the time-reversed field $TF$ equals $\Lambda F$ for some such $\Lambda$.
\end{theorem}

\begin{proof}
If $\bE\cdot\bB=0$ and $\bB\neq0$, the half-turn about $\bB$ fixes $\bB$ and reverses every vector perpendicular to it, sending $(\bB;\bE)$ to
$(\bB;-\bE)=P(\bB;\bE)$; if $\bB=0$ a half-turn about an axis perpendicular to $\bE$ does the same. Conversely the Pfaffian is invariant
under proper Lorentz transformations and changes sign under $P$, so $\Pf F=-\Pf F$. For (c), $T=CP$ on fields and $CF=-F$ has the same
invariant $q$; since $\lG\cong\so{3,\C}$ acts transitively on each level set of $q$ with the origin removed, $-F$ lies in the proper
Lorentz orbit of $F$, and (c) is equivalent to (b).
\end{proof}

A non-orthogonal field is therefore chiral in Kelvin's sense, a screw of definite handedness, and the sign of $\bE\cdot\bB$ is that
handedness; it is odd under $P$ and $T$ and even under $C$, $PT$ and $CPT$. This is a statement about a field at a point. A configuration
can be orthogonal everywhere and still chiral as a whole: a static Beltrami magnetic field with $\bE=0$ carries a nonzero helicity
density. The global chirality of a magnetic configuration is its helicity, $P$-odd and $T$-even, and $\bE\cdot\Btext$, $P$-odd and
$T$-odd, is minus half its rate of change.

For the Maxwell energy the breaking is never spontaneous. Since $u^2-|\mathbf S|^2=\tfrac14|q|^2$ for the energy density $u=\tfrac12(|\bE|^2+|\bB|^2)$ and the
Poynting vector $\mathbf S$,
\begin{equation}
\label{eq:EBbound}
u\ \ge\ \tfrac12|q|\ \ge\ |\bE\cdot\bB| ,
\end{equation}
and the minimum of the energy is the zero field; there is no degenerate family of lowest-energy configurations among which
$\bE\cdot\bB$ could choose a sign. In the quantum theory the corresponding statement is the theorem of Vafa and Witten: in a vector-like
gauge theory at $\theta=0$, zero temperature and zero density, parity is not broken spontaneously through a Bose order parameter such
as $F\tilde F$ \cite{VafaWitten84}. When $\bE\cdot\Btext$ is imposed from outside the vacuum responds, for instance by rotating the
chiral condensate of quantum chromodynamics through the electromagnetic triangle anomaly \cite{CaoHuang15}; that is the response of a
vacuum to an explicit breaking, not a breaking of its own. Both statements have hypotheses that are part of them. The classical one uses the
quadratic Maxwell energy, whose minimum is the zero field; it does not cover modified actions whose energy has minima at nonzero fields. The quantum
one needs a vector-like theory with a positive Euclidean measure ($\theta=0$, zero density); it says nothing about chiral sectors such as the
electroweak block, or about actions with an explicit $\bE\cdot\bB$ term.

Duality behaves differently. The rotation $X\mapsto e^{\vartheta G}X$ leaves the energy--momentum tensor invariant and multiplies $q$ by
$e^{2i\vartheta}$, so for $q\neq0$ only $\vartheta\in\pi\Z$ fixes the invariants, while for $q=0$ every rotation does. The fields whose
invariants are unchanged by duality are the null ones, not the orthogonal ones, and orthogonality is not a symmetric point of duality.
On $\R^{3,3}$, moreover, $(e^{\vartheta G})^{\!\top}\eta\,e^{\vartheta G}=\cos2\vartheta\,\eta+\sin2\vartheta\,\eta G$, so $e^{\vartheta G}$ is an
isometry of $\eta$ only for $\vartheta\in\pi\Z$. Duality is broken explicitly by electric sources, as \cref{sec:mixed} showed, and by the
vacuum of a theory with charged particles. For a Lagrangian $L(\cF,\cG)$ of the invariants $\cF=\tfrac12(|\bE|^2-|\Btext|^2)$ and
$\cG=\bE\cdot\Btext$, with $\mathbf D=\partial L/\partial\bE$ and $\mathbf H=-\partial L/\partial\Btext$,
\[
\mathbf D\cdot\mathbf H-\bE\cdot\Btext=\cG\bigl(L_\cF^2-L_\cG^2-1\bigr)-2\cF L_\cF L_\cG ,
\]
which vanishes for the Maxwell and Born--Infeld Lagrangians, the duality-invariant ones \cite{GaillardZumino81,GibbonsRasheed95}. For the
weak-field Euler--Heisenberg Lagrangian \cite{HeisenbergEuler36,Dunne04} it is $-\tfrac{8\alpha^2}{15m^4}\cF\cG+O(\alpha^3)$: the duality
defect of the vacuum of quantum electrodynamics is proportional, at leading order, to $\bE\cdot\Btext$. For every parity-even Lagrangian
$L=h(\cF,\cG^2)$ it carries the explicit factor $\cG$, so it vanishes on every orthogonal field.

% --- source: tex/s14_orth#5 (Non-orthogonality does not generate mass)
\subsection{Non-orthogonality does not generate mass}
\label{sec:orth:nomass}

Mass can enter in four ways: as the rest energy of the field itself, as a photon mass, as a mass of charged matter in the background, and
through an $\bE\cdot\bB$-type term in an action. None of them is produced by non-orthogonality.

\emph{Rest energy.} For $q\neq0$ the field has a local rest frame, moving with velocity $\mathbf S/(u+\tfrac12|q|)$, in which the Poynting
vector vanishes and the energy density is $\tfrac12|q|=u\sqrt{1-|\mathbf v_E|^2}$, $\mathbf v_E=\mathbf S/u$; null fields transport their
energy at the speed of light and have no rest frame. The rest-frame energy density is the modulus of $q$, orthogonality its phase, and
duality turns the phase at fixed energy, momentum and stress. For a finite configuration with total energy $U$ and momentum $\mathbf P$,
$M^2=U^2-|\mathbf P|^2$ is positive whenever $q\neq0$ on a set of positive measure, but $M$ does not detect orthogonality: the Cauchy data
$(\bE,\Btext)=(0,\mathbf b)$ and $(\mathbf b\sin\vartheta,\mathbf b\cos\vartheta)$ for a compactly supported divergence-free $\mathbf b$ have the same
energy density and zero Poynting vector at every point, hence the same $M$, while $\bE\cdot\Btext$ is $0$ for the first and
$\tfrac12\sin2\vartheta\,|\mathbf b|^2$ for the second.

\emph{Photon mass.} Every Lorentz-covariant linear map from two-forms to vectors is zero; more strongly, a continuous covariant
vector-valued function of the value of $F$ at a point vanishes, because a non-null field is fixed by a rotation and a boost about a common
axis that have no common fixed vector, and every tensor power of $(1,0)\oplus(0,1)$ contains only representations $(j_1,j_2)$ with integer
$j_1,j_2$. A mass term in $\partial_\mu F^{\mu\nu}=j^\nu$ must therefore involve the potential, other fields or derivatives of $F$; no
function of the two invariants can serve. Derivatives give a massive mode that is a ghost, as Podolsky's Lagrangian shows. A constant
nonlinear background moves the light cone without opening a gap, $\det W(s\omega,s\bk)=s^6\det W(\omega,\bk)$, and to first order the two
phase velocities are $1-c_i\Theta^{\mu\nu}k_\mu k_\nu$ for $L=\cF+c_1\cF^2+c_2\cG^2$, which is duality invariant and blind to orthogonality.
Written into the field equations, the binding rotation of \cref{sec:neutral} gives the frequencies $\kappa\pm|\bk|$ to the two circular
polarisations, a frequency offset that the time-dependent duality rotation $z=e^{-i\kappa t}w$ removes.

\emph{Charged matter.} The invariant $\epsilon_E$ vanishes exactly when $\bE\cdot\bB=0$ and $|\bB|\ge|\bE|$, so every non-orthogonal field
has $\epsilon_E>0$ and makes the vacuum of spinor electrodynamics decay into pairs \cite{Schwinger51}, converting field energy into
particles whose mass is a given parameter. For massless fermions the same field pumps chirality, $dN_5/dt=(e^2/2\pi^2)\int\bE\cdot\Btext$,
the axial anomaly \cite{Adler69,BellJackiw69,NielsenNinomiya83}; a fermion mass relaxes it, and none is produced. The background that does
catalyse a fermion mass is a pure magnetic field, which is orthogonal \cite{GusyninMiranskyShovkovy95}.

\emph{Topological terms.} In four dimensions the abelian term $\theta\,\bE\cdot\Btext$ is a total derivative and invisible to the field
equations; a varying $\theta$ makes it a source, as in axion electrodynamics \cite{Wilczek87}, and a constant gradient gives
birefringence \cite{CarrollFieldJackiw90}. For a timelike gradient both branches remain gapless; for a spacelike gradient $p^\mu$ the dispersion relation
$(k\cdot k)^2+(k\cdot k)(p\cdot p)-(k\cdot p)^2=0$ gives, at zero wave vector, $\omega=0$ on one branch and $\omega=|\mathbf p|$ on the other. That gap is a
Lorentz-violating dispersion fixed by an external pseudoscalar, not a Proca mass, and it does not arise without the varying $\theta$. In every known
mechanism in which a term of this type makes mass, something is
essential that a single abelian field on a topologically trivial space lacks: the Chern--Simons term in three dimensions \cite{DJT82}, the
anomaly in the two-dimensional Schwinger model \cite{Schwinger62}, the topological susceptibility of a non-abelian theory with light
fermions for the $\eta'$ \cite{Witten79eta,Veneziano79}, or a pseudoscalar partner for the axion \cite{Weinberg78,Wilczek78}. The frame field
of this arena is a single abelian field strength with no fermions and no pseudoscalar partner, so no route from $\bE\cdot\bB$ to mass is
open within what the arena has built.

The Dirac field supplies an exact counterpart, collected in \cref{tab:diracmaxwell}. Its scalar and pseudoscalar densities
$\sigma=\bar\psi\psi$ and $\omega=\bar\psi i\gamma_5\psi$ turn under a chiral rotation $e^{i\alpha\gamma_5}$ by the angle $2\alpha$, as
$(|\bB|^2-|\bE|^2,2\bE\cdot\bB)$ turn under duality; a plane wave has $\omega=0$ as a single mode has $\Im q=0$; and a pseudoscalar mass enters
the mass shell only through $\sqrt{m_s^2+m_p^2}$. The pseudoscalar density is a phase, not a mass, and $\bE\cdot\bB$ has exactly its place.

\begin{table}[t]
\centering
\small
\begin{tabular}{@{}ll@{}}
\toprule
Maxwell field & Dirac field\\
\midrule
duality $e^{\vartheta G}$, turning $q$ by $2\vartheta$ & chiral rotation $e^{i\alpha\gamma_5}$, turning $\sigma+i\omega$ by $2\alpha$\\
energy--momentum tensor invariant & vector current $J^\mu$ invariant\\
$u^2-|\mathbf S|^2=\tfrac14|q|^2$ & $J^\mu J_\mu=\sigma^2+\omega^2$\\
$|\bB|^2-|\bE|^2$ and $2\,\bE\cdot\bB$ & $\sigma$ and $\omega$\\
null field, $q=0$ & Weyl spinor, $\sigma=\omega=0$\\
one mode: $q$ real & one plane wave: $\omega=0$\\
two modes whose planes span spacetime: $\Im q\neq0$ & both signs of frequency: $\omega\neq0$\\
$\arg q$, twice the Rainich complexion & $\arg(\sigma+i\omega)$, the Yvon--Takabayasi angle\\
\bottomrule
\end{tabular}
\caption{The invariants of a field strength and the bilinears of a Dirac spinor. Each line is an identity or a worked example; the
parallel is exact at the level of these identities and is not an identification of the two fields.}
\label{tab:diracmaxwell}
\end{table}

In the interior dynamics of the next part, orthogonality is broken by the bound motion itself. On an orbit without momentum the flow
is the duality rotation $z\mapsto e^{-i\kappa\tau}z$, so $\bE\cdot\bB(\tau)=\tfrac12|q_0|\sin(\phi_0-2\kappa\tau)$, which vanishes four times per
period, and the rest-frame field turns from magnetic through screws to electric and back at the fixed rest-frame energy density
$\tfrac12|q_0|$. That is the content of \cref{fig:qplane} read along a trajectory, and \cref{sec:neutral} derives it. It is a rotation of
the phase of $q$ at fixed modulus, and by everything in this section it moves no mass.

% --- source: tex/s14b_energy#0 ((head))


% --- source: tex/s14b_energy#1 (The undivided energy and its differentiation)
\section{The undivided energy and its differentiation}
\label{sec:energy}

The energy density of a field point, $\tfrac12(|\bB|^2+|\bE|^2)$, is a single quantity: it treats the magnetic and the electric triple
alike and it is unchanged by duality. A field point also has a rest energy, the energy density seen in a frame in which its Poynting vector
vanishes. This section asks how the second is contained in the first: which quadratic energy the arena singles out and what its symmetry is,
how that energy differentiates into a Lorentz-invariant rest part and a part carried by motion without any interaction or second field, where
the conversion constant $c$ enters, and which of the arena's own structures can carry the name of a mass. The relation $E=mc^2$
\cite{Einstein1905} is inherited; what is at stake is only how the arena realises its ingredients. Most of the answers are inherited
electrodynamics written in the arena's variables, and are said to be so; the arena adds which of its forms plays which role. Throughout, $c$
enters only as the conversion between the two triples of \cref{thm:onec}, never as a symmetry.

\paragraph{Notation.} We use the field-strength reading of \cref{thm:emfield}: $X=\sum_i(B_ie_i+E_ie_{i+3})$, $z=\bB+i\bE\in\C^3$, the
binding root $G$ acts as multiplication by $i$, and $q(z)=z\cdot z$ has $\Re q=\eta(X,X)=|\bB|^2-|\bE|^2$ and $\Im q=2\,\bE\cdot\bB=-\eta(X,GX)$.
With the Hermitian product $h(z,w)=\sum_j\bar z_jw_j$ we put
\begin{equation}
\label{en:defs}
\begin{gathered}
E_6:=\tfrac12|z|^2=\tfrac12\bigl(|\bB|^2+|\bE|^2\bigr),\qquad E_0:=\tfrac12|q(z)|,\\
\mathbf S:=\bE\times\bB=\tfrac i2\,\bar z\times z=-\mathbf S_{\rm t},\qquad \mathbf S_{\rm t}:=\bE\times\Btext .
\end{gathered}
\end{equation}
$E_6$ is the energy density $u$ of \cref{sec:orth}, $\mathbf S$ the Poynting vector of the frame (\cref{ex:circular}) and $\mathbf S_{\rm t}$ the
textbook one. Applied to a frame field $X$ these are the frame field's energies; for the momentum field $\pi=\eta^{-1}p$ of the interior dynamics studied below, \eqref{eq:splitL}
and \cref{sec:closure} we write $E_6(\pi)$, $E_0(\pi)$ and $\mathbf S(\pi)$. Units are Heaviside--Lorentz with $c=1$ until
\cref{sec:energy-c}. The generators $L_i,K_i$ of $\lG$ (\cref{thm:emfield}) act by $\delta\bB=e_i\times\bB$, $\delta\bE=e_i\times\bE$ and by
$\delta\bE=e_i\times\bB$, $\delta\bB=-e_i\times\bE$ respectively, so that $\boldsymbol\omega\cdot L+\bv\cdot K$ acts on $z$ as
$z\mapsto(\boldsymbol\omega+i\bv)\times z$. We write $\mathfrak p_G$ for the part of $\so{3,3}$ anticommuting with $G$, so that
$\so{3,3}=\lG\oplus\mathfrak p_G$, and the Lorentz-class inertia of \cref{thm:inertiareq} as $N=(\nu_1\eta^{-1}+\nu_2G\eta^{-1})/m$, its
coefficients renamed so that the letter $b$ is free for a gauge generator in \cref{sec:energy-transport}.

% --- source: tex/s14b_energy#2 (One undivided energy, with pointwise symmetry U(3))
\subsection{\texorpdfstring{One undivided energy, with pointwise symmetry $U(3)$}{One undivided energy, with pointwise symmetry U(3)}}
\label{sec:energy-unique}

\begin{theorem}[Uniqueness of the undivided energy]
\label{en:unique}
\begin{enumerate}[label=(\roman*),leftmargin=2.2em]
\item The symmetric forms on $V$ invariant under the diagonal rotations are spanned by $|\bB|^2$, $|\bE|^2$ and $\bB\cdot\bE$, and they split as
$\R\,\id\oplus\operatorname{span}\{\eta,\eta G\}$, the first summand commuting and the second anticommuting with $G$.
\item A quadratic energy invariant under the diagonal rotations and under the duality rotations $e^{\vartheta G}$ is a multiple of $E_6$, which is
positive definite.
\item The forms invariant under all of $\lG$ are $\operatorname{span}\{\eta,\eta G\}$, that is $|\bB|^2-|\bE|^2$ and $\bE\cdot\bB$; none of them is
definite.
\end{enumerate}
\end{theorem}

\begin{proof}
The diagonal rotations act on $V=V^+\oplus V^-$ as $\R^2\otimes\R^3$ with $SO(3)$ on the second factor, which is absolutely irreducible, so by
Schur's lemma \cite{GoodmanWallach} the invariant forms are $N=\bigl(\begin{smallmatrix}\alpha&\gamma\\ \gamma&\beta\end{smallmatrix}\bigr)\otimes\id$.
With $G=\bigl(\begin{smallmatrix}0&-1\\1&0\end{smallmatrix}\bigr)\otimes\id$ one finds $\tfrac12(N-GNG)=\tfrac12(\alpha+\beta)\id$ and
$\tfrac12(N+GNG)=\tfrac12(\alpha-\beta)\eta-\gamma\,\eta G$. Since $G^{\!\top}G=\id$ and $G^{\!\top}=-G$, $N$ is invariant under $e^{\vartheta G}$ iff
$G^{\!\top}N+NG=0$ iff $[N,G]=0$, which gives (ii). For (iii), the boost law gives
$\delta(\alpha|\bB|^2+\beta|\bE|^2+2\gamma\,\bB\cdot\bE)=-2(\alpha+\beta)\,\bv\cdot(\bE\times\bB)$, the $\gamma$ term being unchanged, so $\beta=-\alpha$.
\end{proof}

The theorem restates \cref{thm:inertiareq}(iii) and \cref{prop:lGhoms} on $V$ rather than on $V^*$, and it is inherited: for every conjugate
root $G'=gGg^{-1}$, $g\in SO(3,3)$, the forms invariant under the compact part of $\mathfrak l_{G'}$ and under $G'$ are the multiples of the
positive form $g^{-\top}g^{-1}$, so the statement could not have failed. It is the uniqueness of the energy density among rotation- and
duality-invariant quadratic forms. What the arena adds is that this form is the Euclidean form of its own basis: $G$ is $\eta$-self-adjoint and
anticommutes with $\Ad_{e_1e_2e_3}|_V=\eta$ (\cref{lem:Gclifford}), so $G^{\!\top}G=\eta G\eta G=\id$. $E_6$ has duality weight zero, while
$e^{\vartheta G}$ turns $(\Re q,\Im q)$ through $2\vartheta$. In this sense $E_6$ is undivided: it distinguishes neither the magnetic triple from the
electric one nor, as \cref{sec:energy-identity} shows, rest from motion.

\begin{proposition}[The pointwise symmetry of the undivided energy]
\label{en:u3}
\begin{enumerate}[label=(\roman*),leftmargin=2.2em]
\item The maps of $V$ that preserve $E_6$ and commute with $G$ form $\uu3$, of dimension $9$: $A=\bigl(\begin{smallmatrix}P&Q\\-Q&P\end{smallmatrix}\bigr)$
acts on $z$ by the anti-Hermitian matrix $P-iQ$.
\item The Lie algebra of $O(3,\C)\times U(1)_G$, the $G$-linear maps preserving $q$ up to a phase, is $\lG\oplus\R G$, of dimension $7$.
\item The intersection of the two is $\so3\oplus\R G$, of dimension $4$, the maximal compact subalgebra of $\lG\oplus\R G$; and $\uu3\cap\lG=\so3$.
\item The stabiliser of the function $E_0$ in $U(3)$ is $U(1)\cdot O(3)=\{e^{i\varphi}O\}$. For instance the unitary map
\[
U=\frac1{\sqrt2}\begin{pmatrix}1&1&0\\ i&-i&0\\ 0&0&\sqrt2\end{pmatrix}
\]
carries the field $(1,0,0)$, at rest with $E_0=E_6=\tfrac12$, to the null field $\tfrac1{\sqrt2}(1,i,0)$ at the same $E_6$.
\end{enumerate}
\end{proposition}

\begin{proof}
(i) Commuting with $G$ forces the block form, which is multiplication by $P-iQ$ on $z$; $A^{\!\top}=-A$ means $P^{\!\top}=-P$ and $Q^{\!\top}=Q$, so
$P-iQ$ is anti-Hermitian, and $3+6=9$. (ii) $q(e^{tM}z)=e^{i\varphi t}q(z)$ gives $M^{\!\top}+M=i\varphi\,\id$, so $M=\tfrac{i\varphi}2+M_0$ with
$M_0\in\so{3,\C}=\lG$ (\cref{thm:GLorentz}), and $i\,\id$ is $G$. (iii) A complex antisymmetric anti-Hermitian matrix is real. (iv) If $U\in U(3)$
preserves $|q|$, the form $z^{\!\top}U^{\!\top}Uz$ vanishes on the cone $q=0$; $q$ is irreducible, so $U^{\!\top}U=e^{i\varphi}\id$, and
$O=e^{-i\varphi/2}U$ is unitary and complex orthogonal, hence real.
\end{proof}

All four parts are inherited Hermitian linear algebra, which makes the word \emph{unitary} literally true of the pointwise energy; part (iv) is
the stabiliser of $|z\cdot z|$ among unitary maps of the Riemann--Silberstein vector \cite{Silberstein1907,BialynickiBirula13}. Read as a
statement about the rest energy, (iv) says exactly in which sense $E_0$ reduces a symmetry: $U(3)$ is reduced to $U(1)\cdot O(3)$, and the
five-dimensional orbit $U(3)/U(1)\cdot O(3)$ parametrises the real structures $z\mapsto\bar z$ up to phase, of which the arena has chosen one
(\cref{en:parity}). But $U(3)$ is a symmetry of the pointwise form $E_6$ only. It is not a symmetry of Maxwell's equations on the frame,
$\partial_tz=-i\curl z$ (\cref{prop:RS}): the map $U$ of (iv) carries the wave $(1,i,0)e^{ik(x_3-t)}$ of \cref{ex:circular} to
$\tfrac{1+i}{\sqrt2}(1,1,0)e^{ik(x_3-t)}$, which does not solve them. Nor is it a symmetry of the split-form interior, since $G=i\,\id$ lies in
$U(3)$ and reverses $\eta$, and $U(3)\cap O(\eta)=O(3)$. Of the closures of \cref{sec:closure} it is a symmetry only for the inertia $N=\id/m$,
because the $U(3)$-invariant symmetric forms are the multiples of $\id$.

% --- source: tex/s14b_energy#3 (The rest energy is the complete invariant of Lorentz transfo)
\subsection{The rest energy is the complete invariant of Lorentz transformations and duality}
\label{sec:energy-orbits}

\begin{theorem}[Orbits]
\label{en:orbits}
\begin{enumerate}[label=(\roman*),leftmargin=2.2em]
\item Each level set $\{|q|=r^2\}$, $r>0$, is one orbit of $O(3,\C)\times U(1)_G$, of real dimension $5$, with two-dimensional stabilisers; every
non-null $z$ is carried to $(r,0,0)$ with $r=\sqrt{2E_0}$.
\item The nonzero null fields form one orbit of $SO(3,\C)$ alone, of dimension $4$, and at a null point $Gz$ is tangent to the Lorentz orbit.
\item Hence $E_0$ is a complete invariant of Lorentz transformations and duality on non-null fields, and $E_6$ is the complete invariant of
$U(3)$, whose orbits are the spheres.
\end{enumerate}
\end{theorem}

\begin{proof}
(i) Write $q(z)=r^2e^{i\varphi}$ and $w=e^{-i\varphi/2}z/r$, a duality rotation, so that $q(w)=1$. Since $q$ is nondegenerate on
$w^\perp=\{f:f\cdot w=0\}$, some $f\in w^\perp$ has $f\cdot f=1$; with $f_3=w\times f$, the identity
$(a\times b)\cdot(a\times b)=(a\cdot a)(b\cdot b)-(a\cdot b)^2$ gives $f_3\cdot f_3=1$, so $A=(w\,f\,f_3)^{\!\top}\in SO(3,\C)$ and
$Ae^{-i\varphi/2}z=re_1$. An element $M_0+\tfrac{i\varphi}2$ fixes $e_1$ iff $\varphi=0$ and $M_0$ is a complex rotation about $e_1$; the stabiliser is
$\so{2,\C}$ and the orbit has dimension $7-2=5$. (ii) A rotation carries a null $z$ to $\lambda(1,i,0)$ with $\lambda>0$. $K_3$ acts by
$\delta z=i\,e_3\times z$, and $i\,e_3\times(1,i,0)=(1,i,0)$, so the boost $e^{\chi K_3}$ with $\chi=-\log\lambda$ finishes. The Lorentz tangent space
at a null $z$ is its $q$-orthogonal complement, which contains $iz$; at $(1,i,0)$, $iz=-e_3\times z=-L_3z$ is the action of a rotation about $e_3$. (iii)
follows from (i) and (ii).
\end{proof}

This is inherited, the classification of a field strength at a point \cite{Synge56,LandauLifshitz2}. Without duality the complete invariant
on nonzero fields is $q$ itself (the transitivity used in the proof of \cref{thm:achiral}), and duality removes its phase, which is twice the
complexion of \cref{prop:planes}. At the null point $(1,i,0)$ the tangent $Gz=-L_3z$ is generated by the rotation about the propagation axis, the identity of the $G$-charge and the
helicity of \cref{ex:Gcharge}.

% --- source: tex/s14b_energy#4 (The differentiation identity)
\subsection{The differentiation identity}
\label{sec:energy-identity}

The relation $u^2-|\mathbf S|^2=\tfrac14|q|^2$ was used in \cref{sec:orth} for the bound \eqref{eq:EBbound}. Its structure is the following.

\begin{proposition}[Differentiation of the undivided energy]
\label{en:split}
For every $z\in\C^3$
\begin{equation}
\label{en:lagrange}
|z|^4=|z\cdot z|^2+|\bar z\times z|^2,\qquad\text{equivalently}\qquad E_6^2=E_0^2+|\mathbf S|^2 .
\end{equation}
The rank-one Hermitian matrix $H=zz^\dagger$ carries all three terms: $\tr H=2E_6$, $(\Im H)_{jk}=\epsilon_{jkl}S_l$, $\tr(H\bar H)=|q|^2=4E_0^2$ and
$H^2=(\tr H)H$. Hence $E_0\le E_6$, with equality iff $\mathbf S=0$, iff $\bar z\parallel z$, iff $\bB\parallel\bE$ or one of them vanishes.
\end{proposition}

\begin{proof}
$\bar z\times z=-2i\,\bE\times\bB$, and $(a\times b)\cdot(c\times d)=(a\cdot c)(b\cdot d)-(a\cdot d)(b\cdot c)$ at $(\bar z,z,z,\bar z)$ gives
$|\bar z\times z|^2=|z|^4-|q|^2$. Next $H_{jk}=B_jB_k+E_jE_k+i\epsilon_{jkl}(\bE\times\bB)_l$ and $\tr(H\bar H)=|\sum_jz_j^2|^2$. The equality case is the
equality case of the Cauchy--Schwarz inequality $|h(\bar z,z)|\le|z|^2$.
\end{proof}

The identity is inherited, a complex Lagrange identity. It is the exact form of the differentiation asked about: the undivided $E_6$ is the
hypotenuse of a right triangle whose legs are a Lorentz-invariant rest part $E_0$ and a motion part $|\mathbf S|$, and the split needs no
interaction and no second field. In Minkowski terms $E_0$ is the Lorentz length of $(E_6,\mathbf S)$, although $(E_6,\mathbf S_{\rm t})$ is not a
four-vector, as the next proposition makes precise. Three quantities must be kept apart. $E_0$ is a pointwise invariant: for a non-null field it
is the energy density in the local frame where $\mathbf S=0$, and it vanishes on null fields. It is not the invariant mass of a localised field
configuration, $M=\bigl((\int E_6)^2-|\int\mathbf S|^2\bigr)^{1/2}$ with $\Theta^{0\mu}$ integrated over space. The vectors $(E_6,\mathbf S)(x)$ are causal, so by the
reverse triangle inequality $M\ge\int E_0\,d^3x$, with equality only when all points share one rest frame. And neither is a generated particle mass: the
mass $m$ of the Dirac form of \cref{sec:dirac} is an input of the Lagrangian. To avoid a clash with the trivector $e_1e_2e_3$ we write $\Theta^{\mu\nu}$ for the
stress--energy tensor.

\begin{proposition}[Rainich identity]
\label{en:rainich}
For the stress--energy tensor of $(\bE,\Btext)$, $\Theta^{00}=\tfrac12(|\bE|^2+|\Btext|^2)=E_6$, $\Theta^{0i}=(\mathbf S_{\rm t})_i$ and
$\Theta^{ij}=\Theta^{00}\delta_{ij}-E_iE_j-B_{{\rm t}i}B_{{\rm t}j}$, one has $\Theta^\mu{}_\mu=0$ and
\begin{equation}
\label{en:rainicheq}
\Theta^\mu{}_\lambda\Theta^\lambda{}_\nu=E_0^2\,\delta^\mu{}_\nu ,\qquad \Theta_{\mu\nu}\Theta^{\mu\nu}=4E_0^2 .
\end{equation}
$\Theta^{\mu\nu}$ is invariant under duality, and for every unit timelike observer $t$ the density $j_t^\mu=\Theta^\mu{}_\nu t^\nu$ is future
causal with $\eta(j_t,j_t)=E_0^2$.
\end{proposition}

\begin{proof}
Duality invariance is a polynomial identity. Both sides of \eqref{en:rainicheq} are covariant, so by \cref{en:orbits} it suffices to take
$z=(r,0,0)$, where $\Theta^\mu{}_\nu=\tfrac{r^2}2\operatorname{diag}(1,1,-1,-1)$; null fields follow by continuity. That $j_t$ is causal and future
directed is the dominant energy condition of the Maxwell field \cite{HawkingEllis73}.
\end{proof}

This is inherited \cite{Rainich1925,Synge56}. The rest energy of a field point is the Lorentz scalar of the tensor whose $00$ component is the
undivided energy; but each observer $t$ sees a different density $j_t$, all of Lorentz length $E_0$, and no single four-vector carries the
energy and momentum of a field point.

% --- source: tex/s14b_energy#5 (The rest energy as a minimum, and two boost factors)
\subsection{The rest energy as a minimum, and two boost factors}
\label{sec:energy-min}

\begin{theorem}[Variational form of the rest energy]
\label{en:min}
\begin{enumerate}[label=(\roman*),leftmargin=2.2em]
\item For non-null $z$, $E_0$ is the minimum of $E_6$ over the Lorentz orbit, and over the orbit of Lorentz transformations and duality,
attained exactly where $\mathbf S=0$; for null $z\neq0$ the infimum $0$ is not attained. Moreover $E_6=E_0/\sqrt{1-v_E^2}$ with $v_E:=|\mathbf S|/E_6<1$.
\item If $\bB\parallel\bE\parallel\mathbf n$ and the boost $e^{\chi K}$ of rapidity $\chi$ has axis orthogonal to $\mathbf n$, then
$E_6=E_0\cosh2\chi$, $|\mathbf S|=E_0\sinh2\chi$ and $v_E=\tanh2\chi=2\beta/(1+\beta^2)$ with $\beta=\tanh\chi$. A boost along $\mathbf n$ changes
neither $E_6$ nor $\mathbf S=0$.
\item The observer moving with velocity $\mathbf v_{\rm rest}=\mathbf S_{\rm t}/(E_6+E_0)$ sees $\mathbf S=0$ and $E_6=E_0$, and
$|\mathbf v_{\rm rest}|=|\mathbf S|/(E_6+E_0)$, which is $\beta$ in the setting of (ii).
\end{enumerate}
\end{theorem}

\begin{proof}
(i) By \eqref{en:lagrange}, $E_6\ge E_0$ with equality iff $\mathbf S=0$, and \cref{en:orbits} supplies points with $\mathbf S=0$ in both orbits; for null
$z$ the boost of \cref{en:orbits}(ii) drives $E_6$ to $0$. Dividing \eqref{en:lagrange} by $E_6^2$ gives the last formula. (ii) Take $\mathbf n=e_1$
and $X=\sigma_1e_1+\sigma_2e_4$. By \cref{ex:emboost} and the same computation for the other basis vectors, $[K_3,e_1]=e_5$, $[K_3,e_5]=e_1$,
$[K_3,e_4]=-e_2$ and $[K_3,e_2]=-e_4$, so $\ad_{K_3}^3=\ad_{K_3}$ and $e^{\chi\ad_{K_3}}=\id+\sinh\chi\,\ad_{K_3}+(\cosh\chi-1)\ad_{K_3}^2$. Hence
$\bB=(\sigma_1\cosh\chi,-\sigma_2\sinh\chi,0)$, $\bE=(\sigma_2\cosh\chi,\sigma_1\sinh\chi,0)$ and $\mathbf S=-E_0\sinh2\chi\,e_3$; likewise for $K_2$,
while $K_1$ fixes $e_1$ and $e_4$. (iii) $|\mathbf S|/(E_6+E_0)=\sinh2\chi/(\cosh2\chi+1)=\tanh\chi$. For the direction, $e^{\chi K_3}$ is the passage to
a frame moving with velocity $-\tanh\chi\,e_3$ (after \cref{thm:emfield}), in which the rest frame moves along $+e_3$, and so does
$\mathbf S_{\rm t}=-\mathbf S$.
\end{proof}

The theorem is inherited \cite{LandauLifshitz2}. Its velocity is the rest-frame velocity of \cref{sec:orth}, where the Poynting vector is the
textbook $\mathbf S_{\rm t}$, as \cref{sec:conventions} stipulates for formulas from the literature; with the frame vector $\mathbf S=\bE\times\bB$
the direction would be reversed. A field point carries two boost factors, $E_6=E_0\cosh2\chi$ against $E=mc^2\cosh\chi$ for a particle, because the
energy density is a component of a second-order tensor and is quadratic in the boost: the rest frame moves with $\beta$, the energy with the
relativistic sum $2\beta/(1+\beta^2)$ of $\beta$ with itself, and an observer riding with the energy velocity sees the same $E_6$ and the reversed
Poynting vector. $E_0$ is the rest-energy density of a field point, not the rest energy of a particle carried by the flux.

% --- source: tex/s14b_energy#6 (Duality on the invariants, and the stabilisers of field poin)
\subsection{Duality on the invariants, and the stabilisers of field points}
\label{sec:energy-duality}

The duality rotation $D_\vartheta=e^{\vartheta G}=\cos\vartheta\,\id+\sin\vartheta\,G$ acts as $z\mapsto e^{i\vartheta}z$: it preserves $E_6$ and
$\Theta^{\mu\nu}$, multiplies $q$ by $e^{2i\vartheta}$, preserves the positive form for every $\vartheta$ and $\eta$ only for $\vartheta\in\pi\Z$
(\cref{sec:orth}). For $q\neq0$ exactly the angles $\vartheta\in\pi\Z$ fix $q$, and for $q=0$ every angle does. If $q\neq0$, \cref{prop:planes}
supplies Lorentz coordinates with $z=(f_m+if_e)\hat e_1$, $f_m,f_e$ real; then $q=(f_m^2+f_e^2)e^{2i\varphi}$ with $\varphi=\arg(f_m+if_e)$, and
$D_{-\varphi}$ makes the field purely magnetic, $z=|f_m+if_e|\,\hat e_1$. The angle must be the argument and not the principal value of
$\arctan(f_e/f_m)$, which gives $-|f_m+if_e|\,\hat e_1$ when $f_m<0$. All of this is the inherited duality rotation of the Maxwell invariants
\cite{Rainich1925,Calkin65,DeserTeitelboim76}.

\begin{proposition}[The rest energy is the doubly invariant modulus]
\label{en:doubly}
Of the four functions $E_6$, $\Re q$, $\Im q$ and $E_0$, the first is invariant under $U(1)_G$ but not under $\lG$, the next two are invariant under
$\lG$ but not under $U(1)_G$, and $E_0=\tfrac12|q|$ is invariant under both. On the space of Lorentz orbits of nonzero vectors, labelled by
$q\in\C$, duality rotates the label through $2\vartheta$; the stabiliser of a nonzero label is $\{D_0,D_\pi\}=\{\pm\id\}$.
\end{proposition}

\begin{proof}
The six generators of $\lG$ commute with $G$ and lie in $\so{3,3}$, so $\eta(X,AX)=\eta(X,GAX)=0$; a boost $K_i$ is symmetric for the positive form
and $X^{\!\top}K_iX\neq0$ for generic $X$. The labelling is the transitivity of $\lG$ on the level sets of $q$ with the origin removed (proof of
\cref{thm:achiral}), and $D_\vartheta$ acts on the label by $q\mapsto e^{2i\vartheta}q$.
\end{proof}

The proposition is inherited. With $\mathbf g=\mathbf S_{\rm t}/c^2$ and units restored, \eqref{en:lagrange} reads $u^2-c^2|\mathbf g|^2=E_0^2$, the
field form of $E^2-p^2c^2=m^2c^4$ (\cref{en:SI}). A nonzero rest energy thus breaks $U(1)_G$ to $\Z_2$ on the space of Lorentz orbits, and $\arg q$ is the angle along the broken
direction. That is a
reading of the proposition, and it comes with two qualifications. The breaking is a property of a configuration, not of a vacuum: the energy is
minimised by the zero field and no potential selects $|q|\neq0$ (\cref{sec:orth}). And the broken symmetry is not one of the interior Lagrangian,
which duality rotates into its $\eta G$ partner.

\begin{theorem}[Stabilisers of field points]
\label{en:stab}
Let $z\neq0$.
\begin{enumerate}[label=(\roman*),leftmargin=2.2em]
\item $\Stab_{\lG}(z)=\{(\zeta z)\times\cdot\,:\,\zeta\in\C\}$ is abelian of real dimension two. If $q\neq0$ and $\hat z=z/\sqrt q$, then
$R=\hat z\times\cdot$ satisfies $R^3=-R$ (a rotation) and $(GR)^3=GR$ (a boost) about the common complex axis $\hat z$: the stabiliser is
$\so2\oplus\so{1,1}$. If $q=0$, every element $N$ of the stabiliser has $N^3=0\neq N^2$: null rotations, $\R^2$.
\item In $\lG\oplus\R G$ a non-null $z$ has the same stabiliser, with no $G$ component. A null $z=\mathbf n_1+i\mathbf n_2$ has a three-dimensional
stabiliser $\cong\mathfrak e(2)$, spanned by the null rotations and $L_{\mathbf n}+G$, with $\mathbf n=\mathbf n_1\times\mathbf n_2/|\mathbf n_1|^2=\bB\times\bE/|\bB\times\bE|$
the propagation axis, which acts on the null rotations with eigenvalues $\pm i$. For $z=(1,i,0)$, $(L_3+G)z=0$ and $e^{\vartheta G}z=e^{-\vartheta L_3}z$.
\item In $\lint$, on the interior four-plane $\langle e_1,e_4,e_5,e_6\rangle$, $\Stab(e_1)=\so3$ and $\Stab(e_1+e_6)=\mathfrak e(2)$, both
three-dimensional.
\end{enumerate}
\end{theorem}

\begin{proof}
(i) $w\times z=0$ iff $w\in\C z$. From $(z\times)^2w=z(z\cdot w)-qw$ one has $(z\times)^3=-q\,(z\times)$; with $\hat z\cdot\hat z=1$ this is $R^3=-R$, and
$G\leftrightarrow i$ gives $(GR)^3=GR$; for $q=0$, $(z\times)^3=0\neq(z\times)^2$. (ii) The condition is $w\times z+i\vartheta z=0$, and dotting with $z$ gives
$i\vartheta q=0$. For null $z$, $\mathbf n_1\cdot\mathbf n_2=0$ and $|\mathbf n_1|=|\mathbf n_2|$ give $\mathbf n\times\mathbf n_1=\mathbf n_2$ and
$\mathbf n\times\mathbf n_2=-\mathbf n_1$, so $\mathbf n\times z=-iz$, $(L_{\mathbf n}+G)z=0$ and $[L_{\mathbf n},z\times\cdot]=(\mathbf n\times z)\times\cdot=-i(z\times\cdot)$.
(iii) is the classical computation of the little groups of a timelike and a null vector.
\end{proof}

The theorem is inherited: (iii) is Wigner's classification \cite{Wigner39,WeinbergQFT}, and (i)--(ii) are the classical stabilisers of a
bivector and the identity of helicity and duality on null fields \cite{PenroseRindler84,Calkin65}, which on the frame are \cref{thm:axis} (the
axis of a null field) and \cref{ex:Gcharge} (the $G$-charge is the helicity). It separates two notions that the word \emph{mass} runs together.
The little group of a momentum is three-dimensional with or without mass, and the stabiliser of a field point is two-dimensional with or without
$q$. A mass changes the type of a Lorentz stabiliser, never its dimension, and breaks no Lorentz generator, since $q$ and the square of a momentum
are invariants. What a nonzero $q$ removes is the duality generator: a null field point is invariant under duality up to a rotation about its
axis, so the group $\mathfrak e(2)$ of Wigner reappears once duality is admitted, with the helicity as its compact generator, while for a non-null
point duality is broken to $\{\pm\id\}$.

% --- source: tex/s14b_energy#7 (Where c enters)
\subsection{\texorpdfstring{Where $c$ enters}{Where c enters}}
\label{sec:energy-c}

Adding the dimensionless $G$ to a derivative along a clock coordinate needs a length $\ell$ and gives the rate $\kappa=c/\ell$ of
\eqref{eq:kappa}. On the frame, $c$ enters as follows.

\begin{proposition}[$c^2$ is an isotropic weight]
\label{en:onec}
Let the temporal triple carry $\bB$ and the spatial triple $\bE$, let $C:V^+\to V^-$ be the conversion and
$G_C=\bigl(\begin{smallmatrix}0&-C^{-1}\\ C&0\end{smallmatrix}\bigr)$, so that $G_C^2=-\id$.
\begin{enumerate}[label=(\roman*),leftmargin=2.2em]
\item $G_C$ commutes with the diagonal rotations iff $C=c\,\id$.
\item For $C=c\,\id$, writing $G_c$ for $G_C$, the forms invariant under the diagonal rotations and under $e^{\vartheta G_c}$ are the multiples of $c^2|\bB|^2+|\bE|^2$.
\item $(\kappa G)^2=-c^2/\ell^2\,\id$, and $\ell$ is not fixed.
\end{enumerate}
\end{proposition}

\begin{proof}
(i) The lower-left block of $[G_C,R\oplus R]$ is $CR-RC$, and Schur's lemma gives $C=c\,\id$; this is \cref{thm:onec}. (ii) For
$N=\bigl(\begin{smallmatrix}\alpha&\gamma\\ \gamma&\beta\end{smallmatrix}\bigr)\otimes\id$,
$G_c^{\!\top}N+NG_c=\bigl(\begin{smallmatrix}2c\gamma&c\beta-\alpha/c\\ c\beta-\alpha/c&-2\gamma/c\end{smallmatrix}\bigr)\otimes\id$, which vanishes iff
$\gamma=0$ and $\alpha=c^2\beta$. (iii) $G^2=-\id$.
\end{proof}

Part (i) is \cref{thm:onec}, fixed there, and it fails for $C=\operatorname{diag}(c_1,c_2,c_3)$ with distinct entries; part (ii) is inherited, and
part (iii) is constructed, a definition, since $G^2=-\id$ is adopted in \cref{thm:neutral} and $\ell$ is chosen. Even the weight is a unit convention: $\Phi=\operatorname{diag}(c,c,c,1,1,1)$
conjugates $G_c$ to $G$ and turns $c^2|\bB|^2+|\bE|^2$ into $|\Phi X|^2$, so with $\bB$ measured as $c\bB$, as in Heaviside--Lorentz units, $c$
disappears from $E_6$ and appears only in the rest-mass density $\rho_0$ below.

\begin{proposition}[SI form]
\label{en:SI}
With $\varepsilon_0\mu_0c^2=1$, $z_c:=c\bB+i\bE$, $u=\tfrac{\varepsilon_0}2|z_c|^2$, $\mathbf S_{\rm t}=\bE\times\Btext/\mu_0$, $\mathbf g=\mathbf S_{\rm t}/c^2$,
$q_c:=z_c\cdot z_c=c^2|\Btext|^2-|\bE|^2-2ic\,\Btext\cdot\bE$ and $u_0:=\tfrac{\varepsilon_0}2|q_c|$, \cref{en:split} applied to $z_c$ gives
\begin{equation}
\label{en:SIeq}
u^2=c^2|\mathbf g|^2+u_0^2,\qquad \rho_0:=u_0/c^2,\qquad [u]=[c\,\mathbf g]=[u_0]=\mathrm{J\,m^{-3}},\quad [\rho_0]=\mathrm{kg\,m^{-3}} .
\end{equation}
\end{proposition}

This is inherited, and calling $\rho_0$ a rest-mass density is a reading, since $(u,c\,\mathbf g)$ is not a four-vector (\cref{en:rainich}).
It is the exact content of the clause, taken up in \cref{sec:energy-readings}, that $c^2$ is the symmetric part of $E=mc^2$, and it is less than the statement suggests. The constant $c$
is the unique conversion between the triples that commutes with their common rotations, which are the residual symmetry of both $E_6$ and $E_0$;
in SI units it enters $E_6$ as the isotropic weight $c^2$ of $|\bB|^2$, which a choice of units removes. It is a conversion, not a symmetry, and
nothing breaks it: $u_0$ uses the same $c$, and $\rho_0=u_0/c^2$ only converts units.

% --- source: tex/s14b_energy#8 (The mass of a whole)
\subsection{The mass of a whole}
\label{sec:energy-whole}

For a free field with compactly supported Cauchy data, $(U,\mathbf P)=\int(u,\mathbf S_{\rm t})\,d^3x$ is a four-vector and its Lorentz length $M$
is positive whenever $q\neq0$ on a set of positive measure, while $M$ does not detect orthogonality (\cref{sec:orth}). The following bounds and
examples locate $M$ relative to the rest energies of the field points.

\begin{proposition}[Mass of a whole]
\label{en:whole}
\begin{enumerate}[label=(\roman*),leftmargin=2.2em]
\item $M^2=U^2-|\mathbf P|^2\ge U\int(u-|\mathbf S|)\,d^3x\ge0$, and in every frame $M\ge\int E_0\,d^3x$, with equality iff the densities
$(u,\mathbf S_{\rm t})(x)$ are all proportional.
\item A plane wave has $M=0$. In SI units the waves $E_x=a\cos(kz\mp\omega t)$, $cB_{{\rm t}y}=\pm a\cos(kz\mp\omega t)$ superpose to
$E_x=2a\cos kz\cos\omega t$, $cB_{{\rm t}y}=2a\sin kz\sin\omega t$. Per cell (unit area, length $2\pi/k$) one has $U=2\pi\varepsilon_0a^2/k$ and
$\mathbf P=0$, so $Mc^2=U$, although each constituent has $q_c=0$ and $U_\pm=|\mathbf P_\pm|c$. The total has $q_c=-2a^2(\cos2kz+\cos2\omega t)$,
$\int_{\rm cell}\Theta^{zz}=U$ and $\int_{\rm cell}\Theta^{xx}=-U\cos2\omega t$, and the integrated rest energy $\int_{\rm cell}u_0$ oscillates between
$\tfrac2\pi U$ and $U$.
\item With $c=1$ and $f(s)=(1-s^2)^2$ for $|s|\le1$, $f(s)=0$ otherwise, the field $E_x=f(z-t)+f(z+t)$, $B_{{\rm t}y}=f(z-t)-f(z+t)$ has, per unit
transverse area, $U=\tfrac{512}{315}$ and $\mathbf P=0$ at all times, so $M=U$, with $(U,P_z)$ a two-vector under boosts along $z$ only. For $|t|\ge1$
the two pulses are disjoint, $q\equiv0$ and $\int E_0\,d^3x=0$.
\end{enumerate}
\end{proposition}

\begin{proof}
(i) Conservation and finite propagation speed make $\int_\Sigma\Theta^{\mu\nu}d\Sigma_\nu$ independent of the spacelike hyperplane $\Sigma$
\cite{Laue1911}. Since $u-|\mathbf S|=E_0^2/(u+|\mathbf S|)\ge0$ and $M^2=(U-|\mathbf P|)(U+|\mathbf P|)$, the first bound follows. Each $(u,\mathbf S_{\rm t})(x)$ is
future causal of Lorentz length $E_0(x)$ (\cref{en:rainich}), and future causal vectors obey the reversed triangle inequality
$|P_1+P_2|\ge|P_1|+|P_2|$, with equality iff they are proportional. (ii) The cell integrals are elementary. For the last statement, with
$s=\sin^2\omega t$ one has $q_c=4a^2(s-\cos^2kz)$, so $u_0=2\varepsilon_0a^2|s-\cos^2kz|$; its cell mean is least at the median $s=\tfrac12$ of the
arcsine distribution of $\cos^2kz$, where it is $\tfrac2\pi U$, and equals $U$ at $s=0$ and $s=1$. (iii) $u=f(z-t)^2+f(z+t)^2$ and $S_{{\rm t}z}=f(z-t)^2-f(z+t)^2$, the
cross terms cancelling, and $q=-4f(z-t)f(z+t)$; finally $2\int_{-1}^1(1-s^2)^4ds=\tfrac{512}{315}$.
\end{proof}

The first bound and item (ii) are inherited. The cell in (ii) is not isolated: $\int\Theta^{zz}=U$ is the pressure on whatever confines it, so its
$(U,\mathbf P)$ is a four-vector only with the walls included, as for radiation in a box \cite{Laue1911}. The inequality $M\ge\int E_0$ and
item (iii) obstruct the reading that the mass of a whole emerges from the interference that makes $q\neq0$: interference suffices for $M>0$, by (i) and the example (ii),
but it is not necessary, by (iii). A whole gets its mass from the non-collinearity of the energy--momenta of its parts, as a pair of photons does.

% --- source: tex/s14b_energy#9 (Splitting the temporal energy lowers the mass)
\subsection{Splitting the temporal energy lowers the mass}
\label{sec:energy-temporal}

There is a second reading of the phrase that the energy is six-dimensional: the three temporal directions of $V^+$ carry three energies, and a
mass would arise by differentiating them. On the arena this reading fails.

\begin{proposition}[Splitting the temporal energy lowers the mass]
\label{en:temporal}
Let $P=\sum_i(\mathsf E_i/c)e_i+\sum_ap_ae_{a+3}$ with $\eta(P,P)=\mu^2c^2$, let $\widehat W$ be the symmetric clock of \cref{prop:clock},
$\mathsf E_W:=(\mathsf E_1+\mathsf E_2+\mathsf E_3)/\sqrt3$, and let $\mathsf E_\perp$ be the part of $(\mathsf E_1,\mathsf E_2,\mathsf E_3)$ orthogonal
to $(1,1,1)$, the energies of the relative times. The mass read by the clock, $m_4^2c^4:=\mathsf E_W^2-|\bp|^2c^2$, is
\[
m_4^2c^4=\mu^2c^4-|\mathsf E_\perp|^2\le\mu^2c^4 .
\]
For the clock $e_1$ instead, $m_4^2c^4=\mu^2c^4-\mathsf E_2^2-\mathsf E_3^2$; for example $\mathsf E=(3,4,0)$, $\bp=(5,0,0)$ and $c=1$ give $\mu=0$ and
$m_4^2=-16$.
\end{proposition}

\begin{proof}
$|(\mathsf E_1,\mathsf E_2,\mathsf E_3)|^2=\mathsf E_W^2+|\mathsf E_\perp|^2$ and $\mu^2c^4=c^2\eta(P,P)=|(\mathsf E_1,\mathsf E_2,\mathsf E_3)|^2-|\bp|^2c^2$.
\end{proof}

This is an obstruction, and it restates \cref{prop:toy} and \cref{ex:tachyons} in a form covariant under the choice of clock. In the momentum
reading the differentiation of the temporal energy subtracts from the mass and can make it tachyonic, and under the surplus constraints of \cref{thm:wlgauge}, that is under \cref{assum:clockred}, $\mathsf E_\perp=0$, nothing is left to differentiate. In the field-strength reading the motion part adds instead,
$E_6\ge E_0$. Only the latter carries the structure of an undivided energy with an internal differentiation.

% --- source: two_readings/r5_synthesis#0 ((head))
% Synthesis of the two-readings note.

% --- source: two_readings/r5_synthesis#1 (Synthesis)
\section{Synthesis}
\label{r5:sec:main}

The four questions of the introduction have the following answers; each is stated with its class, and the sections hold the proofs.

\paragraph{(1) Lagrangians.} For the classical reading the free worldline of the arena gives the six-momentum and the shell
$E^2=c^2|\bp|^2+M^2c^4-c^2|\mathbf q|^2$, and the interior shifted Lagrangian gives $\ddot Y=-\kappa^2Y$ with canonical momentum obeying
$\dot\pi=\kappa G\pi$ (restated, \cref{r1:prop:free,r1:prop:shift}). The Dirac Lagrangian $\tfrac12\psi^{\!\top}(K^\mu\partial_\mu-mT)\psi$ gives
$\gamma^\mu\partial_\mu\psi=-mJ\psi$, its current and its stress tensors (restated and inherited, \cref{r1:prop:dirac}). The energies are
governed by the Maxwell Lagrangian $-\tfrac12\Re q$, whose Hamiltonian is the undivided energy $E_6$ and whose first-order form makes
duality an off-shell symmetry (inherited, \cref{r1:prop:maxwell,r1:prop:E6}). The rest energy $E_0=\tfrac12|q|$, taken as a Lagrangian, is
Lorentz invariant and conformal but not duality invariant in the sense of Gaillard and Zumino, has an indefinite energy, and is singular on
radiation; its natural place is the first-order term of the ModMax family. Neither $E_0$ alone nor any Hamiltonian $aE_6+bE_0$ with
$a^2-b^2\neq1$ is Lorentz invariant (fixed, \cref{r1:thm:L0,r1:prop:ham}).

\paragraph{(2) The two $X$'s.} $X_{\rm cl}$ is a displacement, with converted times on $V^+$ and lengths on $V^-$; $X_{\rm EM}=(\bB,\bE)$ is
the pair labelling the six components of the field two-form. They share the vector space and the Clifford action and differ in units,
natural Lorentz algebra ($\lint$ on a four-plane against $\lG$ on the two-form), canonical conjugates and observables. No linear
map equivariant for the interior Lorentz algebra relates them in either direction, whatever dictionary is chosen (fixed,
\cref{r1:thm:type}). The two readings must therefore be interpreted separately.

\paragraph{(3) The four-momentum.} With a clock $\hat n$ on the temporal triple, $E_{\hat n}=\R \hat n\oplus V^-$ has signature $(1,3)$; the projection
$P_{\hat n}=(E/c,\bp)$ of the six-momentum is a four-vector of $\mathrm{SO}^+(1,3)_{\hat n}$, with $\bp$ conjugate to the three lengths and $E/c$ to the
time read by $\hat n$, and it is conserved for every dynamics independent of $(ct_{\hat n},\mathbf x)$ (fixed, \cref{r2:prop:cov,r2:prop:cons}). Of three mass shells, the ambient one is the statement of \cref{sec:energy} that splitting the temporal energy lowers the mass, the first-class reduction is the surplus-time constraint of \cref{sec:constraints} with the clock $\hat n$, and the third is constructed (restated and
constructed, \cref{r2:prop:shellI,r2:prop:shellII,r2:prop:shellIII}). The four-dimensional mass read by a clock satisfies
$m_4(\hat n)^2c^4=\mu^2c^4-|\hat n\times\mathbf E|^2$, never above the six-dimensional value (fixed, \cref{r2:prop:rotate}). The field has no linear
covariant map to a four-vector; its four-momentum is the integral of its stress tensor, a four-vector only for an isolated whole, and
the photon's $P=\hbar k$ is not derived in the arena (fixed, inherited and open, \cref{r2:sec:field}). The anti-isometry no-go holds for
every clock and obstructs the identification of the interior momentum with the six-momentum, not the duality of the field
(fixed, \cref{r2:prop:nogo,r2:prop:leak}).

\paragraph{(4) The pair on the octants.} The octants are the joint eigenlines of the pole operators, the Fock basis of $S\cong\Lambda^*(L)$.
The field-octant map places the electric member of the pair on the three octants of weight one and the magnetic member on the three of
weight two, uniquely up to a scale and a quarter turn of duality once rotations and Maxwell's equations are imposed; $\bE$ is then a
one-form and $\bB$ a two-form (fixed and constructed, \cref{r3:thm:phi,r3:prop:maxunique}). Maxwell's four equations are one first-order
equation on $S$, with the two Gauss laws on the two poles and the evolution laws on the field octants (fixed, inherited in substance,
\cref{r3:thm:maxwell}); the spatial derivatives in it multiply the temporal generators, a property of the representation and not a
derivation of space from time. The mass form pairs each electric octant with its magnetic antipode and gives $-\bE\cdot\bB$ (fixed,
\cref{r3:prop:energies}). Clifford multiplication by the field point mixes the two members of the pair and does not realise it
(\cref{r3:prop:poles}). Boosts and the duality rotation act on the octants only through odd operators (obstructed for even operators,
\cref{r3:thm:noeven,r3:prop:odd}). The construction represents the classical field strength; it quantises nothing.

\paragraph{The bridge.} The element shared by the two readings is $T=e_1e_2e_3=\hat nJ(\hat n)$: for every clock it is the Dirac mass form of the
particle reading, it lifts the frame parity, and on the octants $-T$ is the duality of the field reading (fixed, \cref{r4:thm:roles}); it is,
up to scale, the only element that is the clock times a Dirac complex structure for every clock (fixed, \cref{r4:thm:unique}). The field
reading does not see the clock and the particle reading does (fixed, \cref{r4:prop:clockdep}). The Dirac action, the octant representation
of the field and the spin lift of $\lG$ are three inequivalent Lorentz representations on one module, whose joint commutant is the
scalars; no bilinear coupling of a Dirac field to the field octants is covariant, while the trilinear couplings are exactly the Pauli
term and its dual, at every clock (fixed and inherited, \cref{r4:prop:ineq,r4:thm:couplings}). The answer to the question whether $T=\hat nJ(\hat n)$
incorporates consistently into both readings is therefore yes at the level of the module: the two readings keep their Lorentz
representations, and they meet in $T$ and in the Pauli coupling. A covariant Dirac--Maxwell dynamics in the arena, with the minimal coupling
through the potential and the charge $J(\hat n)$, a dynamics that selects the clock, and the quantisation of either reading remain open.

